\documentclass[11pt]{article}

\usepackage[T1]{fontenc}
\usepackage[utf8]{inputenc}
\usepackage{lmodern}
\usepackage[margin=0.78in]{geometry}
\usepackage{graphicx}
\usepackage{booktabs}
\usepackage{tabularx}
\usepackage{array}
\usepackage{microtype}
\usepackage{xcolor}
\usepackage{caption}
\usepackage[numbers,sort&compress]{natbib}
\usepackage{hyperref}
\usepackage{fancyhdr}
\usepackage{enumitem}
\usepackage{longtable}
\usepackage{listings}
\usepackage{placeins}
\usepackage{newunicodechar}
\newunicodechar{’}{\textquoteright}

\definecolor{molblue}{RGB}{20,96,126}
\definecolor{molteal}{RGB}{0,126,133}
\definecolor{callout}{RGB}{236,247,249}

\hypersetup{
  colorlinks=true,
  linkcolor=molblue,
  citecolor=molblue,
  urlcolor=molblue,
  pdftitle={Molarium: Agent-Generated Adaptive Software for Molecular Design},
  pdfauthor={Rafal P. Wiewiora and Woody Sherman}
}

\captionsetup{font=small,labelfont=bf}
\setlength{\parindent}{0pt}
\setlength{\parskip}{0.45em}
\setlist{nosep}

\pagestyle{fancy}
\fancyhf{}
\fancyhead[L]{\footnotesize\color{gray}Molarium: Agent-Generated Adaptive Software for Molecular Design}
\fancyhead[R]{\footnotesize\color{gray}Wiewiora and Sherman}
\fancyfoot[C]{\thepage}
\renewcommand{\headrulewidth}{0.4pt}

\makeatletter
\renewcommand\section{\@startsection{section}{1}{\z@}%
  {-2.2ex \@plus -1ex \@minus -.2ex}%
  {1.1ex \@plus .2ex}%
  {\normalfont\Large\bfseries\color{molblue}}}
\makeatother

\lstset{basicstyle=\ttfamily\footnotesize,breaklines=true,columns=fullflexible,keepspaces=true,literate={’}{{\textquoteright}}1}

\newcommand{\safeincludegraphics}[2][]{%
  \IfFileExists{#2}{\includegraphics[#1]{#2}}{%
    \fbox{\parbox[c][1.8in][c]{0.92\linewidth}{\centering\small Figure file not included in this source package.}}}%
}

\newcommand{\callout}[1]{%
  \begin{center}
  \fcolorbox{molteal}{callout}{\parbox{0.90\linewidth}{\centering\bfseries #1}}
  \end{center}
}

\newcommand{\appendixsection}[2]{%
  \clearpage
  \refstepcounter{section}%
  \setcounter{subsection}{0}%
  \fancyhead[L]{\footnotesize\color{gray}Molarium --- Appendix \Alph{section}}%
  \begin{center}
    {\Large\bfseries\color{molblue}APPENDIX \Alph{section}}\\[0.55em]
    {\LARGE\bfseries\color{molblue}#1}
  \end{center}
  \vspace{0.8em}
  \addcontentsline{toc}{section}{Appendix \Alph{section}: #1}%
  \label{#2}%
}

\title{\textcolor{molblue}{\textbf{Molarium: Agent-Generated Adaptive Software for Molecular Design}}\\[0.5em]
\large A perspective on user-directed adaptive software, persistent molecular state, and scientific reproducibility in drug discovery}
\author{%
\begin{tabular}{@{}ccc@{}}
\textbf{Rafal P. Wiewiora} & \textbf{Woody Sherman}\\
\normalsize PsiThera & \normalsize PsiThera\\
\small\texttt{rafal.wiewiora@psithera.com} &
\small\texttt{woody.sherman@psithera.com}
\end{tabular}}
\date{September 6, 2026}

\begin{document}
\maketitle
\thispagestyle{fancy}

\begin{abstract}
Commercial drug discovery software has traditionally bundled three sources of value: the interface, the scientific models, and the specialized expertise required to run the models through the interface, thereby translating outputs into design decisions. Large language models, coding agents, and tool-using agents are beginning to unbundle this paradigm. Interfaces, analyses, workflow logic, and substantial portions of scientific software can increasingly be generated to meet the needs of a particular scientist, project, and question, while agents can automate sophisticated workflows that once required specialist users.

This shift repositions where competitive advantage is created in computational drug discovery. More intuitive interfaces and faster execution alone do not fix model inaccuracies or broaden the applicability domain. As the underlying tooling becomes increasingly accessible and fluid, strategic differentiation will shift away from familiarity with a particular software package and toward superior preditive algorithms, more generalizable models, proprietary experimental data, accumulated discovery intelligence, and the human discernment needed to frame the right questions and evaluate the answers.

Molarium, available at \url{https://molarium.org}, provides a case study of this transition in the context of molecular design software. Created during an intensive weekend effort of a tireless drug-discovery scientist and an AI coding agent, this open-source molecular workbench operates exclusively in a standard web browser. It integrates local molecular graphics, editing, conformer generation, energy minimization, molecular dynamics, machine-learned potentials, and protein-structure prediction workflows. Beneath its rapidly adaptable interface, Molarium maintains a versioned molecular state together with user actions, method specifications, provenance, and validation evidence. Decoupling an adaptable interface from an enduring scientific core points toward a broader architecture for generated adaptive software in drug discovery.
\end{abstract}

\textbf{Keywords:} agentic software development; computational chemistry; molecular design; WebGPU; scientific validation; provenance; discovery intelligence

\section{From fixed software products to generated adaptive software}

In computational drug discovery, success depends on deploying the right methodology at the right time, supported by appropriate parameters, explicit assumptions, and sound scientific context. Historically, this work has been concentrated among computational specialists and has required substantial operational overhead. Docking depends on protein preparation, protocol selection, constraints, and pose evaluation. Molecular dynamics requires complex system setup, solvation, ion placement, force-field assignment, compute infrastructure, equilibration, production simulation, potential for enhanced sampling, trajectory analysis, visualization, and interpretation in the context of project-specific decisions. Free-energy calculations add further complexities to  molecular simulations with alchemical perturbations, sampling windows, convergence assessment, error analysis, and alignment with experimental conditions. Comparable layers of infrastructure and expertise surround cheminformatics, quantum chemistry, ADME property prediction, generative molecular design, and protein structure prediction.

Historically, commercial drug discovery software has created considerable value by simplifying and integrating many of these activities. Graphical interfaces reduced technical barriers, workflow engines connected fragmented programs, and software packages encoded operational details needed to apply methods consistently. A single static inferface (perhaps updated quarterly) coupled scientific questions to a fixed software environment. User-specific requests competed with broader corporate product roadmaps that attempted to balance the needs of a diverse user base. Even when feasible, implementing new methods from the literature or integrating open-source software often required software engineering support. Scientific users adapted their questions to the capabilities and conventions of the available tools within their preferred/available graphical interface rather than adapting the software to their needs.

Large language models and coding agents are beginning to change this relationship. A scientist with deep drug-discovery expertise but limited software programming experience can increasingly describe an application, analysis, visualization, or workflow in natural language and have software constructed around the scientific need. Tool-using agents can operate established methods, build complex workflows, implement methods from the literature, handle failures, assemble results into insights that support decisions, and even optimize processes based on iterative feedback. Earlier systems such as ChemCrow and ChemGraph established proof-of-concept for language models coordinating chemistry tools and computational workflows \citep{bran2024chemcrow,boiko2023coscientist,pham2026chemgraph,macknight2025rethinking}. While their practical scope was limited compared with current coding agents, they helped establish the feasibility of coupling language models to executable scientific tools.

The consequence is a shift in how drug hunters can interact with software. Scientists can increasingly customize software around project-specific questions, available data, and their preferred way of reasoning about a problem. A medicinal chemist optimizing a macrocycle may need a different interface and set of calculations than a structural biologist testing a mechanistic hypothesis or a program team deciding which compounds to advance from a virtual screening campaign. In this model, the interface becomes an adaptive and potentially transient view over a more durable scientific substrate. Changes to the interface do not imply changes to the underlying record of the work: molecular states, computational methods, parameters, intermediate results, user actions, and software versions should persist and link through an auditable provenance history. This separation between a mutable graphical layer and a persistent scientific state is important for reproducibility, inspection, and reconstruction of prior decisions. Although Molarium presents as a relatively simple browser-based application, it is built on a substantially richer underlying architecture for managing versioned molecular state, tracking computational provenance, and executing adapted implementations of established molecular modeling algorithms efficiently within the constraints of a local web browser.

Molarium’s development directly reflects this paradigm. It began as a need for a better environment to view and modify molecular structures; the first version evolved through active scientific use over a weekend. Scientific problems drove software changes: unrealistic geometries led to refined protocols; unconverged energies prompted explicit numerical reference checks; ambiguities about execution backends led to clearer identification of the active engine; and unvalidated methods became explicit boundaries rather than implicit assumptions. The software evolved with the scientific questions it was built to address. Appendix~\ref{app:architecture} describes the execution architecture, numerical validation, and an executable example of this design history. Appendix~\ref{app:manual} provides a module-by-module reference distinguishing current, experimental, repository-only, proposed, and retired functionality.

\section{Molarium and the value of a persistent molecular state}

Molarium is a local-first molecular viewer, builder, and simulation workbench released under the Apache License 2.0. Calculations execute on the user's local device using WebAssembly or WebGPU, allowing molecular graphics, editing, cheminformatics, conformer exploration, molecular mechanics, molecular dynamics, machine-learned potentials, and selected protein-structure prediction workflows to coexist in a standard browser environment. The browser reduces installation and orchestration overhead and avoids remote job submission, cloud-compute charges, and application licensing while keeping computation inside the interactive design loop. WebAssembly provides a route for mature compiled scientific libraries to run locally, while WebGPU exposes modern accelerators through a portable browser API \citep{haas2017wasm,webgpu2025}.

\begin{figure}[!htbp]
  \centering
  \safeincludegraphics[width=0.90\linewidth]{figures/fig1_molarium_interface.png}
  \caption{Molarium in its normal View workspace after browser-local Sage/CCD preparation of the experimental fragment F1-bound KRAS--SOS1 complex (PDB 6EPM). The official RCSB Chemical Component Dictionary topology supplies the BQ5/F1 bond orders and the chemically correct 2D depiction shown at lower left. The complete deposited protein assembly is shown as cartoons, while BQ5/F1 and protein residues within 5~\AA{} are rendered atomically to localize the allosteric design pocket. Solvent, cofactors, and hydrogens are hidden for clarity. This scene provides structural context; the prospective hit-to-BAY-293 trajectory in Fig.~\ref{fig:sos1story} began independently from the 5OVE/AXE structure alone.}
  \label{fig:views}
\end{figure}

\FloatBarrier

One of the advanced architectural capabilities is the preservation of a common scientific state. Molarium maintains explicit topology, coordinates, protonation state, chemical edits, preparation history, trajectories, model identity, and provenance as properties of the molecular session rather than allowing each view or calculation to construct its own interpretation of the molecule. The interface can therefore change rapidly without redefining the scientific object beneath it.

Three task-specific views illustrate this separation. ``View'' emphasizes the molecular system, including molecular representations of protein, ligand, waters, ions, contacts, measurements, and selected residues. ``Design'' exposes chemical and geometrical modifications. ``Simulate'' emphasizes preparation, simulation, and analysis. The controls differ because the scientific questions differ; topology, coordinates, method identity, and history remain attached to the same evolving molecular state. Detailed implementation and replay examples are provided in Appendix~\ref{app:architecture}, while the state mutations, exports, and failure behaviour of the corresponding user-facing modules are specified in Appendix~\ref{app:manual}.

\begin{figure}[htbp]
  \centering
  \safeincludegraphics[width=\linewidth]{figures/fig2_sos1_hit_to_bay293.png}
  \caption{Hit-only SOS1 design replayed in the Molarium interface through the public Chemist Actions API. (A) The prepared 5OVE/AXE hit (teal) and its pale local pocket are centered before prospective design; 5OVE/AXE was the sole coordinate-bearing input. The camera established here remains fixed. Between panels A and B, the replay labels Phe890 and Lys898 while the first predicted pyrazole is inspected. (B) Sequential scaffold rewriting and growth produces compound 21 (purple) in the Phe890-in pocket. The cyan halo identifies the edited ligand region; gold identifies Phe890; magenta dotted connectors are severe ligand--protein short contacts recalculated from the live coordinates. (C) A discrete side-chain search rotates Phe890 out of the growth path. Only the six moved Phe890 heavy atoms are haloed, so the conformational change can be compared directly with B without a camera refit. (D) The final predicted BAY-293 complex is shown after coupled WebGPU relaxation with transient change marks cleared. The replay ends at this prediction and never opens 5OVI. Independent post-freeze comparison with 5OVI subsequently gave a receptor-aligned, graph-symmetry-minimized ligand RMSD of 1.215~\AA{} and Phe890 $\chi_1$ values of $-172.43^{\circ}$ in the prediction and $-166.55^{\circ}$ in the crystal; those validation coordinates did not seed the replay.}
  \label{fig:sos1story}
\end{figure}

The executable replay is available at \url{https://molarium.org/sos1-hit-to-bay293}.

This separation becomes important when interfaces are inexpensive to modify. If every generated view can redefine a molecule, select a different charge model, change protonation, alter preparation, or reinterpret a trajectory, flexibility becomes a source of scientific ambiguity. Generated adaptive software therefore requires a stronger boundary between presentation and scientific state than a conventional fixed interface typically provides.

Molarium extends this principle by preserving the sequence of molecular states and user actions that produced the current state. Edits, selections, manipulations, calculations, and other interactions form an event history rather than leaving only the final coordinates. The resulting architecture is closer to version control than to a traditional molecular project file. Each molecular state can be associated with the actions that produced it, the calculation applied to it, and the evidence available at that point in the design process.

Thia concept is analogous to ``Git for molecules'', although a drug design history is richer and more complex than tracking revisions of source code. A design history includes structural changes together with evolving hypotheses, model outputs, uncertainty, and experimental evidence. In a discovery setting, such a record could capture which analogs were considered, which structures or properties influenced a design, which alternatives were rejected, and what was learned when a compound was synthesized and tested. This creates a foundation for recording medicinal chemistry design trajectories rather than storing only the compounds that survived them. Much of the information that explains how a program learned lies in alternatives that were considered and abandoned, yet conventional compound databases usually only preserve experimental endpoints.

\FloatBarrier

\begin{figure}[htbp]
  \centering
  \safeincludegraphics[width=0.94\linewidth,trim=0 15 0 0,clip]{figures/fig2_architecture.png}
  \caption{Molarium's local-first architecture. Named computational engines act on an explicit molecular state while provenance and validation evidence remain attached. The same principle provides a foundation for versioned molecular histories in which scientific actions can be replayed and compared without allowing a new interface to redefine the underlying state.}
  \label{fig:architecture}
\end{figure}

The browser also provides a useful trust boundary for proprietary work. Molarium's ``Local Lab'' mode can restrict network access and verify reviewed application assets while keeping supported calculations on the local device. Because privacy controls depend on deployment configuration, this architecture makes data movement explicit and auditable. Organizations can evaluate and enforce how molecular data are handled rather than treating transfer to remote infrastructure as an invisible property of the software stack. Appendix~\ref{app:manual} specifies the connected and Local Lab network policies and makes clear that local calculation and fully offline operation are distinct claims.

\section{Scientific contracts: provenance, validation, and falsifiability}

The emergence of rapidly adaptive software introduces a scientific vulnerability. Agentically generated workflows become especially risky when ease-of-use obscures the underlying scientific methods. As the time required to build a custom analysis view drops, that view still needs to expose the molecule, model, parameters, assumptions, and evidence that produced the result. Reducing the friction of software development without those safeguards increases the rate at which scientific ambiguity and errors can be introduced.

We use the term ``scientific contract'' for the collection of invariants and evidence that gives a computational output its meaning. A molecular-mechanics result cannot be defined by the generic label ``molecular mechanics.'' Its identity depends on the molecular state, charge assignment, parameterization, energy function, treatment of nonbonded interactions, boundary conditions, constraints, integrator, numerical implementation, and precision. A learned potential depends on its architecture, weights, preprocessing, supported chemistry, and training domain. A protein-structure prediction depends on checkpoint selection, sequence and template information, inference protocol, and other model-specific choices. The surrounding interface may adapt dynamically, but these scientific identities must remain explicit and protected from silent drift.

Molarium preserves distinctions among computational pathways even when they operate in the same workspace. RDKit supports chemical perception and conformer generation; OpenFF Sage defines a molecular-mechanics model; OpenMM Reference provides an independent numerical route for selected validation tests; STORMM-inspired WebGPU code advances homogeneous replicas; ANI-2x provides a machine-learned potential with a defined chemical domain; and OpenFold supports a separate protein-structure prediction pathway \citep{riniker2015etkdg,boothroyd2023sage,eastman2024openmm8,cerutti2024stormm,devereux2020ani2x,ahdritz2024openfold}. These methods can be compared or chained without implying that their energies, assumptions, or domains are interchangeable. Appendix~\ref{app:architecture} specifies the implemented numerical pathways and their validation boundaries; Appendix~\ref{app:manual} records how those pathways are exposed, invoked, and constrained in the current workbench.

The same principle applies to workflow provenance. Reproducible computational science requires retention of the data, computational process, execution environment, and provenance that produced a result \citep{cohenboulakia2017workflows,wilkinson2016fair}. More recent work formalizes workflow invariants and testable reproducibility criteria rather than treating provenance as passive metadata \citep{pritchard2025reproducibility}. Generated adaptive software increases the importance of this requirement because the workflow itself may be created or modified during the scientific interaction rather than existing as a single release version defined in advance.

This leads to a stronger requirement: generation of a scientific workflow should be accompanied by generation of the conditions under which its claims could be falsified. If an agent implements a conformer-search strategy, the output should include the benchmark that tests whether conformational coverage improves, the independent reference used for comparison, the quantitative tolerance that defines agreement, the applicability limits of the test, and explicit failure conditions. If an agent adds a numerical kernel, it should generate the corresponding independent comparison and acceptance criteria. These are better understood as scientific contract tests than as ordinary software unit tests because they define what scientific claim the implementation can support.

\begin{figure}[htbp]
  \centering
  \safeincludegraphics[width=0.94\linewidth]{figures/fig3_build_loop_compact.png}
  \caption{The scientist--agent build loop. Scientific needs and objections are translated into software changes, while each new capability is paired with claim-defining tests, reference evidence, applicability limits, and provenance. The ability to generate a workflow and the ability to specify how its claim could fail are treated as parts of the same scientific result.}
  \label{fig:buildloop}
\end{figure}

Molarium's development repeatedly followed this pattern. A chemically implausible geometry exposed a problem in representation and editing semantics and a questionable energy prompted comparison with an independent implementation. Similarly, a discrepancy between CPU and GPU execution revealed both a unit error and an incorrect assumption about which backend was running. Requesting constrained dynamics led to implementation of SHAKE/RATTLE and comparison against a reference pathway \citep{ryckaert1977shake,andersen1983rattle}. In each case, the useful scientific result consisted of the feature together with a test that established the boundary of the claim.

Fluent agentic software makes this distinction more consequential, since problems are not always obvious and do not necessarily result in errors. For example, a unit error can still produce a smooth trajectory and independent implementations can agree if they inherit the same flawed inputs. Successful execution and an attractive visualization therefore occupy the bottom of an evidence hierarchy rather than establishing that a molecular prediction is correct.

\begin{figure}[htbp]
  \centering
  \safeincludegraphics[width=0.94\linewidth]{figures/fig4_evidence_ladder_fixed.png}
  \caption{An evidence ladder for generated scientific software. Successful execution establishes less than numerical agreement; numerical agreement establishes less than retrospective predictive performance; retrospective performance establishes less than prospective transfer; and prospective accuracy is distinct from measurable improvement in drug-discovery decisions.}
  \label{fig:evidence}
\end{figure}

When software construction is difficult, the labor of building a computational method can obscure the separate obligation to validate it. As implementation becomes easier, the scientific responsibility becomes harder to ignore. The workflow, its provenance, and the conditions that could falsify its claims should therefore be treated as a single scientific object.

\section{Local computation changes the interaction model}

A web browser cannot always replace production molecular simulation software or high-performance computing. Large periodic systems, long trajectories, and large-scale screening studies can require substantially more compute resources than a personal machine can provide. Molarium currently addresses calculations that are sufficiently bounded to remain inside the scientist's interactive design session, although remote job submission to a HPC cluster or the cloud for specific calculations would be trival to add.

The practical benefit of local compute is the reduced decision latency. Consider a medicinal chemist examining a bound ligand and asking whether a proposed three-dimensional edit creates an avoidable steric problem or opens a plausible alternative conformation of the protein binding pocket. In a fragmented workflow, the scientist may export coordinates, request or run a separate calculation, move results back into a viewer, and reconstruct the context in which the question arose. In Molarium, the editing, relaxation, conformational exploration, and trajectory inspection all run locally within the same session, tracking the molecular state, the calculation, associated parameters, the design question, and decisions.

Local execution also changes economics for these interactive calculations. The marginal cost of many small calculations approaches the cost of hardware already on the user's desk, without requiring a remote orchestration cycle or a separate software license for each supported operation. This advantage is most relevant in the rapid design loop, where a calculation may be useful because it answers a focused question in minutes rather than because it achieves the maximum possible throughput on a benchmark.

The lower operational barrier increases the importance of defining the project questions appropriately and choosing calculations deliberately. A medicinal chemist deciding between two substituent orientations may gain value from a constrained local refinement that changes which analog is synthesized next, while running hundreds of additional simulations after the decision is already insensitive to the result adds computational activity without increasing scientific leverage. The operational efficiency helps the scientist focus on the right question, apply the appropriate method, and prioritize molecules that have a realistic chance to advance the program.

Local execution also redistributes control. The scientist can determine which analysis should exist, how data should be juxtaposed, and which calculation should be available at the point of decision. Molarium illustrates this shift: a scientist did not need to wait for every useful workflow to appear on a conventional product roadmap. Missing interfaces, algorithms, and methods can be added while the underlying question was still active.

\section{What becomes scarce when the tool layer becomes cheap}

A concrete medicinal-chemistry example helps separate the layers of value. Suppose a project team is considering a substitution intended to improve a ligand's interaction with a newly identified pocket while preserving permeability. The tool layer provides the editor, visualization, coordinate handling, workflow logic, and connection to the relevant calculations. The model layer determines whether the predicted pose, conformational preference, interaction energy, or permeability estimate is credible. The discovery-intelligence layer supplies the project-specific context: prior SAR, assay conditions, structural interpretation, known liabilities, synthetic constraints, and the reason the team is considering that particular design at that point in the program. Coding agents can make the first layer easier to construct. The value of the second and third layers depends on scientific accuracy and accumulated knowledge rather than interface complexity.

The tool layer includes graphical interfaces, visualization, format conversion, wrappers, workflow logic, reports, job control, data plumbing, and bounded numerical implementation. These capabilities remain important and require maintenance, but their marginal reproduction cost is falling quickly as coding agents improve. A scientist can increasingly generate a project-specific view or connect an existing computational method without waiting for the corresponding feature to appear in a commercial release.

The model layer includes force fields, quantum-mechanical methods, docking and scoring functions, free-energy models, machine-learned potentials, protein-structure predictors, molecular representations, property models, parameters, weights, uncertainty calibration, and the evidence that establishes their useful domain. Easier invocation does not improve predictive quality. A docking model with an inadequate scoring function remains inadequate when an agent invokes it. A systematic force-field error can become more damaging when automated workflows propagate it across thousands of calculations. A machine-learned property model does not gain extrapolative validity because a language model describes its output fluently, and a protein-structure predictor does not become reliable for a particular conformational state because it is available through a better interface.

Differences among models therefore become more consequential as the mechanics of invoking them become less scarce. Drug discovery frequently requires extrapolation into new chemistry, structural states, and biological contexts rather than interpolation within a well-sampled training distribution \citep{bender2021realistic1,bender2021realistic2,wigh2022representations}. Faster access can accelerate good models and weak models alike.

The discovery-intelligence layer includes proprietary experimental data, negative results, project history, assay context, structural interpretation, design rationale, internal protocols, expert judgment, and program objectives. Molarium intentionally contains little of this layer because it is a public workbench built primarily from public scientific components. The same architecture could be connected to internal structures, assay data, pharmacokinetics, target biology, medicinal-chemistry reasoning, and program objectives. That project-specific context can determine which computation is relevant and how much weight its output should receive.

\begin{figure}[htbp]
  \centering
  \safeincludegraphics[width=0.94\linewidth]{figures/fig5_value_layers_fixed.png}
  \caption{Three interdependent layers of value in generated adaptive molecular design. Implementation and orchestration are becoming increasingly reproducible by coding agents. Scientific models remain differentiated by accuracy, representation, domain coverage, calibration, and validation. Proprietary discovery intelligence combines experimental data with the context and decision history that make those data useful.}
  \label{fig:value}
\end{figure}

Versioned molecular history provides a bridge into discovery intelligence. Traditional compound databases record what was made and measured well, but they rarely preserve what was considered, why one design was preferred, which model or structure influenced the decision, or what evidence was visible at the time. A persistent action history can capture that decision trajectory.

For example, a chemist might propose adding a substituent after inspecting a pocket in a protein--ligand structure. The project record could retain the parent structure, the molecular edit, the structural hypothesis, the conformational or energetic calculations viewed before synthesis, competing designs that were rejected, and the subsequent experimental measurements. If the design later fails because an assumed interaction was absent or a liability dominated the profile, that information remains associated with the decision that created the molecule. Future scientists and agents can then ask which assumptions repeatedly led to productive designs and which repeatedly failed. This is decision provenance: preserving the relationship among hypotheses, interventions, predictions, decisions, and outcomes rather than storing only final compounds and assay values.

Agentic systems provide a mechanism to operationalize some of this accumulated knowledge through scientific skills. Here, a skill means a version-controlled, machine-executable procedure that combines a defined tool or model with an organization's accepted inputs, assumptions, validation criteria, applicability limits, failure handling, and escalation rules. A protein-preparation skill might encode how crystallographic waters, protonation states, alternate conformations, unusual residues, and ambiguous cases are handled. A free-energy skill might define accepted perturbations, convergence criteria, diagnostics, and failure modes. The durable asset is the scientific procedure and its validation, rather than the natural-language prompt used to invoke it.

\section{The changing role of scientific expertise}

As software becomes easier to construct and operate, scientific judgment becomes more valuable. Traditional computational expertise includes substantial operational knowledge: preparing inputs, locating settings, translating files between programs, recovering failed jobs, and reproducing standard analyses. An increasing fraction of this work can be encoded in software and agents. The broader automation literature has long recognized that reducing routine operation can increase the importance of human expertise when systems encounter ambiguous or abnormal conditions \citep{bainbridge1983ironies}.

The highest-value expertise shifts toward selecting the right question, choosing an appropriate level of model fidelity, defining acceptable tradeoffs between cost and accuracy, recognizing implausible results, selecting meaningful controls, and determining which evidence would change a project decision. Scientists closest to the underlying discovery problem---whether medicinal chemists, structural biologists, mechanistic biologists, pharmacologists, or computational scientists---can increasingly interact directly with relevant data and calculations rather than routing every question through a software intermediary.

This shift changes the leverage of computational specialists. Less specialist time needs to be spent constructing repetitive interfaces, transferring files, or manually operating routine workflows. More can be spent on model development, molecular physics, parameterization, sampling, uncertainty, validation, and difficult edge cases where established workflows fail. Once that expertise is made explicit, part of it can be encoded into validated scientific skills that broaden access without implying that the underlying science has become simple.

Molarium's development illustrates this division of labor. The coding agent could inspect a repository, refactor software, implement numerical kernels, add tests, and repeat the loop rapidly. The scientist had to recognize chemically implausible results, decide whether a reference was genuinely independent, determine whether a discrepancy mattered, choose acceptable approximations, and identify what evidence would change a decision. The agent compressed the distance between recognizing a scientific need and changing the software; scientific judgment determined which changes mattered.

The bottleneck therefore moves toward epistemology: deciding which approximation is acceptable, whether a faster method sacrifices information that matters to the decision, when an independent control is required, whether a reference is truly independent, whether a model is operating outside its validated domain, what uncertainty must remain visible, and which experiment would discriminate between competing hypotheses. Faster implementation increases the value of asking the right question.

\section{Limitations and outlook}

Molarium is a case study rather than a controlled experiment in software productivity or drug-discovery performance. Its development involved a small number of scientists, one evolving codebase, and a particular coding-agent workflow. The project did not include a randomized comparison with a conventional software-development team, and it cannot establish a universal productivity multiplier or determine what fraction of commercial scientific software can be regenerated reliably.

The open-source repository can nevertheless support a useful community experiment. Contributors can identify computational capabilities described in papers or public specifications, attempt independent agent-assisted reconstructions, test them against appropriate references, and record which parts can and cannot be reproduced reliably. Over time, this could become a living benchmark for the regeneration of molecular-modeling software, with success defined by scientific contract tests rather than by code generation alone.

Molarium's scientific scope is deliberately bounded. Current browser-native molecular mechanics does not reproduce the full capability of mature native simulation environments. General periodic simulation, particle-mesh Ewald electrostatics, barostats, membranes, unrestricted protein parameterization, arbitrary plugins, and other advanced capabilities remain outside the current browser-native scope. ANI-2x remains valid only within its declared chemical domain, and browser-based OpenFold support is narrower than general production protein-structure prediction workflows. Generated wrappers and easier access must preserve, expose, and enforce these applicability limits rather than implying broader scientific validity. Appendix~\ref{app:architecture} catalogues the demonstrated implementation and numerical boundaries, and Appendix~\ref{app:manual} provides the corresponding operational status, failure modes, and end-user limits.

A second limitation concerns the distinction between implementation fidelity and predictive validity. If a browser implementation reproduces an OpenMM reference energy and force to a specified tolerance, that result establishes that the implementation is executing the declared calculation correctly for the tested case. It does not establish that the shared force field predicts the experimental behavior of a new molecule. Likewise, a conformer workflow can reproduce its intended algorithm and explore more structures without improving recovery of biologically relevant conformations. Implementation validation and model validation answer different questions and require different evidence.

The future suggested by Molarium is one in which the interface between scientist and computation becomes more fluid while the scientific state, models, assumptions, evidence, and decision history become more explicit. Project teams can construct temporary views around specific questions; internal models can coexist with open-source methods and commercial services; and the user experience no longer needs to be owned by the provider of each underlying calculation.

The durable infrastructure in such an environment is concrete. A user should be able to inspect the active molecular state and its history, the exact model and parameter versions used for each calculation, declared applicability limits, validation status, assumptions, provenance, and links between predictions and subsequent decisions. These quantities should remain available regardless of which generated view is open. An adaptive interface can change rapidly because the scientific contract underneath it remains stable and inspectable.

\section{Conclusion}

Molarium began as a request for a better environment in which to inspect and manipulate molecules. It developed rapidly into a local-first molecular-design workbench incorporating topology-aware editing, protein preparation, conformer exploration, molecular mechanics, molecular dynamics, learned potentials, protein-structure prediction, browser-native GPU computation, provenance, and reference validation. The speed of that development is technically interesting, but the broader lesson concerns the architecture of scientific software and the location of durable value.

Software can increasingly be generated around the scientific question and the user's specific preferences. The quality of the underlying models, validation, and decision history becomes more consequential as the interface becomes cheaper and more fluid. The scarce assets shift toward trustworthy models, proprietary discovery intelligence, and the human judgment required to know what to ask and what to believe.

Molarium also suggests that molecular-design environments should preserve more than the latest structure. A versioned record of molecular states, user actions, calculations, assumptions, and results can capture the trajectory by which scientific decisions were made. Coupled to explicit scientific contracts and falsification tests, that history can make generated adaptive software more reproducible, extensible, and scientifically accountable.

When software is difficult to construct, implementation effort can hide the distinction between building a method and validating it. When implementation becomes easy, the remaining scientific obligations become harder to ignore. The opportunity is to place adaptable computational capability directly in the hands of scientists who understand the problem while preserving the models, evidence, and decision history that determine whether the answer should be trusted.

The tool layer is becoming cheap and adaptive. Scientific advantage moves toward better models, better discovery intelligence, and better decisions.

\section*{AI assistance and availability}

GPT-5.6 Sol was used extensively as a coding agent during the development of Molarium under human scientific direction, inspection, testing, and revision. Language-model assistance was also used in restructuring and drafting this manuscript. The authors remain responsible for the scientific claims, validation, source attribution, and final text.

Molarium is available at \url{https://molarium.org}. Its source code and technical documentation are publicly available under the Apache License 2.0 at \url{https://github.com/rafwiewiora/molarium}. Bundled third-party software, model parameters, force fields, and scientific data retain their separately identified terms. Quantitative benchmark results are scoped by the relevant implementation state, fixtures, model and asset hashes, browser and hardware environment, numerical precision, and timing protocol. Appendix~\ref{app:architecture} provides the architecture, numerical validation, and executable design-history evidence underlying these results; Appendix~\ref{app:manual} records the inspected operating surface, module boundaries, failure behaviour, and export/persistence contracts.

\appendix

\appendixsection{Molarium architecture, numerical validation, and executable design history}{app:architecture}
\begin{small}
\setlength{\emergencystretch}{2em}

This appendix specifies the implementation and validation substrate underlying the architectural claims in the main text. It defines the molecular-state contract shared across chemistry perception, parameterization, numerical execution, visualization, and provenance; describes the OpenMM/WebAssembly and WebGPU execution paths; reports fixed-input numerical validation and workload-dependent performance measurements; and documents the content-addressed replay and decision-history mechanisms used by Molarium. The final section provides an executable SOS1 design-history example. Numerical agreement is treated as implementation validation of a declared calculation and is kept distinct from validation of the underlying physical or statistical model.

\subsection{Execution architecture}

Molarium maintains a single molecular state keyed by stable atom identifiers across ingestion, chemical perception, parameterization, calculation, visualization, and replay. RDKit provides graph interpretation, conformer generation, SMARTS matching, and ligand-level preparation. A JavaScript SMIRNOFF parameterizer applies OpenFF Sage 2.1.0 to generate an explicit numerical recipe containing particle properties, bonded terms, nonbonded terms, exceptions, and constraints. Preparation fails on missing required parameters rather than substituting inferred terms. The accepted recipe is consumed by the OpenMM/WebAssembly reference path and the direct or replica-parallel WebGPU paths, allowing backend comparisons on identical topology, coordinates, parameters, and atom ordering.

The primary data path is:

\begin{center}\small
PDB or SMILES input\\
$\downarrow$\\
molecular graph, stable atom identities, and coordinates\\
$\downarrow$\\
RDKit interpretation $\rightarrow$ OpenFF Sage numerical recipe\\
$\downarrow$\\
OpenMM/WebAssembly reference $\;|\;$ direct WebGPU $\;|\;$ replica-parallel WebGPU\\
$\downarrow$\\
viewer, action audit, calculation record, and design history
\end{center}

The same molecular object remains authoritative as calculations return coordinates, forces, energies, or trajectories. Table~\ref{tab:appA_browser_impl} summarizes the current implementation.

\begin{table}[htbp]
\centering\small
\caption{Current browser implementation.}
\label{tab:appA_browser_impl}
\begin{tabularx}{\linewidth}{@{}p{0.20\linewidth}p{0.24\linewidth}X@{}}
\toprule
\textbf{Task} & \textbf{Software used} & \textbf{Implementation role}\\
\midrule
Molecular interpretation & RDKit 2025.03.4, WebAssembly & SMILES/graph handling, ETKDGv3, SMARTS matching, MMFF94, explicit UFF fallback, and Gasteiger charges\\
Force-field assignment & OpenFF Sage 2.1.0 & Generates particles, bonds, angles, torsions, nonbonded terms, exceptions, and constraints\\
Numerical reference & OpenMM 8.2.0 Reference, WebAssembly & Double-precision energies and forces, L-BFGS minimization, and short dynamics used for fixed-input reference comparisons\\
Direct local mechanics & WebGPU/WGSL, 32-bit floating point & Single-system force/energy evaluation, OBC2/ACE, minimization, and short dynamics\\
Replica-parallel mechanics & STORMM-inspired WebGPU engine & Advances independent replicas sharing one topology and parameter layout; paired-word fixed-point accumulation; current limit 512 atoms per replica\\
State and provenance & JavaScript modules, background workers, Canvas, Web Crypto & Preserves atom identity and records action, calculation, replay, and SHA-256-linked provenance\\
\bottomrule
\end{tabularx}
\end{table}

Internal units are nm, ps, daltons, kJ/mol, and kJ/mol/nm; the interface presents angstroms and kcal/mol where appropriate. Numerical engines execute in reusable Web Workers so rendering and user interaction remain responsive. Requests carry job identifiers, and failed workers are discarded and recreated. Results return as compact coordinate, force, energy, or trajectory arrays mapped to the originating atom identities.

For ordinary small molecules, the browser parameterizer applies Sage 2.1.0 patterns and labelled Gasteiger charges \citep{boothroyd2023sage,gasteiger1980charges}. Prepared Rosemary validation fixtures retain their exported numerical recipes and NAGL charges. CPU/GPU comparisons therefore operate on the same prepared graph and numerical inputs rather than independently parameterized systems.

\subsubsection{OpenMM WebAssembly reference}

OpenMM 8.2.0 is cross-compiled with Emscripten 6.0.6-git and CMake 3.31.0 \citep{eastman2017openmm7,eastman2024openmm8}. A thin C interface constructs an OpenMM \texttt{System} from the serialized numerical recipe and returns energies, forces, minimized coordinates, or short trajectories. Source, patches, interface code, and built assets carry recorded fingerprints so the deployed reference path can be associated with a specific implementation state.

The reference configuration is nonperiodic and uses the double-precision OpenMM Reference platform. Supported operations include energy and force evaluation, L-BFGS minimization, short dynamics, optional X--H constraints, 2 fs Langevin-middle integration, a cutoff configuration used for validation, and OBC2/ACE with mbondi2 radii \citep{onufriev2004gb}.

The WebAssembly build exposed a threading-specific setup failure for constrained dynamics: native OpenMM creates helper CPU threads during constraint initialization, whereas the browser build is single-threaded. A serial patch removes the wait path for this build without modifying native OpenMM.

Fixed-input parity was evaluated in two stages. The same C interface compiled natively and for WebAssembly agreed to a maximum energy difference of $2.84 \times 10^{-14}$ kcal/mol and relative force RMS of $3.76 \times 10^{-15}$. Five hash-selected protein--ligand poses were then prepared through browser Sage and evaluated with matched numerical inputs. Maximum energy differences were $1.24 \times 10^{-4}$ kcal/mol in vacuum and $2.76 \times 10^{-4}$ kcal/mol with OBC2; maximum relative force RMS error was $1.16 \times 10^{-5}$ in both modes. These tests establish implementation parity for the specified fixtures and protocols.

\subsubsection{Direct WebGPU mechanics}

The direct WebGPU path evaluates the prepared force-field equations in WGSL using 32-bit floating-point arithmetic. Force ownership is atomwise: one invocation owns the output force for one atom, and compact incidence lists identify the bonded terms that contribute to that atom. Pair and valence terms are recomputed where required so each invocation writes its force once, avoiding floating-point atomic updates, which WGSL does not provide \citep{webgpuWGSL2026}.

The deployed nonbonded representation stores excluded and specially scaled pairs explicitly and mixes ordinary Lennard--Jones parameters on demand. An earlier all-pairs metadata representation required $32N^2$ bytes and imposed a practical allocation ceiling; the sparse exception representation removes that quadratic metadata cost. OBC2/ACE is evaluated in three kernel passes for Born radii, coordinate derivatives, and solvent/nonpolar energy and force accumulation. Optional X--H constraints are graph-coloured so non-overlapping SHAKE/RATTLE corrections can execute concurrently \citep{ryckaert1977shake,andersen1983rattle}. Seeded Langevin dynamics supports 1 fs flexible and 2 fs constrained integration.

Pocket relaxation uses a bounded steepest-descent procedure with zero inverse mass on fixed receptor atoms; only the ligand and explicitly selected pocket atoms move. Current operating limits are listed in Table~\ref{tab:appA_webgpu_limits}.

\begin{table}[htbp]
\centering\small
\caption{Operating limits in the current direct WebGPU worker.}
\label{tab:appA_webgpu_limits}
\begin{tabularx}{0.95\linewidth}{@{}p{0.32\linewidth}X@{}}
\toprule
\textbf{Setting} & \textbf{Current value}\\
\midrule
Principal force workgroup & 64 atomwise calculations\\
Time step & 1 fs flexible; 2 fs with X--H constraints\\
Constraint solve & Four graph-coloured sweeps by default; at most 32\\
Minimization & 750 iterations by default; at most 2,000; each displacement capped at 0.001 nm\\
Direct dynamics request & At most 5,000 steps; longer requests use the ensemble route\\
Validation cutoff & 1.0 nm cutoff; 0.2 nm skin; rebuild every 20 steps\\
\bottomrule
\end{tabularx}
\end{table}

At the interactive system sizes tested here, a conventional Verlet neighbour list did not improve performance. With a 0.2 nm skin rebuilt every 20 steps, 5,000 constrained OBC2 Trp-cage steps were 46\% slower and 1,000 ubiquitin steps were 17\% slower than full-range nonbonded evaluation on an Apple M1 Pro. The cutoff route is therefore retained for validation but is not exposed in the product controls.

GPU benefit depends on amortizing dispatch and data movement. A single Trp-cage energy/force evaluation required 25.5 ms in WebGPU and 5.7 ms in OpenMM Reference. Across 1,000 dynamics steps, WebGPU reached 214.9 versus 102.8 ns/day for 304-atom Trp-cage and 65.18 versus 7.11 ns/day for 1,231-atom ubiquitin. These measurements use warmed, nonperiodic, all-pairs workloads. The direct path does not currently implement periodic boundaries, particle-mesh Ewald electrostatics, explicit solvent, barostats, virtual sites, or arbitrary OpenMM custom forces.

\subsubsection{Replica-parallel WebGPU}

The replica engine targets throughput across independent copies of a shared topology. One replica maps to one GPU workgroup, in contrast to the atomwise decomposition used by the direct engine. The memory layout and fixed-point accumulation strategy draw on STORMM's stacked-system design, while the browser implementation is independent WGSL code \citep{cerutti2024stormm}. Replicas share topology and parameters but do not interact.

Concurrent accumulation uses paired 32-bit words to represent signed 64-bit fixed-point values because WGSL exposes 32-bit integer atomics. The implementation uses $2^{22}$ counts per kcal/mol for energies, $2^{18}$ counts per kcal/mol/\AA{} for forces, and $2^{32}$ counts per \AA{} for coordinates. This replaced 32-bit accumulators that overflowed silently in highly clashed configurations and produced bit-for-bit replay for identical seeds on the tested adapter.

Integration uses velocity Verlet, seeded Langevin thermalization, and SHAKE/RATTLE. Constraint projection is serial within a replica while replicas execute in parallel; the default relative tolerance is $10^{-5}$ with at most 32 iterations. Commands contain at most 50 steps. Persistent status words convert non-finite values, accumulator overflow, coordinate clamps, and failed constraints into terminal errors. Current limits are 512 atoms per replica, all-pairs nonbonded evaluation, and identical atom count, topology, and parameters across replicas.

Force evaluation uses a separate 32-bit floating-point coordinate working copy. Translating a test system 500 \AA{} from the origin changed energies by approximately $3 \times 10^{-4}$ relative, identifying the working coordinate representation rather than the fixed-point store as the remaining translation-sensitivity boundary. Eliminating this source of error requires recentering before force evaluation or a wider working representation.

Replica throughput also crosses over with workload size. In warm tests, OpenMM Reference was 13.3-fold faster for one C16 molecule and 3.5-fold faster for one 27-water system. WebGPU was 10.3-fold faster for 1,024 C16 replicas and 24.1-fold faster for 256 water replicas. These measurements compare aggregate replica throughput; the two engines do not use identical integrators.

\subsection{Reference-driven implementation and validation}

The numerical implementation is developed against a shared executable specification. OpenMM and the browser kernels receive the same molecular graph, atom ordering, coordinates, parameters, units, exceptions, constraints, and named energy terms. This enables term-resolved differential testing rather than comparison of independently prepared end states. The implementation loop is:

\begin{enumerate}[leftmargin=*]
\item freeze the molecular graph, atom ordering, coordinates, force-field parameters, and units used by the reference calculation;
\item implement one energy/force term and its memory representation at a time;
\item evaluate the identical numerical input with OpenMM;
\item compare named energy components and Cartesian force components;
\item localize the first discrepancy, correct it, and retain the minimal reproducing case as a regression test; and
\item execute the resulting WebGPU and WebAssembly paths in the browser to include worker, adapter, serialization, and UI integration boundaries.
\end{enumerate}

Coding agents operate inside this harness by modifying implementation code, tests, and adapters. The physical-model equivalence criteria---including atom order, solvent model, cutoffs, constraints, units, charge provenance, and acceptable numerical tolerances---are specified explicitly and remain prerequisites for interpreting any backend comparison.

Several observed failures led directly to architectural changes (Table~\ref{tab:appA_failures}).

\begin{table}[htbp]
\centering\small
\caption{Observed failures that changed the implementation.}
\label{tab:appA_failures}
\begin{tabularx}{\linewidth}{@{}p{0.26\linewidth}p{0.36\linewidth}X@{}}
\toprule
\textbf{Observation} & \textbf{Design change} & \textbf{Retained check}\\
\midrule
Pair metadata grew as $32N^2$ & Store only exceptional pairs; mix ordinary pairs on demand & Large-system allocation and parity tests\\
32-bit force/energy sums overflowed & Paired-word 64-bit fixed-point accumulation and fault flags & Clashed system at approximately 2 million kcal/mol\\
Coordinates wrapped beyond 128 \AA{} & 64-bit fixed-point coordinate store plus floating-point working copy & Whole-system translation to 500 \AA{}\\
Standard cutoff reduced performance & Retain for validation; omit from product controls & Fixed Trp-cage/ubiquitin timing protocols\\
Apparent engine mismatch & Verify atom order, input hashes, and solvent mode before numerical debugging & Native/WebAssembly parity and matched-solvent comparisons\\
\bottomrule
\end{tabularx}
\end{table}

Regression coverage is organized around the interfaces where failures were observed: equation-level fixtures for bonded and nonbonded terms; representation tests for clashes, constraints, overflow, and translated coordinates; backend comparisons on identical prepared inputs; protein fixtures that exercise special pairs and worker transfer; browser tests for adapter support and worker recovery; and fixed performance protocols that retain slower but numerically valid designs. Fixed random seeds and stable atom identifiers make failing cases repeatable. The replica engine terminates on fixed-point overflow or unconverged constraints; the direct engine records final constraint error for comparison against an acceptance threshold.

The demonstrated portability boundary is fixed-topology, nonperiodic molecular mechanics with regular array state and a small number of synchronization points. WebGPU is effective when work decomposes into local or pairwise operations and repeated evaluations amortize dispatch. Algorithms requiring 64-bit floating-point arithmetic, frequent device-wide synchronization, rapidly changing graphs, large adaptive sparse structures, or mature FFT/sparse-linear-algebra libraries require a different decomposition under current WGSL \citep{webgpuWGSL2026}. Particle-mesh Ewald electrostatics, heterogeneous reactive systems, and some long-timescale production protocols are therefore outside the demonstrated browser-native boundary. For those workloads, the architecture supports escalation to a WebAssembly reference, hybrid local path, or external specialist engine while preserving the originating molecular state and protocol provenance.

\subsection{Workflow orchestration and local execution boundary}

Molarium distinguishes local calculation from offline operation. In normal connected mode, supported numerical calculations execute on the user's device, while explicit PDB/CCD retrieval or sequence-alignment searches contact the services identified in the interface. Local Lab serves a reviewed, versioned checkout from localhost under a restrictive Content Security Policy. External structure and alignment controls are disabled, and an egress-canary test fails if a synthetic private-data marker reaches an external request. A hosted application can change the assets served on a subsequent visit; the stronger reproducibility and privacy boundary is a reviewed checkout with locally pinned assets.

Current workflow surfaces include:

\begin{itemize}[leftmargin=*]
\item \textbf{Local small-molecule design.} Synchronized 2D/3D construction, graph validation, RDKit or direct-WebGPU calculations, aligned trajectory/energy inspection, and export. OpenMM/WebAssembly is used for internal validation and fixed-scaffold operations rather than as a selectable Simulate-panel method.
\item \textbf{Protein--ligand pose propagation.} A trusted pose is captured with stable atom lineage; changed ligand branches are sampled while the receptor remains fixed; optional contact hypotheses constrain search; and a selected pose can be applied and relaxed. An experimental induced-fit action releases the ligand and all atoms of protein residues within a 6 \AA{} pocket shell while outer atoms remain fixed.
\item \textbf{Agent-operated design.} Agents invoke the versioned Chemist Actions API rather than arbitrary page functions. The API exposes the same bounded state-changing verbs available to a chemist---including select, edit, search, apply, relax, and undo---plus read-only inspection. Requests execute serially and record sequence, timestamps, arguments, status, and duration.
\end{itemize}

Two extension paths are represented in the architecture but are not complete product workflows. A local-to-specialist path can export a fingerprinted checkpoint and explicit protocol for periodic solvent, long trajectories, free energy, or electronic-structure calculations, with returned results attached as evidence to the originating branch; transport, scheduling, signatures, and trust policy are not implemented. A collaborative review model can branch by chemotype, protonation state, pose, or author and record a selected branch with explicit parentage; coordinate averaging is not used as a merge operation.

\subsection{Content-addressed molecular and decision provenance}

Molarium separates operational, numerical, and decision provenance. A molecular snapshot records one explicit state. Content-addressed commits link that state to parent states, action scripts, hypotheses, and evidence. Branches preserve alternative design routes, and decisions record which route progressed, stopped, or was deferred. These semantics are implemented as a custom content-addressed ledger rather than a Git repository or a general structural merge algorithm.

The current live application retains action records in memory; clearing the scene clears that audit. Durable provenance therefore requires explicit export of API moves or a replay log. The append-only design-campaign ledger, commit, branch, decision, and verification functions exist in the repository and support standalone builders and reviews, while one-click promotion of a live session into ledger commits is not currently exposed in the main interface.

\begin{table}[htbp]
\centering\small
\caption{Separation of operational, numerical, and decision provenance.}
\label{tab:appA_provenance}
\begin{tabularx}{\linewidth}{@{}p{0.24\linewidth}p{0.27\linewidth}X@{}}
\toprule
\textbf{Record} & \textbf{Question} & \textbf{Contents}\\
\midrule
Chemist Actions audit & What was requested? & Public action, arguments, timing, status, and compact after-state\\
Calculation record & What numerical protocol ran? & Inputs, method, parameters, outputs, and hash-linked events\\
Design campaign & Why did a state advance or stop? & Snapshots, branches, parents, hypotheses, evidence, and decisions\\
\bottomrule
\end{tabularx}
\end{table}

Canonical JSON and SHA-256 fingerprints bind records to exact serialized representations. The \path{molarium.design-campaign/v1} schema stores snapshots, commits, branches, events, and decisions. Each commit identifies its selected snapshot, parent commit(s), optional action script, hypotheses, evidence, tags, and branch. Events are append-only and contain the previous event fingerprint, actor class (human, agent, system, or import), and source. Supported decision states are \texttt{progressed}, \texttt{not-progressed}, \texttt{deferred}, \texttt{failed}, \texttt{duplicate}, \texttt{superseded}, and \texttt{archived}. Finalization appends a completion event and hashes the campaign; subsequent record changes fail verification. Hashes provide alteration detection for the recorded representation, not authorship, chemical equivalence, or scientific validity.

Registered design routes use a separate \path{molarium.registered-design-route/v1} schema that specifies the coordinate boundary, hash-pinned starting state, and ordered graph edits. The SOS1 route uses this schema. Route actions are \texttt{designRoute.load}, \texttt{designRoute.applyStep}, and \texttt{designRoute.inspect}; campaign-ledger operations remain separate. Cross-schema validation rejects a route presented as a campaign ledger or a campaign ledger presented as a route. Frozen predictions made before the schema split retain their original input hash, with migration records storing both hashes and accepting only the schema-identifier substitution that reconstructs the frozen bytes.

A decision-to-state linkage has the following form; full identifiers are SHA-256-derived values rather than the abbreviated labels shown here.

\begin{lstlisting}
{
  "moleculeCommit": {
    "snapshotId": "snapshot:<selected-child>",
    "parents": ["commit:<parent>"],
    "actionScriptId": "script:<designer-moves>",
    "hypothesisIds": ["event:<pocket-must-open>"]
  },
  "decisionEvent": {
    "kind": "decision.recorded",
    "actorId": "chemist.01",
    "payload": {
      "targetCommitId": "commit:<selected-child>",
      "disposition": "progressed",
      "reasonCodes": ["contact-loss"],
      "rationale": "Phe-out removes growth-induced clashes.",
      "evidenceIds": ["event:<rotamer-search>"]
    }
  }
}
\end{lstlisting}

Executable replays use \texttt{molarium.chemist-action-script/v1}. Scripts contain ordered calls to the allow-listed Chemist Actions API rather than arbitrary code or direct coordinate replacement. Designer Moves JSON files must be smaller than 2 MB; each serialized API request and its argument object are capped at 8 MiB for ASCII JSON, eight nested levels, and 2,048 JSON nodes. Actions can name returned stable identifiers for reuse by later steps. The runner executes one action at a time through \texttt{api.execute(\{action,args\})}, derives a repeatable request ID from the script fingerprint, and stops on the first failure. Replay results use \texttt{molarium.chemist-action-replay/v1} and record the source fingerprint and each step outcome. Presentation actions can alter focus or styling but cannot change molecular graph or coordinates except through the same allow-listed state-changing actions.

The SOS1 replay is identified by a stable public route and by exact release and content fingerprints. A human-readable URL can resolve to a newer release, whereas the release commit, route-file fingerprint, and canonical action-script fingerprint identify the reviewed bytes and normalized JSON content. The live workflow can capture a starting state, record exploration actions, apply and relax selected states, and export action and numerical records. Formal campaign commits and decision events are represented by the ledger layer described above.

\subsection{SOS1 hit-to-BAY-293 executable replay}

The SOS1 example follows the analogue series reported by Hillig and colleagues \citep{hillig2019sos1}. The supplied design information consisted of the 5OVE/AXE coordinate-bearing starting complex and the published analogue graphs for AWT, AWZ, AWW, and AXH. The later analogue poses and receptor conformations were withheld during prediction. Each analogue was propagated from the preceding frozen prediction rather than from its corresponding later crystal structure.

The selected route contains 33 recorded API actions: 29 molecular or protocol operations and four workspace switches. The full exploration audit contains 89 completed API calls, including three Phe890/pose branches, nine undo calls, and deterministic reselection. The rendered interface sequence adds 16 display-only actions for 49 replay steps. Later crystal structures were opened only after four prediction checkpoints and the exploration audit had been hashed. Display and focus actions are non-mutating.

The analogue transitions involve substantial graph rewrites while preserving explicit atom lineage (Table~\ref{tab:appA_lineage}).

\begin{table}[htbp]
\centering\small
\caption{Heavy-atom lineage across the supplied SOS1 analogue graphs.}
\label{tab:appA_lineage}
\begin{tabular}{@{}lrrrr@{}}
\toprule
\textbf{Transition} & \textbf{Retained} & \textbf{Deleted} & \textbf{Added} & \textbf{Candidate mappings}\\
\midrule
AXE--AWT & 23 & 4 & 6 & 4\\
AWT--AWZ & 19 & 10 & 12 & 1\\
AWZ--AWW & 16 & 15 & 15 & 1\\
AWW--AXH & 15 & 16 & 17 & 1\\
\bottomrule
\end{tabular}
\end{table}

Each ``64-chain'' search comprises 64 independently seeded local pose searches, distributed over eight browser workers in the recorded replay. Parameterization without motion assigns the numerical force-field representation while preserving coordinates; only the subsequent relaxation changes coordinates. The SOS1 route uses whole-system Sage 2.1.0 with deterministic Gasteiger charges. AWT, AWZ, AWW, and AXH contain 7,929, 7,933, 7,935, and 7,937 particles, respectively. This charge assignment is part of the recorded protocol and is distinct from standard Sage workflows using AM1-BCC or NAGL charges.

Before flexible relaxation, Molarium stores ligand bond lengths and rejects a relaxed pose if any heavy-atom bond is both more than 0.45 \AA{} longer and 1.25-fold its assigned target, or both more than 0.35 \AA{} shorter and below 0.75-fold its target. Rejected relaxation restores the previous pose rather than accepting an improved score with a broken molecular geometry.

\begin{table}[htbp]
\centering\small
\caption{Sequence of molecular states in the SOS1 replay. Predictions, rather than future crystal structures, seed the next step.}
\label{tab:appA_sos1_beats}
\begin{tabularx}{\linewidth}{@{}p{0.07\linewidth}p{0.25\linewidth}X@{}}
\toprule
\textbf{Step} & \textbf{Operation} & \textbf{State/evidence}\\
\midrule
0 & Load and prepare & 5OVE/AXE only; capture the compound 1 reference and coordinate boundary\\
1 & Rewrite to compound 17 (AWT) & Run 64 independent reference-guided pose searches; apply selected pose; parameterize; relax\\
2 & Rewrite to compound 18 (AWZ) & Propagate from predicted AWT, not 5OVF; apply, parameterize, relax\\
3 & Grow compound 21 (AWW) & New ring enters the Phe890-in volume; fixed-core representative gains four clashes\\
4 & Resolve pocket state & Enumerate discrete Phe890 rotamers; select outward branch; retain the moved receptor as the next reference\\
5 & Predict BAY-293 (AXH) & Grow from frozen AWW; 64-chain pose search; apply; parameterize without motion; relax; freeze\\
6 & Evaluate holdouts & Compare AWT/AWZ/AWW/BAY-293 with 5OVF/5OVG/5OVH/5OVI\\
\bottomrule
\end{tabularx}
\end{table}

The Phe890 conformational step is treated as a discrete state search because local minimization does not reliably cross the relevant side-chain torsional barrier. Separate rotamer basins are generated before local ligand--pocket relaxation. Feasible branches are ranked first by clashes introduced by ligand growth, then by pose score and finite relaxed energy. The nominal $-60^{\circ}$ branch adds four clashes above the fixed-core baseline and is rejected; the $-180^{\circ}$ branch adds none and is selected; the $+60^{\circ}$ branch adds two and remains feasible. The outward Phe890 $\chi_1$ basin is therefore selected before 5OVH or 5OVI is revealed.

For post-freeze structural evaluation, predicted receptors are aligned to their holdouts using 23, 25, 25, and 27 pocket C$\alpha$ atoms for AWT through AXH. Ligand RMSD is then calculated over 29, 31, 31, and 32 heavy atoms, minimizing over 2, 1, 1, and 1 graph-symmetry mappings, respectively; the ligand itself is not fitted. Results are reported in Table~\ref{tab:appA_sos1_eval}.

\begin{table}[htbp]
\centering\small
\caption{Post-freeze structural evaluation. Torsions are predicted/crystal in degrees; ligand RMSD is receptor-aligned and graph-symmetry-minimized.}
\label{tab:appA_sos1_eval}
\begin{tabular}{@{}lllll@{}}
\toprule
\textbf{Prediction} & \textbf{Holdout} & \textbf{Ligand RMSD (\AA{})} & \textbf{Phe890 $\chi_1$} & \textbf{Phe890 $\chi_2$}\\
\midrule
AWT & 5OVF & 0.516 & -77.6 / -71.5 & -71.7 / -71.7\\
AWZ & 5OVG & 1.528 & -73.4 / -75.6 & -92.0 / -87.9\\
AWW & 5OVH & 2.137 & -172.2 / +175.8 & +116.1 / +76.3\\
BAY-293 & 5OVI & 1.215 & -172.4 / -166.5 & +108.7 / +71.7\\
\bottomrule
\end{tabular}
\end{table}

The replay recovers the experimentally observed outward Phe890 $\chi_1$ basin before exposure to 5OVH or 5OVI. The terminal $\chi_2$ orientation is less accurate, differing by approximately $40^{\circ}$ for AWW and $37^{\circ}$ for BAY-293. The corresponding ligand RMSDs range from 0.516 to 2.137 \AA{} under receptor-aligned, symmetry-minimized evaluation.

The interface replay executes from a blank canvas through the public Chemist Actions API. Candidate-generation operations leave coordinates unchanged until an explicit Apply action. The replay therefore represents discrete state transitions rather than interpolated trajectories or molecular dynamics. Its 48-action rendered sequence lasts 62.58 s while the underlying audit records approximately 908.5 s of operation durations; timing remains attached to the audit and render manifest, and the rendered sequence is not used as a performance benchmark. Comparison of fragment-bound 6EPM with the 5OVE starting structure shows that the inhibitory Tyr884 arrangement is already present at the coordinate boundary; the hit-only route does not predict a Tyr884 transition.

\end{small}


\appendixsection{Working with Molarium: scientific workflows for humans and agents}{app:manual}
% Appendix B inlined from appendix-b-workflows.tex for a single editable source file.
\begingroup
\setlength{\parskip}{5.5pt}
\setlength{\emergencystretch}{2em}
\renewcommand{\arraystretch}{1.08}
\small

\setlist[itemize]{leftmargin=14pt,labelsep=5pt,itemsep=1pt,topsep=2pt,parsep=0pt}
\setlist[enumerate]{leftmargin=18pt,labelsep=5pt,itemsep=1.5pt,topsep=2pt,parsep=0pt}

\newcommand{\code}[1]{\texttt{#1}}
\newcommand{\lead}[1]{\textbf{#1}\hspace{0.35em}}
\newcommand{\subhead}[1]{%
  \par\nopagebreak[3]\vspace{2.5pt}%
  {\normalsize\bfseries #1}\par\vspace{-1.5pt}}

\newcommand{\workflowsection}[4]{%
  \par\nopagebreak[4]\vspace{7.5pt}%
  {\normalsize\bfseries #1\hspace{1.1em}#2}\par\vspace{2pt}%
  \lead{Status.}#3.\par\vspace{-3.5pt}%
  \lead{The question.}#4\par}

\newcolumntype{K}{>{\bfseries\raggedright\arraybackslash}p{1.24in}}
\newcolumntype{V}{>{\raggedright\arraybackslash}p{5.53in}}
\newenvironment{skillcontract}[1]{%
  \par\nopagebreak[4]\vspace{2pt}%
  {\normalsize\bfseries Agent skill contract: \code{#1}}\par\vspace{1pt}%
  \footnotesize
  \begin{longtable}{@{}K V@{}}
  \toprule
  \textbf{Contract element} & \textbf{Requirement}\\
  \midrule
  \endfirsthead
  \toprule
  \textbf{Contract element} & \textbf{Requirement}\\
  \midrule
  \endhead
  \midrule
  \multicolumn{2}{r}{\footnotesize Continued on next page}\\
  \endfoot
  \bottomrule
  \endlastfoot
}{%
  \end{longtable}
  \small}
\newcommand{\contractrow}[2]{#1 & #2\\}

\newenvironment{smallcomparison}[1]{%
  \par\vspace{2pt}\footnotesize
  \begin{longtable}{@{}#1@{}}\toprule
}{%
  \bottomrule\end{longtable}
  \small}

Molarium is designed so that a scientist working through the graphical interface and an agent operating through the Chemist Actions API act on the same molecular state and encounter the same scientific boundaries. This appendix therefore describes Molarium by \textbf{scientific task}, not by software module.

Each workflow begins with the question a scientist is trying to answer. The human procedure explains how to perform the task and what to inspect before accepting the result. The accompanying \textbf{agent skill contract} expresses the same procedure as preconditions, permitted actions, evidence, success criteria, and stopping conditions. The two views are intentionally paired: a procedure should be understandable by a scientist and sufficiently explicit that an agent cannot silently change its meaning.

\lead{Implementation basis.}This appendix describes the inspected browser implementation prepared on 2 September 2026 for the public Molarium release. It distinguishes current browser behavior from executable but experimental behavior. Repository-only and proposed functionality is not presented as an operating instruction. User-facing workspace names are View, Design, and Simulate; the serialized values remain \code{view}, \code{build}, and \code{run} so existing action scripts remain replayable.

\workflowsection{B.1}{One molecular state, two interfaces}{Current browser surface}{What changes the molecule, and what merely changes how I am looking at it?}

Molarium separates the authoritative molecular state from the interface used to manipulate it. A scientist can work in \textbf{View}, \textbf{Design}, and \textbf{Simulate}; an agent can invoke the same guarded scientific operations through the versioned Chemist Actions API. Both routes act on the same graph, coordinates, selections, calculation records, and pending chemistry transaction. Browser-native transport---choosing a local file, writing to the clipboard, downloading a file, or entering fullscreen---remains outside the molecular API; the state-changing operation before or after that boundary has a named action.

Rotating the camera, changing a representation, focusing a component, or showing contacts changes only the presentation. Loading a structure, applying a protonation state, finishing a chemistry transaction, accepting a docking candidate, running a coordinate-producing Simulate job, or selecting or playing a saved calculation frame changes the current coordinates or graph. Refined-pose and side-chain alternatives remain candidates until explicitly applied.

\lead{The common operating rule.}Inspect, act, verify. Generate alternatives freely, but change the authoritative molecular state deliberately. A valid operation is not automatically a scientifically appropriate one; acceptance still depends on the question being asked.

\subhead{Starting from a structure}

Start with \code{?blank} when a reproducible workflow should begin from an empty canvas. Otherwise the application loads its bundled launch molecule. A structure can be pasted as XYZ or PDB text, uploaded as XYZ or PDB, entered as a SMILES string or PDB identifier in connected mode, or selected from the prepared examples. The current upload router should not be relied on for ordinary MOL-file ingestion, even though a V2000 parser exists elsewhere in the application.

After loading, inspect atom count, formula, formal charge, connected components, graph connectivity, and the 3D structure. For a selected small-molecule component, compare the synchronized 2D depiction with the 3D graph. For PDB input, remember that only the first model is loaded and alternate locations are resolved by the implemented occupancy policy. For XYZ input, bonds are inferred from proximity and should be checked rather than assumed.

\subhead{State and persistence}

Loading a new molecule resets incompatible calculation, docking, selection, and preparation state. The working session is in memory: Clear, page reload, or script import can discard unsaved work. Share copies a URL rather than serializing the current molecule, and Save produces a PNG rather than a reusable molecular record. Export the scientific record needed by the next step before replacing or clearing the scene.

\begin{skillcontract}{operate\_molecular\_state}
\contractrow{Goal}{Load and inspect a starting structure without confusing a visual change, candidate, or imported approximation with an accepted molecular state.}
\contractrow{Preconditions}{A supported input or prepared example is available; connected-mode requirements are known.}
\contractrow{Inputs}{Input type and contents; intended component; optional network policy; requested inspection fields.}
\contractrow{Procedure}{Start blank when reproducibility requires it; load through a supported route; inspect graph, components, charge, and coordinates; report ambiguities before any mutation.}
\contractrow{State mutation}{A successful load replaces the current molecule and resets dependent state. View and selection operations do not alter graph chemistry.}
\contractrow{Evidence}{Source identity or input hash where available; atom/component summary; explicit warnings about inferred bonds, alternate locations, or network retrieval.}
\contractrow{Success}{The loaded graph and coordinates correspond to the intended starting object and all material ambiguities have been surfaced.}
\contractrow{On failure}{Leave the prior scene in place where the interface does so; report the parse, fetch, or format error; do not reinterpret an unsupported file silently.}
\contractrow{Boundary}{No autosave or complete-session archive; an image or share link is not a reusable molecular state.}
\end{skillcontract}

Structure text, identifiers, and prepared examples have explicit loading actions. A local file picker only supplies bytes to \code{session.loadStructure}; it is not an alternate molecular loader:

\begin{lstlisting}
await api.execute({action:"session.loadStructure", args:{
  format:"pdb", content:pdbText, name:"starting complex"}})
// alternatives:
await api.execute({action:"session.loadIdentifier",
  args:{value:"5OVE", kind:"pdb"}})
await api.execute({action:"session.loadFixture",
  args:{fixtureId:"trp-cage"}})
await api.execute({action:"session.inspect",
  args:{scope:"all", includeCoordinates:false, maximumAtoms:500}})
await api.execute({action:"view.setDisplay",
  args:{representation:"both", showHydrogens:false,
        showPocketAtoms:true, showInteractions:true}})
\end{lstlisting}

Registered prospective routes and public structure stories retain their narrower guarded loaders (\code{designRoute.load} and \code{structureStory.load}).

\workflowsection{B.2}{Prepare a protein for local molecular design}{Experimental current surface}{I have an experimental protein structure. What must happen before I can safely edit a ligand or run a local calculation?}

A deposited PDB structure is not automatically a simulation-ready system. Atoms may be missing, alternate locations may have been chosen, protonation can be ambiguous, crystallographic waters may or may not matter, and ligands or unusual residues may exceed the supported parameterization. Molarium's preparation workflow is deliberately conservative: it makes a bounded set of assumptions explicit, performs repairs it knows how to perform, and stops when the structure exceeds those capabilities.

\subhead{Human workflow}

\begin{enumerate}
\item Load the PDB structure and inspect the protein, ligand, waters, other components, and the region around the intended binding site before preparing anything.
\item Open \textbf{Protein Preparation}. Choose pH, histidine policy, supported missing-heavy-atom repair, ligand treatment, water treatment, and gap policy.
\item Run \textbf{Prepare structure}. Molarium constructs an internal preview that inventories residues, missing atoms, sequence gaps, ligands, waters, clashes, and invalid coordinates. Within scope, it can retrieve CCD information, reconstruct supported missing side-chain atoms, add standard hydrogens and terminal oxygen, choose histidine states, and scan polar-hydrogen orientations.
\item An unresolved blocker stops the operation and exposes the preview for review. A blocker-free preview is immediately parameterized and applied; there is no normal second Apply step.
\item Inspect every applied change, warning, and attached audit. Download the preparation report if the decisions must survive the browser session.
\end{enumerate}

\subhead{What should I inspect?}

\begin{itemize}
\item Ligand identity, connectivity, completeness, and whether it should be retained or excluded.
\item Histidine and other local protonation choices near the ligand or catalytic region.
\item Crystallographic waters that bridge interactions or define the local geometry.
\item Missing atoms, chain breaks, unresolved regions, reconstructed geometry, and severe clashes near the region of interest.
\item Whether the resulting force-field and charge assumptions match the scientific question.
\end{itemize}

\lead{Acceptance is narrower than success.}A clean preview means Molarium completed its declared preparation procedure. It does not certify a production molecular-dynamics system. The implementation explicitly records \code{readyForProductionSimulation: false}.

Preparation blocks unsupported residues, unsupported missing atoms, incomplete retained ligands, disallowed gaps, invalid coordinates, and severe clashes. It does not currently construct missing loops or membranes, model general metal coordination, parameterize covalent ligands, or build periodic explicit-solvent boxes. Arbitrary proteins use an experimental Sage/Gasteiger route.

\begin{skillcontract}{prepare\_protein}
\contractrow{Goal}{Convert the loaded PDB into an explicitly reviewed molecular state suitable for a supported local Molarium calculation.}
\contractrow{Preconditions}{A PDB-derived structure is loaded; the intended downstream use is known; ligand and water roles can be reviewed.}
\contractrow{Inputs}{pH; histidine policy; missing-heavy-atom repair policy; ligand policy; water policy; gap policy.}
\contractrow{Procedure}{Inspect the starting structure; run preparation; review any blocking preview or the applied result; export the report when persistence matters; stop on unresolved blockers.}
\contractrow{State mutation}{Prepared graph and coordinates replace the prior state; preparation audit and parameterization are attached; one undo snapshot is created.}
\contractrow{Evidence}{Options, detected issues, repairs, warnings and blockers, preparation report, and resulting parameterization identity.}
\contractrow{Success}{The declared procedure completes without unresolved blockers and its scope matches the intended local calculation.}
\contractrow{Failure behavior}{Report the blocker and stop. Do not manufacture missing atoms, parameters, loop geometry, or unsupported chemistry.}
\contractrow{Escalate when}{The problem requires loops, membranes, metal coordination, covalent chemistry, periodic solvent, or another unsupported preparation capability.}
\contractrow{Does not establish}{Production simulation readiness or physical accuracy of the chosen downstream model.}
\end{skillcontract}

Representative public actions are:

\begin{lstlisting}
await api.execute({action:"protein.prepare", args:{
  pH:7.4, histidine:"auto", repairMissingHeavy:true,
  ligandPolicy:"ccd", waterPolicy:"crucial", gapPolicy:"block"}})
await api.execute({action:"session.inspect",
  args:{scope:"pocket", includeCoordinates:false, maximumAtoms:500}})
\end{lstlisting}

The preparation action already parameterizes and applies a blocker-free result. \code{protein.parameterize} is a separate alternative for an already prepared or subsequently edited complex, not the next mandatory call. The exact argument vocabulary returned by \code{api.describe()} is the runtime authority; this sequence illustrates the contract rather than replacing that manifest.

\workflowsection{B.3}{Choose the protonation state you are designing}{Current ligand workflow; protein choices within experimental preparation}{What charged and protonated species does the molecular graph actually represent?}

Protonation is part of chemical identity, not a cosmetic display choice. It changes formal charge, hydrogen count, hydrogen-bonding patterns, parameterization, and often the meaning of every calculation that follows. Molarium exposes two related but distinct workflows: ligand-state enumeration through RDKit, and protein protonation decisions inside the bounded PDB-preparation workflow.

\subhead{Ligand workflow}

\begin{enumerate}
\item Load the ligand from a SMILES string or select the ligand component whose identity you intend to change.
\item Choose a pH from 0 to 14 and enumerate candidate states. RDKit applies a versioned Dimorphite-style site table, edits formal charges and hydrogens, sanitizes each result, and returns an ordered list with estimated weights.
\item Inspect each candidate as chemistry: total charge, changed atoms, tautomeric or aromatic consequences, and compatibility with known assay conditions or the bound environment.
\item Apply one state deliberately. Molarium re-embeds the chosen species, replaces the graph and coordinates, and invalidates state that depended on the former topology.
\end{enumerate}

\subhead{Protein workflow}

For a PDB structure, set the pH and histidine policy during preparation. The implemented policies include automatic selection or explicit HID, HIE, and HIP choices, followed by inspection of polar-hydrogen orientations. Treat residues near a ligand, metal, catalytic center, or salt-bridge network as scientific decisions rather than defaults.

\lead{What the reported ligand weights mean.}They are independent Henderson-Hasselbalch estimates from the implemented site model. They are not microscopic pKa predictions, coupled-site populations, or evidence for the state selected by a protein binding site.

\begin{skillcontract}{select\_protonation\_state}
\contractrow{Goal}{Make the intended protonation and formal-charge state explicit before editing or calculation.}
\contractrow{Preconditions}{A valid ligand or PDB structure is loaded; the relevant pH or experimental condition is known or stated as an assumption.}
\contractrow{Inputs}{pH; candidate-state selection; for proteins, histidine policy and local residues requiring review.}
\contractrow{Procedure}{Enumerate or preview; compare candidates; inspect charge and local interactions; apply exactly one accepted state; re-inspect the resulting graph.}
\contractrow{State mutation}{Applying a ligand state replaces graph and coordinates. Protein protonation changes are committed through preparation. Prior parameterization becomes obsolete.}
\contractrow{Evidence}{pH, enumeration or preparation output, chosen state, total charge, changed sites, and stated rationale when more than one state is plausible.}
\contractrow{Success}{The authoritative graph explicitly represents the intended state and downstream methods will see the corresponding atoms and charges.}
\contractrow{On failure}{Report invalid pH, sanitization, embedding, or worker errors; do not partially apply or choose a different state merely because it runs.}
\contractrow{Boundary}{Enumeration weights do not resolve coupled microscopic states or protein-induced pKa shifts.}
\end{skillcontract}

Ligand-state enumeration and application use the same bounded worker and graph replacement as the visible pH controls:

\begin{lstlisting}
await api.execute({action:"ligand.enumerateProtonation",
  args:{pH:7.4, maximumStates:16}})
await api.execute({action:"ligand.applyProtonation", args:{index:0}})
await api.execute({action:"session.inspect",
  args:{scope:"ligand", includeCoordinates:false}})
\end{lstlisting}

Enumeration does not select a chemically correct state. The chosen index remains a reviewable scientific decision; protein protonation is carried separately within \code{protein.prepare}.

\workflowsection{B.4}{Edit atoms, bonds, charges, and stereochemistry in 3D}{Current browser surface}{How do I change the chemistry without corrupting the molecular graph or losing the ability to undo the operation?}

The synchronized 2D and 3D views are two ways of selecting the same graph. Chemistry-changing actions are grouped into a transaction so that a related set of edits can be validated and committed as one scientific step. Protein covalent atoms are protected from ordinary ligand editing.

\subhead{Human workflow}

\begin{enumerate}
\item Enter \textbf{Design} and select the relevant atom or bond in either view. Confirm the stable atom identity and local environment before changing it.
\item Apply the intended element replacement, bond order or aromaticity change, formal charge, explicit hydrogen, or stereochemical operation. Related operations remain in one pending chemistry transaction.
\item For coordinate-only geometry changes, select a connected path of two, three, or four atoms to set a distance, angle, or dihedral. Molarium moves the connected side defined by the graph and blocks ambiguous ring cuts.
\item Choose \textbf{Finish chemistry}. RDKit sanitizes the graph, hydrogens are reconciled, eligible local coordinates are polished, reference contacts are remapped where applicable, obsolete parameterization is invalidated, and a single undo snapshot is created.
\item If validation fails, correct the still-pending transaction or choose \textbf{Discard} to restore the pre-edit state. Use Undo/Redo only after a transaction has been committed.
\end{enumerate}

\subhead{What should I inspect?}

\begin{itemize}
\item Valence, formal charge, aromaticity, hydrogen count, stereochemistry, and the intended bond topology.
\item Whether the 2D depiction and 3D graph agree after sanitization.
\item New close contacts, distorted local geometry, and movement outside the edited region.
\item Whether topology-dependent records---parameterization, docking mapping, or calculation results---were invalidated or remapped as expected.
\end{itemize}

A synchronous mutation error restores the transaction snapshot. If the chemistry is valid but optional coordinate polishing fails, the valid graph can remain committed and the cleanup error is recorded. Fused-aromatic edits that cannot preserve the aromatic graph are rejected rather than approximated.

Distance, angle, dihedral, atom translation, and graph-chemistry controls all have named actions. A representative graph edit follows; coordinate-only edits use \code{geometry.setInternalCoordinate} or \code{geometry.translateAtoms} and remain separate from \code{chemistry.finish}.

\begin{skillcontract}{edit\_molecular\_graph}
\contractrow{Goal}{Apply an explicit, reviewable graph or local-geometry change while preserving atomic lineage and undoability.}
\contractrow{Preconditions}{The intended editable component is identified; no unrelated transaction is pending; the target atoms or bond are unambiguous.}
\contractrow{Inputs}{Selected stable atom identifiers; edit type and value; optional connected geometry path; intended transaction boundary.}
\contractrow{Procedure}{Inspect target; open or join a transaction; apply supported edits; finish once; inspect sanitized graph and local coordinates.}
\contractrow{State mutation}{Graph edits remain provisional until Finish. Finish commits one new state and invalidates obsolete parameterization; Discard restores the starting state.}
\contractrow{Evidence}{Pre-edit identity, ordered public actions, changed atoms and bonds, sanitization outcome, cleanup outcome, and final inspection.}
\contractrow{Success}{The final graph encodes the intended chemistry, sanitizes successfully, and passes local visual and structural checks.}
\contractrow{On failure}{Keep the transaction pending for correction or discard it. Never bypass sanitization or inject coordinates directly.}
\contractrow{Boundary}{Local graph validity does not establish synthetic feasibility, tautomer preference, binding affinity, or global conformational plausibility.}
\end{skillcontract}

\begin{lstlisting}
await api.execute({action:"view.setMode", args:{mode:"build"}})
await api.execute({action:"build.setTool", args:{tool:"select"}})
await api.execute({action:"selection.replace", args:{atomIds:[anchorId]}})
await api.execute({action:"chemistry.addAtom",
  args:{attachedToAtomId:anchorId, element:"Cl"}})
await api.execute({action:"chemistry.finish", args:{}})
await api.execute({action:"geometry.setInternalCoordinate",
  args:{atomIds:[a,b,c,d], value:172.0, moveConnected:true}})
await api.execute({action:"session.inspect",
  args:{scope:"ligand", includeCoordinates:true, maximumAtoms:200}})
\end{lstlisting}

The serialized value \code{build} selects the user-facing Design workspace; it is retained for action-script compatibility.

\workflowsection{B.5}{Grow an analogue by adding and positioning fragments}{Current Design workspace; pose-aware extensions experimental}{How do I add the substituent I have in mind while preserving the useful context of the parent structure?}

Fragment growth combines a graph edit with an initial 3D placement. Those are separate scientific problems: the new molecule must first be chemically correct, and its coordinates must then be a plausible starting point for relaxation or pose maintenance. Molarium's Design workspace stages a fragment template and applies its attachment as one undoable edit. Subsequent atom, bond, charge, hydrogen, or stereochemical changes use the pending chemistry transaction described above.

\subhead{Human workflow}

\begin{enumerate}
\item If the parent is bound and its pose matters, capture the reference pose \textbf{before} beginning the topology change. Retain any explicit core or interaction choices needed for later pose maintenance.
\item Select the ligand attachment atom and stage the desired fragment or template. Confirm which atoms are new, which parent atom remains the anchor, and which bond is to be formed.
\item Apply the attachment as one undoable fragment edit. Perform any associated leaving-group, bond-order, formal-charge, hydrogen, or stereochemical changes as a subsequent pending chemistry transaction.
\item Use distance, angle, and dihedral controls to give the new group a sensible initial orientation. These controls establish a starting geometry; they do not score the binding site.
\item Finish any subsequent pending chemistry, then inspect the attachment, local valence, stereochemistry, ring geometry, clashes, and preserved parent coordinates. Relax the structure or invoke reference-guided pose maintenance as a separate step.
\end{enumerate}

\subhead{How should I judge the result?}

A successful fragment operation means that the intended graph was built and sanitized. It does not mean that the chosen rotamer, ring conformation, or pocket orientation is preferred. For a solvent-exposed edit, local minimization may be enough to remove construction strain. For an edit that changes pocket contacts or releases a ring, use a workflow that samples alternatives and retains the relationship to the experimental reference.

\lead{Enumeration boundary.}Development modules can represent bounded edit plans, but general combinatorial enumeration, library management, automated campaign selection, and automatic MCS transfer of an independently imported analogue are not current workbench workflows.

\begin{skillcontract}{grow\_fragment}
\contractrow{Goal}{Construct one intended analogue and a reviewable starting geometry while retaining parent atom lineage and, when requested, reference-pose context.}
\contractrow{Preconditions}{Editable ligand and attachment atom identified; fragment chemistry specified; reference captured first if pose continuity matters.}
\contractrow{Inputs}{Parent and attachment atom identifiers; fragment/template; new bond; associated deletions, charge or stereochemical edits; optional geometry targets.}
\contractrow{Procedure}{Capture reference if needed; stage and attach the fragment; perform and finish any associated atom/bond transaction; set only the geometry needed for a starting pose; inspect; relax or sample separately.}
\contractrow{State mutation}{Attachment commits the fragment graph and an undo snapshot immediately. Finishing a later atom/bond transaction commits those additional changes. A later pose or optimization changes coordinates separately.}
\contractrow{Evidence}{Parent and product inspection, atom mapping, edit actions, transaction outcome, initial placement rationale, and any preserved reference/contact record.}
\contractrow{Success}{The attached analogue graph is correct on inspection; any subsequent chemistry transaction is sanitized; the initial geometry has no obvious construction pathology and is suitable for the chosen next workflow.}
\contractrow{On failure}{A failed attachment leaves the prior molecule in place. Correct or discard any subsequent pending chemistry transaction. Do not treat failed cleanup as permission to accept an invalid graph or arbitrary pose.}
\contractrow{Boundary}{Fragment growth is not combinatorial design, synthesis planning, global docking, or proof that the new group is energetically favored.}
\end{skillcontract}

Generic fragment-template staging and attachment use \code{fragment.stage} and \code{fragment.attach}; individual atom/bond edits and registered \code{designRoute.applyStep} operations remain available. For a registered, hash-pinned prospective route, an agent can stage a predefined product graph. \code{attachmentAtomId} is required only when that route declares designer-directed \code{spatialIntent}; the SOS1 route does not:

\begin{lstlisting}
await api.execute({action:"fragment.stage",
  args:{smiles:"c1ccncc1", name:"pyridyl", attachmentIndex:0}})
await api.execute({action:"fragment.attach",
  args:{attachedToAtomId:anchorId}})
\end{lstlisting}

\begin{lstlisting}
await api.execute({action:"designRoute.load",
  args:{routeId:"sos1-hit-only"}})
await api.execute({action:"view.setMode", args:{mode:"build"}})
await api.execute({action:"protein.prepare", args:{
  pH:7.4, histidine:"auto", repairMissingHeavy:true,
  ligandPolicy:"ccd", waterPolicy:"retain", gapPolicy:"cap"}})
await api.execute({action:"pose.captureReference",
  args:{mode:"propagate"}})
await api.execute({action:"designRoute.applyStep",
  args:{stepId:"scaffold-rewrite"}})
await api.execute({action:"designRoute.inspect", args:{}})
\end{lstlisting}

Registered routes constrain prospective inputs and permitted edits; they are not retrospective campaign ledgers.

\workflowsection{B.6}{Relax a structure, run bounded dynamics, and inspect motion}{Current and experimental methods with different scopes}{I have made a structure. Does its geometry survive relaxation, and what moves during a short calculation?}

Choose the method from the scientific question, not from a generic expectation that every engine answers the same thing. Molarium exposes several numerical pathways whose energies are not interchangeable. Record the selected model, solvent, constraints, and protocol before interpreting the result.

\begingroup
\footnotesize
\begin{longtable}{@{}p{1.31in}p{0.78in}p{1.77in}p{1.85in}@{}}
\toprule
\textbf{Question} & \textbf{Method} & \textbf{What it provides} & \textbf{Principal boundary}\\
\midrule
\endfirsthead
\toprule
\textbf{Question} & \textbf{Method} & \textbf{What it provides} & \textbf{Principal boundary}\\
\midrule
\endhead
Small-molecule cleanup or short force-field motion & RDKit & Energy, geometry optimization, and short dynamics using the RDKit force-field path. & The reported 0.1 fs dynamics path is a workbench calculation, not a production trajectory.\\
Sage energy, relaxation, or bounded single-system MD & Direct WebGPU & Browser-GPU energy, forces, geometry, and saved frames from the numeric Sage system. & Experimental; f32; direct dynamics capped at 5,000 steps; no PBC or PME.\\
Many replicas of one topology & WebGPU ensemble & Homogeneous replica dynamics with energies, frames, timing, and constraint diagnostics. & Experimental; at most 512 atoms per replica, 1,024 replicas, and 100,000 steps; nonperiodic all-pairs model.\\
Electronic energy or geometry for an eligible neutral molecule & ANI-2x & Eight-model ANI-2x ensemble energy, spread, analytical forces, and optional optimized coordinates. & Experimental; connected neutral singlets with explicit H; H/C/N/O/S/F/Cl only; at most 96 atoms.\\
Reference calculation inside another workflow & OpenMM 8.2 Reference & Double-precision internal parameterization, scoring, relaxation, and validation path. & Internal subsystem, not a selectable general Simulate method.\\
\bottomrule
\end{longtable}
\endgroup

\subhead{Human workflow}

\begin{enumerate}
\item Finish any pending chemistry. Confirm that the molecule is eligible for the chosen engine and, where required, can be completely parameterized.
\item In \textbf{Simulate}, choose the job and compatible method. Set the step count, vacuum or OBC2/ACE implicit solvent, X-H constraints where supported, saved frames, and ensemble options.
\item Run once. Inspect the engine or model identity, force field, charge model, solvent, constraints, initial and final energy, convergence or diagnostic status, elapsed time, and all warnings.
\item For dynamics, step through saved frames or replicas and inspect the region relevant to the design question. Display alignment does not overwrite raw calculation coordinates.
\item Review or export coordinates only when the result passes the declared acceptance check. Coordinate-producing jobs immediately update the current coordinates when they complete, and selecting or playing a saved frame writes that frame into the current molecule. Export the starting coordinates first when reversibility matters. Current-frame XYZ is available; saved frames are a bounded sample, and the workbench has no generic full-trajectory-file export.
\end{enumerate}

\subhead{Interpreting the model}

Automatic Sage parameterization uses deterministic Gasteiger charges rather than AM1-BCC or NAGL. Unsupported elements, invalid chemistry, or missing required parameters fail instead of silently dropping terms. The exposed workflows are nonperiodic and do not provide PME, a barostat, explicit-solvent setup, generic production protein simulation, or binding free energy.

\begin{skillcontract}{relax\_or\_simulate}
\contractrow{Goal}{Use a declared numerical model to test local geometry or bounded motion and retain enough metadata to interpret the result.}
\contractrow{Preconditions}{No chemistry transaction is pending; method eligibility and complete parameterization are established; an acceptance observable is specified.}
\contractrow{Inputs}{Job; method; steps; solvent; constraints; saved frames; fixed or movable atoms where supported; ensemble settings.}
\contractrow{Procedure}{Validate eligibility; record model and protocol; export the starting coordinates when reversibility matters; run one compatible job; inspect energies, diagnostics, and frames.}
\contractrow{State mutation}{Energy jobs do not move the molecule. Successful coordinate-producing jobs immediately replace current coordinates, and selecting a saved frame also writes current coordinates; raw frames remain calculation records.}
\contractrow{Evidence}{Engine/model version and hashes where recorded, force field, charge model, solvent, constraints, timestep, energies, convergence/diagnostics, timing, frames, and warnings.}
\contractrow{Success}{The worker returns a finite, internally valid result and the predefined structural check is satisfied within the method's declared scope.}
\contractrow{On failure}{Keep the prior graph and coordinates; record the returned error; correct eligibility, options, or assets before retrying.}
\contractrow{Boundary}{A minimized pose or short stable trace is not an affinity estimate, equilibrium ensemble, free-energy calculation, or production trajectory.}
\end{skillcontract}

The current public API exposes a bounded Design optimizer:

\begin{lstlisting}
await api.execute({action:"view.setMode", args:{mode:"build"}})
await api.execute({action:"optimization.run",
  args:{method:"induced-fit-webgpu"}})
await api.execute({action:"session.inspect",
  args:{scope:"pocket", includeCoordinates:true, maximumAtoms:500}})
\end{lstlisting}

General Simulate jobs use \code{calculation.run}; the action accepts only the installed engine/job combinations and the same bounded options as the interface. Selecting a returned frame is explicit because it writes those coordinates into the current molecule:

\begin{lstlisting}
await api.execute({action:"calculation.run", args:{
  job:"dynamics", method:"stormm",
  options:{stormmSystem:"current", steps:1000,
           savedFrameCount:20, replicaCount:32}}})
await api.execute({action:"calculation.selectFrame", args:{index:19}})
\end{lstlisting}

\workflowsection{B.7}{Explore conformations with Conformer Arena}{Experimental current surface}{What other conformations can this molecule adopt, and how sensitive is the answer to the generation and refinement method?}

Conformer search is an alternative-generation problem. Molarium separates seed generation, physical refinement, clustering, comparison, and selection so that a low-energy-looking structure is not confused with proof that the ensemble is complete or Boltzmann weighted.

\subhead{Human workflow}

\begin{enumerate}
\item Start from a chemically finished small molecule with explicit hydrogens and a sensible protonation state. Recenter the structure before tight STORMM positional comparisons, because fixed-point spatial precision degrades far from the working origin.
\item Choose \textbf{WebGPU ensemble} and \textbf{Conformer search}. Set the requested conformer count, effort, clustering cutoff, solvent/constraint options, and whether to open the comparison arena.
\item Molarium generates deterministic RDKit conformer seeds, up to the implemented cap of 256. The STORMM protocol then performs staged relaxation, heating at 600 K, cooling through 450 K and 300 K, and final relaxation across homogeneous replicas.
\item Inspect the returned ensemble before selecting a representative: energy range, distinct RMSD clusters, cluster populations, recurring torsions, ring states, constraint diagnostics, and obvious duplicates or failures.
\item Use Conformer Arena to compare supported pathways on the same molecule and settings. Ask whether the methods discover the same basins, not merely whether they rank one shared conformer similarly.
\item Select a conformer only after inspection. Selection changes the displayed/current coordinates. Export the ensemble as SDF when it fits V2000 record limits; retain method and clustering metadata separately.
\end{enumerate}

\subhead{What does the result establish?}

The workflow shows which basins the declared finite search found and how the supported methods behaved on that search. The anneal/minimize protocol is a search heuristic, so replica counts and cluster frequencies are not thermodynamic populations. Failure to find a conformation is not evidence that it is inaccessible, and energy differences from different force fields are not directly comparable.

\begin{skillcontract}{explore\_conformers}
\contractrow{Goal}{Generate, refine, cluster, and compare a finite set of conformational alternatives for one molecular topology.}
\contractrow{Preconditions}{Chemistry and protonation are finalized; molecule is eligible; intended use and comparison metric are stated.}
\contractrow{Inputs}{Conformer count; effort; cluster cutoff; seed; solvent and constraints; arena choice; acceptance or diversity criteria.}
\contractrow{Procedure}{Generate RDKit seeds; run declared refinement protocol; cluster and inspect; compare basins in Arena when requested; select or export without rewriting raw records.}
\contractrow{State mutation}{Search produces an ensemble and raw calculation evidence. Selecting a frame or conformer writes its coordinates into the current molecule.}
\contractrow{Evidence}{Seed and method identity, requested and returned counts, energies, cluster assignments, diagnostics, timings, and exported ensemble where applicable.}
\contractrow{Success}{A finite set of valid, distinct candidates is produced and its diversity is sufficient for the stated downstream purpose.}
\contractrow{On failure}{Report parameterization, eligibility, replica, overflow, or pair-work limits; reduce scope or escalate rather than hiding missing replicas.}
\contractrow{Boundary}{Not exhaustive conformational sampling, a Boltzmann ensemble, solution free energies, or proof of bioactive-conformation probability.}
\end{skillcontract}

Conformer search and Arena use \code{calculation.run}; result selection uses \code{calculation.selectConformer}. Arena axes, filtering, and sorting are presentation controls exposed through \code{calculation.setConformerView}; they do not create new conformers:

\begin{lstlisting}
await api.execute({action:"calculation.run", args:{
  job:"conformers", method:"stormm",
  options:{conformerCount:32, conformerEffort:"standard",
           conformerClusterRms:0.75, conformerArena:true}}})
await api.execute({action:"calculation.selectConformer", args:{rank:0}})
\end{lstlisting}

\workflowsection{B.8}{Predict a short single-chain protein structure}{Experimental current surface; connected mode only}{Can Molarium create a protein structure from sequence, and is this the co-folding workflow?}

\lead{This is not currently co-folding.}The implemented workflow predicts one protein entity from one sequence. It does not accept a ligand for joint protein-ligand prediction, does not predict a multimer, and does not maintain a starting ligand pose. Calling it co-folding would overstate the present capability.

The Fold Protein panel combines a configured remote MSA service with local OpenFold ONNX inference. The sequence is transmitted to the selected MSA endpoint; model inference then runs in the browser, preferring WebGPU and falling back to WASM. Local Lab disables the external MSA path.

\subhead{Human workflow}

\begin{enumerate}
\item Enter one amino-acid sequence using the standard residue alphabet or X. The current maximum length is 128 residues, using 64- or 128-residue feature buckets.
\item Inspect and accept the configured MSA endpoint and network boundary. First use may require approximately 540 MB of model assets.
\item Run and monitor MSA and inference progress. The current browser model executes three recycle passes; this count is not user-configurable. Cancel aborts the active MSA/prediction controller.
\item On success, Molarium loads the prediction as the current molecule. Inspect chain completeness, residue mapping, gross geometry, and suitability for the downstream question before replacing other work.
\item Export the predicted structure as PDB if it must persist. Preparation and local relaxation are separate decisions.
\end{enumerate}

The current path does not use templates or perform Amber relaxation. It creates no server-side job ledger. Fetch, archive, feature, model, or inference failure produces no partial prediction. A returned structure is a model output, not experimental evidence or a prepared simulation system.

\begin{skillcontract}{predict\_protein\_structure}
\contractrow{Goal}{Predict coordinates for one short protein sequence using the declared MSA service and local OpenFold model.}
\contractrow{Preconditions}{Connected mode; configured MSA endpoint; sequence length at most 128; network and model-asset use accepted.}
\contractrow{Inputs}{One sequence and MSA endpoint; the current model uses three recycle passes.}
\contractrow{Procedure}{Validate sequence; disclose external transmission; obtain and validate MSA; run the matching ONNX bucket; inspect result; export before replacing it.}
\contractrow{State mutation}{A successful prediction replaces the current molecule. Failure does not load a partial structure.}
\contractrow{Evidence}{Input sequence and endpoint, model/bucket and executed recycle count, progress/failure record, prediction metadata, and exported PDB where required.}
\contractrow{Success}{A complete finite prediction is returned for the declared input and passes basic structural inspection for the intended next step.}
\contractrow{On failure}{Report MSA, archive, rate-limit, maintenance, timeout, feature, asset, or inference failure; do not fabricate a partial structure.}
\contractrow{Boundary}{No ligand co-folding, multimer prediction, templates, Amber relaxation, experimental validation, or production preparation.}
\end{skillcontract}

The workflow is exposed as \code{protein.predict}; \code{protein.cancelPrediction} aborts the active request. The MSA endpoint remains an explicit network boundary rather than a hidden service call:

\begin{lstlisting}
await api.execute({action:"protein.predict", args:{
  sequence:fastaText, msaEndpoint:"https://msa.example/api"}})
// while active:
await api.execute({action:"protein.cancelPrediction", args:{}})
\end{lstlisting}

\workflowsection{B.9}{Maintain an experimental pose while changing ligand chemistry}{Experimental current surface}{I trust the experimental pose of my starting ligand. How can I search for a plausible pose of an edited analogue without discarding what I already know?}

Reference-guided docking is a local analogue pose-maintenance workflow. It preserves stable atom identities and selected interactions from a prepared reference complex, searches the degrees of freedom introduced by an edit, rejects infeasible candidates, and ranks distinct feasible poses in a rigid local receptor site. It is intentionally narrower than global docking.

\subhead{Human workflow}

\begin{enumerate}
\item Prepare and inspect the reference protein-ligand complex. In \textbf{Design}, capture the reference before editing. The capture stores ligand lineage, receptor atoms within 8~\AA, reference contacts, selected hydrogen bonds, settings, and provenance.
\item Use the default \textbf{Pose propagation} mode to map surviving reference atoms and seed the edited region. Use \textbf{Expert selected core} only when you can justify at least three connected, non-collinear heavy atoms as a fixed reference.
\item Edit and finish the analogue. Review remapped or unavailable contacts, transformed ring regions, atom mapping, and any constraints before starting the search.
\item Run the constrained search. Molarium verifies receptor and contact integrity, releases eligible transformed rings, generates deterministic feature and torsion seeds, freshly parameterizes the edited ligand, relaxes feasible candidates, and scores the distinct results.
\item Inspect the ranked candidates, not only rank 1. Compare mapped-core displacement, preserved or replaced interactions, ligand strain, clashes, ring and side-chain choices, and the reasons candidates were rejected.
\item Apply one candidate explicitly. Candidate generation does not change the authoritative molecule. Application rechecks topology and receptor integrity before replacing coordinates. Export the labbook as JSON and human-readable Markdown/notes.
\end{enumerate}

\subhead{What is being scored?}

The ranker uses a rigid receptor site within 8~\AA, Lennard-Jones and Coulomb ligand-receptor cross terms with relative dielectric 4, and relative vacuum ligand strain. It omits general receptor relaxation, water displacement, entropy, binding free energy, and broad macrocycle or fused-ring concerted search. The separate induced-fit Design optimizer is not part of this ranker and is not an affinity calculation.

\lead{Negative results are results.}An invalid core, altered receptor, unmappable required atom, unavailable contact feature, parameterization failure, or absence of feasible candidates produces an explicit failed or negative record. Molarium does not silently apply the least-bad pose.

\begin{skillcontract}{maintain\_reference\_pose}
\contractrow{Goal}{Generate and review plausible local poses for one edited analogue while retaining explicit relationships to a trusted reference complex.}
\contractrow{Preconditions}{Prepared complex; trusted reference pose; reference captured before edit; chemistry finished; receptor unchanged; required contacts available.}
\contractrow{Inputs}{Pose-propagation or expert-core mode; selected core or contacts; cleanup/contact policy; candidate count; analogue graph.}
\contractrow{Procedure}{Capture; edit; finish; verify mapping and receptor integrity; generate and score feasible candidates; inspect alternatives; apply one explicitly or retain a negative result.}
\contractrow{State mutation}{Capture and search create reference and candidate records. Only Apply replaces the current coordinates.}
\contractrow{Evidence}{Protocol/settings snapshot, hashes, atom mapping, contact decisions, constraints, rejected candidates, scores, selected pose, and chained labbook.}
\contractrow{Success}{At least one feasible distinct candidate survives the declared constraints and a human or agent can justify any applied selection from the recorded evidence.}
\contractrow{On failure}{Apply nothing; preserve the negative labbook; repair the reference, mapping, constraints, or chemistry before rerunning.}
\contractrow{Boundary}{Local analogue pose maintenance only---not global docking, automatic MCS pose transfer, receptor-ensemble docking, or affinity/free-energy prediction.}
\end{skillcontract}

\begin{lstlisting}
await api.execute({action:"pose.captureReference",
  args:{mode:"propagate"}})
// perform and finish the intended molecular-graph edit
await api.execute({action:"pose.refine",
  args:{searchChains:16}})
await api.execute({action:"pose.apply", args:{index:0}})
await api.execute({action:"session.inspect",
  args:{scope:"pocket", includeCoordinates:true, maximumAtoms:500}})
\end{lstlisting}

For a receptor-side-chain alternative, the required order is:

\begin{lstlisting}
await api.execute({action:"pose.enumerateSidechainRotamers",
  args:{receptorAtomId:sidechainAtomId, maximumCandidates:32}})
await api.execute({action:"pose.applySidechainRotamer", args:{index:5}})
await api.execute({action:"pose.updateReceptorReference", args:{}})
await api.execute({action:"pose.refine", args:{searchChains:64}})
\end{lstlisting}

Enumeration itself remains non-mutating; the selected receptor coordinates are applied explicitly and then made the new reference before another ligand-pose search.

\workflowsection{B.10}{Record, replay, and hand off a design process}{Current action/replay and Design History surfaces}{I have done something scientifically useful. What must I export so that another scientist or an agent can understand and reproduce it?}

Molarium records several different kinds of evidence. They answer different questions and should not be treated as interchangeable. The live action audit says what was requested; a calculation or workflow record says what method ran and what it returned; an action script says how to replay completed public operations; and Design History represents a durable sequence of content-addressed states and decisions. The current workbench does not automatically combine every scientific record into one complete project archive.

\begingroup
\footnotesize
\begin{longtable}{@{}p{1.23in}p{2.76in}p{2.61in}@{}}
\toprule
\textbf{Record} & \textbf{Question answered} & \textbf{Persistence and limit}\\
\midrule
\endfirsthead
\toprule
\textbf{Record} & \textbf{Question answered} & \textbf{Persistence and limit}\\
\midrule
\endhead
Chemist Actions audit & What recognized public action was requested, with what arguments and outcome? & Session memory; serialized calls; bounded history, normally 500 entries; cleared with the scene unless exported through a script path.\\
Preparation or calculation record & What system, model, settings, energies, frames, warnings, or audit produced this result? & Attached to the in-memory state and available through workflow-specific reports or exports.\\
Docking labbook & How were mapping, feasibility, scoring, rejection, and selection handled? & Exportable JSON and Markdown/notes with chained hashes.\\
Designer action script & Which completed public actions should be replayed? & Exportable JSON; replay invokes the same guarded routes; not a complete UI or binary session snapshot.\\
Registered design route & What prospective, hash-pinned edit/evaluation sequence is permitted? & Validated route schema; distinct from retrospective history.\\
Design History campaign & What content-addressed structures, decisions, commits, and branches form the design history? & Live browser surface with locally persisted commits, branches, merges, decisions, integrity events, import, and export.\\
\bottomrule
\end{longtable}
\endgroup

\subhead{Replay workflow}

Designer Moves import begins from a blank canvas and performs script-level validation of schema, recognized action names, binding order, and forbidden shortcut keys. Individual action arguments and molecular preconditions are checked when each step executes. Conversion from an audit includes only completed public actions and excludes campaign administration. A standalone script therefore begins with an explicit reproducible loader such as \code{session.loadStructure}, \code{session.loadIdentifier}, \code{session.loadFixture}, or \code{designRoute.load}; it cannot depend on an unrecorded starting scene.

\begin{enumerate}
\item Export any unsaved molecular or workflow records before importing a script; import intentionally clears the existing scene.
\item Create a Designer action script from completed recognized public actions. Failed actions remain audit evidence but are not silently converted into successful replay steps.
\item Import the validated script, restart it, and Play or step through it. Replay uses the same Chemist Actions routes, runtime validations, and workers as manual operation and produces a separate replay audit.
\item Compare the resulting graph, coordinates, decisions, and errors with the intended endpoint. Export the replay audit and scientific records needed by the recipient.
\item For a registered story or route, verify schema, referenced content, and SHA-256 before execution. Treat a hash chain as integrity evidence, not an authenticated author identity or proof of scientific truth.
\end{enumerate}

Replay provides play/pause, previous/next-step review, stable captions, and a separate audit. Failed audit entries remain evidence but are not converted silently into successful replay steps.

\subhead{Live Design History}

Design History is a live browser feature. \code{campaign.create} creates a locally persisted campaign and atomically commits the visible molecule. \code{campaign.commitCurrent} stores a complete graph-and-coordinate snapshot together with completed public actions since the prior commit. Branch switching restores the exact head molecule and refuses to proceed when the visible molecular state is dirty. Merge records the currently visible result with two parents; it is not an automatic structural merge. Decisions attach disposition, rationale, reason codes, and optional evidence to a commit. Verification recomputes content hashes, the append-only event chain, commit relationships, and branch heads derived from that chain.

Campaign JSON and the active branch are stored in browser IndexedDB. Reload restores the graph and coordinates at the active branch head; it does not restore parameter tables, docking candidates, calculation frames, view state, or the complete workbench session. \code{campaign.import} verifies the canonical JSON and requires a restorable graph/coordinate snapshot before mutating live state; \code{campaign.export} returns that canonical JSON. Closing a campaign clears the active pointer without deleting its commits. A live commit can attach a script of completed molecular Chemist Actions; the complete molecular snapshot remains authoritative. Campaign-management actions are excluded from molecule replay scripts so replay cannot recursively administer its own history.

\subhead{Choose the export that matches the handoff}

\begingroup
\footnotesize
\begin{longtable}{@{}p{1.20in}p{2.48in}p{2.92in}@{}}
\toprule
\textbf{Export} & \textbf{Carries} & \textbf{Does not carry}\\
\midrule
\endfirsthead
\toprule
\textbf{Export} & \textbf{Carries} & \textbf{Does not carry}\\
\midrule
\endhead
XYZ current frame & Displayed coordinates & Full topology metadata, history, or provenance\\
SDF conformer ensemble & Exportable conformer graph and coordinates & Unrelated session state or complete search history\\
PDB protein & Imported or predicted protein coordinates and supported annotations & Complete application or preparation provenance\\
Preparation JSON & Policies, issues, audit, and readiness status & A production-readiness guarantee\\
Docking JSON/Markdown & Settings, mapping, candidates, failures, hashes, and selection & Author authentication or binding free energy\\
Action script JSON & Public actions and bindings & Arbitrary code, direct coordinates, or complete view state\\
PNG or share URL & A visual snapshot or current/registered URL & A reusable molecule or serialized live session\\
\bottomrule
\end{longtable}
\endgroup

\begin{skillcontract}{export\_and\_handoff}
\contractrow{Goal}{Preserve the smallest complete set of molecular, procedural, numerical, and decision records needed by the next scientist or agent.}
\contractrow{Preconditions}{The authoritative result and its intended downstream use are identified; unsaved state has not yet been cleared or replaced.}
\contractrow{Inputs}{Recipient and purpose; required molecular format; action/replay needs; calculation, preparation, docking, or campaign evidence to retain.}
\contractrow{Procedure}{Inventory records; export reusable molecular state and method evidence; create an action script when replay is required; commit decision history when project continuity is required; verify hashes and replay; describe missing pieces explicitly.}
\contractrow{State mutation}{Export is non-mutating. Script import is destructive to the unsaved scene and must follow export of work that must be retained.}
\contractrow{Evidence}{Exported files, schema and version identifiers, hashes where recorded, replay audit, and a statement of assumptions and unresolved failures.}
\contractrow{Success}{The recipient can identify the starting state, reproduce guarded actions where supported, inspect the method and result, and distinguish accepted decisions from candidates.}
\contractrow{On failure}{Stop on schema, action, binding, hash, or referenced-content mismatch. Do not describe a partial replay as equivalent to a completed audit.}
\contractrow{Boundary}{No complete project bundle or automatic capture of browser-native transport, calculation binaries, or every transient view gesture. An action script is not a substitute for the authoritative snapshot.}
\end{skillcontract}

\begin{lstlisting}
await api.execute({action:"campaign.create", args:{
  campaignId:"sos1-local-design", title:"SOS1 local design",
  initialCommitMessage:"Register 5OVE/AXE starting state"}})
// chemist or agent performs guarded molecular actions
await api.execute({action:"campaign.commitCurrent", args:{
  message:"Accept Phe890-out design state", tags:["prospective"]}})
await api.execute({action:"campaign.recordDecision", args:{
  disposition:"progressed",
  rationale:"No new clashes after pocket-state search"}})
await api.execute({action:"campaign.verify", args:{}})
\end{lstlisting}

\workflowsection{B.11}{How skill contracts, API actions, and Design History fit together}{Current conceptual and executable relationship}{How does a written agent contract become an authorized action and then a durable scientific decision?}

The three layers serve different purposes and should not be collapsed:

\begin{enumerate}
\item \textbf{The skill contract defines scientific intent.} It states preconditions, permissible operations, required evidence, success, failure, and escalation. Here these names are proposed narrative operating contracts stored with the manuscript. They are not installed runtime skills, Chemist Actions, schemas, or executable enforcement.
\item \textbf{The Chemist Actions API defines executable authority.} \code{api.describe()} lists installed action names and bounded arguments. Calls are validated, serialized through one queue, and audited with sequence, status, result or error, and duration. An agent cannot gain extra molecular authority by phrasing a request differently.
\item \textbf{Design History defines durable meaning over time.} It commits complete molecular states and can link them to an explicitly partial public-action script, branches, decisions, evidence, actors, and append-only integrity records. The snapshot answers ``what state was accepted?''; the action script answers ``which guarded operations were captured and can be replayed?''; and the decision answers ``why was it accepted?''
\end{enumerate}

The practical agent pattern is:

\begin{lstlisting}
const manifest = api.describe();          // discover installed authority
await api.execute({action:"session.inspect", args:{scope:"all"}})
await api.execute({action:"geometry.setInternalCoordinate", args:{
  atomIds:[a,b,c,d], value:172.0, moveConnected:true}})
await api.execute({action:"session.inspect", args:{scope:"all"}})
await api.execute({action:"campaign.commitCurrent",
  args:{message:"Accepted state after declared check"}})
await api.execute({action:"campaign.recordDecision",
  args:{disposition:"progressed", rationale:"..."}})
await api.execute({action:"campaign.verify", args:{}})
\end{lstlisting}

The manifest is the executable authority: a narrative skill contract cannot add an action, relax argument validation, or bypass a molecular precondition. Scientific review remains necessary even when an operation is exposed to both interface and agent.

\lead{Handoff rule.}Export the structure when the recipient needs coordinates; export the calculation record when the recipient must judge the method; export the Designer Moves script together with its required starting state when the recipient must replay operations; and commit to Design History when the recipient must recover branches, decisions, and accepted molecular states. Hashes detect alteration of the serialized record. They do not establish authorship, chemical equivalence, or scientific validity.

\lead{The appendix in one sentence.}For every task: make the molecular state explicit, choose a method whose scope matches the question, inspect before applying a candidate, retain the evidence needed to reproduce the decision, and stop when the declared scientific boundary is crossed.

\subhead{Capabilities deliberately outside these workflows}

The current workbench does not provide missing-loop construction, membrane or general metal preparation, covalent-ligand parameterization, periodic explicit solvent, PME, production-scale trajectories, binding free energy, global docking, general remote CUDA execution, GFN/xTB, general combinatorial enumeration, or automatic campaign selection. Their absence is part of the operating contract rather than a prompt to approximate them silently.

{\footnotesize
This task-centered appendix intentionally omits maintainer file paths, source-line numbers, worker lifecycle details, JSON nesting limits, deployment internals, and the complete repository test list. Those remain valuable implementation and audit documentation, but they are not required for a scientist or agent to understand how to perform and judge the workflows above.\par}
\endgroup


\appendixsection{Editorial history: recovered paragraph versions}{app:paragraph-history}
\begingroup
\small
\setlength{\emergencystretch}{2em}
\lstset{language={},basicstyle=\ttfamily\footnotesize,breaklines=true,breakatwhitespace=false,columns=fullflexible,keepspaces=true,showstringspaces=false,aboveskip=0.4em,belowskip=0.7em,literate={’}{{'}}1}

\textit{Editorial working appendix, prepared 7 September 2026. This is an evidence archive for another editing agent, not a proposed rewrite or a scientific validation record.}

\subsection{Scope, provenance, and use}

This appendix preserves the recovered alternatives for 38 changed paragraph or structured-block groups. Each group's current version is marked explicitly. The other 190 baseline records have only one distinct recovered version and are indexed at the end rather than duplicating text already present in the manuscript. Historical LaTeX is shown verbatim so that old citations, figures, labels, and formatting cannot modify the active manuscript. The appendix is self-contained within main.tex; it does not require a separate Markdown file.

The available history is incomplete. The local Git manuscript history inspected for this reconstruction starts with an import on 6 September 2026. The 4--5 September commits inspected in an earlier session are no longer available in this checkout. The uploaded ZIP and selected full paragraphs from earlier conversation previews recover additional wording, but not every intermediate revision. Obtain the earlier project revision export to complete the history.

Git identifies the snapshot author as Crixet Sandbox, not an individual editor. The user's report that Woody revised the manuscript is context, not evidence attributing a particular paragraph to him. Sources A and B are archived drafts, not a proven chronological chain. ZIP member timestamps have no timezone and are not reliable edit times; conversation timestamps record when wording was observed. Version numbers below indicate presentation order, not established editing chronology.

Wording, typos, and LaTeX commands are retained. Whitespace-only differences are deduplicated. Nonidentical paragraph correspondences are explicitly marked as inferred or reviewed by topic; splits and merges are not automatically reconstructed. Source E preserves an earlier Figure 1 caption without its figure wrapper. Source A refers to an external Appendix B file absent from the ZIP; its missing contents cannot be recovered from that source.

Do not restore historical text indiscriminately. Preserve the requested Apache License 2.0 wording unless explicitly directed otherwise. Old broken references and misspellings are evidence, not recommendations. One substantive conflict needs review: archived Appendix A describes live campaign actions and IndexedDB persistence, whereas the current version describes repository-level ledger functions with live promotion not exposed. Both claims are retained without deciding which is correct.

Baseline source line numbers refer to main.tex immediately before this appendix was inserted; lines in the main text and Appendices A and B remain unchanged. The baseline SHA-256 is:

\begin{lstlisting}
61e57f2d6944b25becb1c033e97c098d1848fa35d90027485f838c8baca6cc8f
\end{lstlisting}

\subsection{Source inventory}

\paragraph{Source A}

ZIP member timestamp 2026-09-02 21:59:02 (timezone unspecified; not an edit time).

\begin{lstlisting}
Molarium_paper_source_2026-09-02.zip :: Molarium_paper_source_2026-09-02/main-original-input.tex
\end{lstlisting}

Extracted records: 119. Source SHA-256:

\begin{lstlisting}
5f639cc0a4bf05e0ae88725dc1957e64b09d29099b1b6252b51cbe211645c909
\end{lstlisting}

\paragraph{Source B}

ZIP member timestamp 2026-09-02 21:06:42 (timezone unspecified; not an edit time).

\begin{lstlisting}
Molarium_paper_source_2026-09-02.zip :: Molarium_paper_source_2026-09-02/Molarium_complete.tex
\end{lstlisting}

Extracted records: 228. Source SHA-256:

\begin{lstlisting}
bb47e79e91b413c1b81ceb1047de0ba7aff0ca03e4987d9e853c75654ec486f6
\end{lstlisting}

\paragraph{Source C}

Observed 2026-09-03T04:51:39.502Z; snapshot time, not edit time.

\begin{lstlisting}
Earlier conversation: six full paragraphs from the main.tex preview accompanying the last-edit question
\end{lstlisting}

Extracted records: 6. Source SHA-256:

\begin{lstlisting}
b715e1807044780fee0c6f1205aeacb258dd77043d990e0448fd1d11c866aaeb
\end{lstlisting}

\paragraph{Source D}

Observed 2026-09-03T04:52:31.548Z; snapshot time, not edit time.

\begin{lstlisting}
Earlier conversation: development paragraph from the preview accompanying 'by anyone'
\end{lstlisting}

Extracted records: 1. Source SHA-256:

\begin{lstlisting}
0a31bdeb0555b4b7cde58d2ea9cb3bf10afa23b9353fd670e02e94b635976140
\end{lstlisting}

\paragraph{Source E}

Observed during 5 September 2026 requests at 00:03:42.437Z and 00:06:45.642Z; request times, not precise edit times.

\begin{lstlisting}
Earlier tool readouts: full Figure 1 caption before Solvent correction, and full paragraph before removal of The before RDKit; known single-word edits reconstructed against the unchanged baseline
\end{lstlisting}

Extracted records: 2. Source SHA-256:

\begin{lstlisting}
255560d9089cdc228a0d8f23de04212795e1b7578124f635fc1ea1f20d0312dd
\end{lstlisting}

\paragraph{Source G1}

2026-09-06T08:33:29Z | author: Crixet Sandbox <sandbox@example.com>.

\begin{lstlisting}
146c4fc1f33c4a8da5948a9e81edcca23de438fa:main.tex
\end{lstlisting}

Extracted records: 228. Source SHA-256:

\begin{lstlisting}
61e57f2d6944b25becb1c033e97c098d1848fa35d90027485f838c8baca6cc8f
\end{lstlisting}

\paragraph{Source LIVE}

Read 2026-09-07; not an edit timestamp.

\begin{lstlisting}
main.tex before insertion of Appendix C
\end{lstlisting}

Extracted records: 228. Source SHA-256:

\begin{lstlisting}
61e57f2d6944b25becb1c033e97c098d1848fa35d90027485f838c8baca6cc8f
\end{lstlisting}

For C, D, and E, hashes cover assembled excerpts with section delimiters, not full historical manuscript files. Their line numbers are excerpt-local. Conversation excerpts are selected rather than exhaustive. LIVE is the frozen baseline before this appendix, not the updated file containing the appendix.

\subsection{Changed-record index}

\begin{longtable}{@{}p{0.10\linewidth}p{0.12\linewidth}p{0.10\linewidth}p{0.60\linewidth}@{}}
\toprule Record & Line & Versions & Location\\\midrule\endhead
R001 & 97 & 2 & Abstract \\
R002 & 99 & 2 & Abstract \\
R003 & 101 & 3 & Abstract \\
R005 & 108 & 3 & From fixed software products to generated adaptive software \\
R006 & 110 & 3 & From fixed software products to generated adaptive software \\
R007 & 112 & 3 & From fixed software products to generated adaptive software \\
R008 & 114 & 3 & From fixed software products to generated adaptive software \\
R009 & 116 & 4 & From fixed software products to generated adaptive software \\
R010 & 120 & 3 & Molarium and the value of a persistent molecular state \\
R011 & 122 & 4 & Molarium and the value of a persistent molecular state \\
R012 & 131 & 2 & Molarium and the value of a persistent molecular state \\
R013 & 133 & 2 & Molarium and the value of a persistent molecular state \\
R014 & 135 & 2 & Molarium and the value of a persistent molecular state \\
R015 & 142 & 2 & Molarium and the value of a persistent molecular state \\
R017 & 146 & 2 & Molarium and the value of a persistent molecular state \\
R018 & 148 & 2 & Molarium and the value of a persistent molecular state \\
R020 & 159 & 2 & Molarium and the value of a persistent molecular state \\
R021 & 163 & 2 & Scientific contracts: provenance, validation, and falsifiability \\
R022 & 165 & 2 & Scientific contracts: provenance, validation, and falsifiability \\
R023 & 167 & 3 & Scientific contracts: provenance, validation, and falsifiability \\
R024 & 169 & 2 & Scientific contracts: provenance, validation, and falsifiability \\
R027 & 180 & 2 & Scientific contracts: provenance, validation, and falsifiability \\
R028 & 182 & 2 & Scientific contracts: provenance, validation, and falsifiability \\
R031 & 195 & 2 & Local computation changes the interaction model \\
R032 & 197 & 2 & Local computation changes the interaction model \\
R034 & 201 & 2 & Local computation changes the interaction model \\
R035 & 203 & 2 & Local computation changes the interaction model \\
R060 & 266 & 2 & Conclusion \\
R062 & 272 & 2 & AI assistance and availability \\
R096 & 421 & 2 & Molarium architecture, numerical validation, and executable design history \\
R099 & 432 & 2 & Molarium architecture, numerical validation, and executable design history \\
R101 & 449 & 2 & Molarium architecture, numerical validation, and executable design history \\
R102 & 451 & 2 & Molarium architecture, numerical validation, and executable design history \\
R105 & 477 & 2 & Molarium architecture, numerical validation, and executable design history \\
R108 & 485 & 2 & Molarium architecture, numerical validation, and executable design history \\
R111 & 505 & 2 & Molarium architecture, numerical validation, and executable design history \\
R113 & 509 & 2 & Molarium architecture, numerical validation, and executable design history \\
R118 & 550 & 2 & Molarium architecture, numerical validation, and executable design history \\
\bottomrule\end{longtable}

\subsection{R001: Abstract}

Baseline main.tex line 97; paragraph.

\paragraph{Version 1 --- earlier/alternative wording}

\begin{itemize}[leftmargin=*]
\item Source A, line 97: inferred correspondence, similarity 0.94; verify before merging.
\item Source B, line 97: inferred correspondence, similarity 0.94; verify before merging.
\end{itemize}

\begin{lstlisting}
Commercial computational drug-discovery software has traditionally bundled three sources of value: the interface, the scientific models, and the specialized expertise required to operate those models and translate their outputs into design decisions. Large language models, coding agents, and tool-using scientific agents are beginning to unbundle this paradigm. Interfaces, analyses, workflow logic, and substantial portions of scientific software can increasingly be generated to meet the needs of a particular scientist, project, and question, while agents can automate sophisticated workflows that once required specialist users.
\end{lstlisting}

\paragraph{Version 2 --- CURRENT}

\begin{itemize}[leftmargin=*]
\item Source G1, line 97: exact text.
\item Source LIVE, line 97: exact text.
\end{itemize}

\begin{lstlisting}
Commercial drug discovery software has traditionally bundled three sources of value: the interface, the scientific models, and the specialized expertise required to run the models through the interface, thereby translating outputs into design decisions. Large language models, coding agents, and tool-using agents are beginning to unbundle this paradigm. Interfaces, analyses, workflow logic, and substantial portions of scientific software can increasingly be generated to meet the needs of a particular scientist, project, and question, while agents can automate sophisticated workflows that once required specialist users.
\end{lstlisting}

\subsection{R002: Abstract}

Baseline main.tex line 99; paragraph.

\paragraph{Version 1 --- earlier/alternative wording}

\begin{itemize}[leftmargin=*]
\item Source A, line 99: inferred correspondence, similarity 0.81; verify before merging.
\item Source B, line 99: inferred correspondence, similarity 0.81; verify before merging.
\end{itemize}

\begin{lstlisting}
This shift repositions where competitive advantage is created in computational drug discovery. Faster execution cannot fix an imprecise model, broaden an applicability domain, or enhance the utility of a weak prediction. As the underlying tooling becomes increasingly accessible and fluid, strategic differentiation should shift toward superior and more generalizable models, proprietary experimental data, accumulated discovery intelligence, and the human discernment needed to frame the right questions and evaluate the answers.
\end{lstlisting}

\paragraph{Version 2 --- CURRENT}

\begin{itemize}[leftmargin=*]
\item Source G1, line 99: exact text.
\item Source LIVE, line 99: exact text.
\end{itemize}

\begin{lstlisting}
This shift repositions where competitive advantage is created in computational drug discovery. More intuitive interfaces and faster execution alone do not fix model inaccuracies or broaden the applicability domain. As the underlying tooling becomes increasingly accessible and fluid, strategic differentiation will shift away from familiarity with a particular software package and toward superior preditive algorithms, more generalizable models, proprietary experimental data, accumulated discovery intelligence, and the human discernment needed to frame the right questions and evaluate the answers.
\end{lstlisting}

\subsection{R003: Abstract}

Baseline main.tex line 101; paragraph.

\paragraph{Version 1 --- earlier/alternative wording}

\begin{itemize}[leftmargin=*]
\item Source A, line 101: inferred correspondence, similarity 0.92; verify before merging.
\item Source B, line 101: inferred correspondence, similarity 0.92; verify before merging.
\end{itemize}

\begin{lstlisting}
Molarium, available at \url{https://molarium.org}, provides a case study of this transition. Created during an intensive weekend through collaboration between a drug-discovery scientist and an AI coding agent, this open-source molecular workbench operates primarily in a standard web browser. It integrates local molecular graphics, editing, conformer generation, molecular mechanics and dynamics, machine-learned potentials, and selected protein-structure prediction workflows. Beneath its rapidly adaptable interface, Molarium maintains a versioned molecular state together with user actions, method specifications, provenance, and validation evidence. Decoupling an adaptable interface from an enduring scientific core points toward a broader architecture for generated adaptive software in drug discovery.
\end{lstlisting}

\paragraph{Version 2 --- earlier/alternative wording}

\begin{itemize}[leftmargin=*]
\item Source C, line 3: inferred correspondence, similarity 0.99; verify before merging.
\end{itemize}

\begin{lstlisting}
Molarium, available at \url{https://molarium.org}, provides a case study of this transition in the context of molecular design software. Created during an intensive weekend effort of a tireless drug-discovery scientist and an AI coding agent, this open-source molecular workbench operates primarily in a standard web browser. It integrates local molecular graphics, editing, conformer generation, energy minimization, molecular dynamics, machine-learned potentials, and protein-structure prediction workflows. Beneath its rapidly adaptable interface, Molarium maintains a versioned molecular state together with user actions, method specifications, provenance, and validation evidence. Decoupling an adaptable interface from an enduring scientific core points toward a broader architecture for generated adaptive software in drug discovery.
\end{lstlisting}

\paragraph{Version 3 --- CURRENT}

\begin{itemize}[leftmargin=*]
\item Source G1, line 101: exact text.
\item Source LIVE, line 101: exact text.
\end{itemize}

\begin{lstlisting}
Molarium, available at \url{https://molarium.org}, provides a case study of this transition in the context of molecular design software. Created during an intensive weekend effort of a tireless drug-discovery scientist and an AI coding agent, this open-source molecular workbench operates exclusively in a standard web browser. It integrates local molecular graphics, editing, conformer generation, energy minimization, molecular dynamics, machine-learned potentials, and protein-structure prediction workflows. Beneath its rapidly adaptable interface, Molarium maintains a versioned molecular state together with user actions, method specifications, provenance, and validation evidence. Decoupling an adaptable interface from an enduring scientific core points toward a broader architecture for generated adaptive software in drug discovery.
\end{lstlisting}

\subsection{R005: From fixed software products to generated adaptive software}

Baseline main.tex line 108; paragraph.

\paragraph{Version 1 --- earlier/alternative wording}

\begin{itemize}[leftmargin=*]
\item Source A, line 108: inferred correspondence, similarity 0.84; verify before merging.
\item Source B, line 108: inferred correspondence, similarity 0.84; verify before merging.
\end{itemize}

\begin{lstlisting}
In computational drug discovery, success depends on deploying the right methodology at the right time, supported by appropriate parameters, explicit assumptions, and sound scientific context. Historically, this work has been concentrated among computational specialists and has required substantial operational overhead. Docking depends on protein preparation, protocol selection, constraints, and critical pose evaluation. Molecular dynamics adds system setup, force-field assignment, compute infrastructure, trajectory analysis, and visualization. Free-energy calculations add perturbation design, sampling sufficiency, convergence assessment, and alignment with experimental conditions. Comparable layers of infrastructure surround cheminformatics, quantum chemistry, property prediction, generative molecular design, and structure prediction.
\end{lstlisting}

\paragraph{Version 2 --- earlier/alternative wording}

\begin{itemize}[leftmargin=*]
\item Source C, line 6: inferred correspondence, similarity 0.97; verify before merging.
\end{itemize}

\begin{lstlisting}
In computational drug discovery, success depends on deploying the right methodology at the right time, supported by appropriate parameters, explicit assumptions, and sound scientific context. Historically, this work has been concentrated among computational specialists and has required substantial operational overhead. Docking depends on protein preparation, protocol selection, constraints, and pose evaluation. Molecular dynamics requires complex system setup, solvation, ion placement, force-field assignment, compute infrastructure, equilibration, production simulation, potential for enhanced sampling, trajectory analysis, visualization, and interpretation in the tonext of project-specific decisions. Free-energy calculations add alchemical perturbations, sampling windows, convergence assessment, error analysis, and alignment with experimental conditions. Comparable layers of infrastructure surround cheminformatics, quantum chemistry, ADME property prediction, generative molecular design, and protein structure prediction.
\end{lstlisting}

\paragraph{Version 3 --- CURRENT}

\begin{itemize}[leftmargin=*]
\item Source G1, line 108: exact text.
\item Source LIVE, line 108: exact text.
\end{itemize}

\begin{lstlisting}
In computational drug discovery, success depends on deploying the right methodology at the right time, supported by appropriate parameters, explicit assumptions, and sound scientific context. Historically, this work has been concentrated among computational specialists and has required substantial operational overhead. Docking depends on protein preparation, protocol selection, constraints, and pose evaluation. Molecular dynamics requires complex system setup, solvation, ion placement, force-field assignment, compute infrastructure, equilibration, production simulation, potential for enhanced sampling, trajectory analysis, visualization, and interpretation in the context of project-specific decisions. Free-energy calculations add further complexities to  molecular simulations with alchemical perturbations, sampling windows, convergence assessment, error analysis, and alignment with experimental conditions. Comparable layers of infrastructure and expertise surround cheminformatics, quantum chemistry, ADME property prediction, generative molecular design, and protein structure prediction.
\end{lstlisting}

\subsection{R006: From fixed software products to generated adaptive software}

Baseline main.tex line 110; paragraph.

\paragraph{Version 1 --- earlier/alternative wording}

\begin{itemize}[leftmargin=*]
\item Source A, line 110: inferred correspondence, similarity 0.75; verify before merging.
\item Source B, line 110: inferred correspondence, similarity 0.75; verify before merging.
\end{itemize}

\begin{lstlisting}
Commercial scientific software has created considerable value by simplifying and integrating many of these activities. Graphical interfaces reduced technical barriers, workflow engines connected fragmented programs, and software packages encoded operational details needed to apply methods consistently. The same integration coupled scientific questions to relatively fixed software environments. Project-specific improvements competed with broader product roadmaps, new analyses often required engineering support, and scientific users frequently adapted their questions to the capabilities and conventions of the available applications.
\end{lstlisting}

\paragraph{Version 2 --- earlier/alternative wording}

\begin{itemize}[leftmargin=*]
\item Source C, line 8: inferred correspondence, similarity 0.95; verify before merging.
\end{itemize}

\begin{lstlisting}
Historically, commercial scientific software has created considerable value by simplifying and integrating many of these activities. Graphical interfaces reduced technical barriers, workflow engines connected fragmented programs, and software packages encoded operational details needed to apply methods consistently. A single static inferface (perhaps updated quarterly) coupled scientific questions to a fixed software environments. User-specific requests competed with broader product roadmaps that attempted to balance the needs of a diverse users base. Implementing new methods from the literature or integrating open-source software often required software engineering support, if even possible. Scientific users frequently adapted their questions to the capabilities and conventions of the available tools within their preferred/available graphical interface.
\end{lstlisting}

\paragraph{Version 3 --- CURRENT}

\begin{itemize}[leftmargin=*]
\item Source G1, line 110: exact text.
\item Source LIVE, line 110: exact text.
\end{itemize}

\begin{lstlisting}
Historically, commercial drug discovery software has created considerable value by simplifying and integrating many of these activities. Graphical interfaces reduced technical barriers, workflow engines connected fragmented programs, and software packages encoded operational details needed to apply methods consistently. A single static inferface (perhaps updated quarterly) coupled scientific questions to a fixed software environments. User-specific requests competed with broader corporate product roadmaps that attempted to balance the needs of a diverse user base. Implementing new methods from the literature or integrating open-source software often required software engineering support, if it was even possible. Scientific users adapted their questions to the capabilities and conventions of the available tools within their preferred/available graphical interface rather than adapting the software to their needs.
\end{lstlisting}

\subsection{R007: From fixed software products to generated adaptive software}

Baseline main.tex line 112; paragraph.

\paragraph{Version 1 --- earlier/alternative wording}

\begin{itemize}[leftmargin=*]
\item Source A, line 112: inferred correspondence, similarity 0.91; verify before merging.
\item Source B, line 112: inferred correspondence, similarity 0.91; verify before merging.
\end{itemize}

\begin{lstlisting}
Large language models and coding agents are beginning to change this relationship. A scientist with deep drug-discovery expertise but limited programming experience can increasingly describe an application, analysis, visualization, or workflow in natural language and have software constructed around the scientific need. Tool-using agents can operate established methods, link calculations, handle routine failures, and assemble results into forms that support decisions. Earlier systems such as ChemCrow and ChemGraph established proof-of-concept for language models coordinating chemistry tools and computational workflows \citep{bran2024chemcrow,boiko2023coscientist,pham2026chemgraph,macknight2025rethinking}. Their practical scope was limited compared with current coding agents, but they helped establish the feasibility of coupling language models to executable scientific tools.
\end{lstlisting}

\paragraph{Version 2 --- earlier/alternative wording}

\begin{itemize}[leftmargin=*]
\item Source C, line 10: inferred correspondence, similarity 0.99; verify before merging.
\end{itemize}

\begin{lstlisting}
Large language models and coding agents are beginning to change this relationship. A scientist with deep drug-discovery expertise but limited software programming experience can increasingly describe an application, analysis, visualization, or workflow in natural language and have software constructed around the scientific need. Tool-using agents can operate established methods, build complex workflows, implemnt methods from the literature, handle failures, assemble results into forms that support decisions, and even optimize processed based on iterative feedback. Earlier systems such as ChemCrow and ChemGraph established proof-of-concept for language models coordinating chemistry tools and computational workflows \citep{bran2024chemcrow,boiko2023coscientist,pham2026chemgraph,macknight2025rethinking}. While their practical scope was limited compared with current coding agents, they helped establish the feasibility of coupling language models to executable scientific tools.
\end{lstlisting}

\paragraph{Version 3 --- CURRENT}

\begin{itemize}[leftmargin=*]
\item Source G1, line 112: exact text.
\item Source LIVE, line 112: exact text.
\end{itemize}

\begin{lstlisting}
Large language models and coding agents are beginning to change this relationship. A scientist with deep drug-discovery expertise but limited software programming experience can increasingly describe an application, analysis, visualization, or workflow in natural language and have software constructed around the scientific need. Tool-using agents can operate established methods, build complex workflows, implement methods from the literature, handle failures, assemble results into insights that support decisions, and even optimize processes based on iterative feedback. Earlier systems such as ChemCrow and ChemGraph established proof-of-concept for language models coordinating chemistry tools and computational workflows \citep{bran2024chemcrow,boiko2023coscientist,pham2026chemgraph,macknight2025rethinking}. While their practical scope was limited compared with current coding agents, they helped establish the feasibility of coupling language models to executable scientific tools.
\end{lstlisting}

\subsection{R008: From fixed software products to generated adaptive software}

Baseline main.tex line 114; paragraph.

\paragraph{Version 1 --- earlier/alternative wording}

\begin{itemize}[leftmargin=*]
\item Source A, line 114: reviewed topic or figure-label match; not proof of edit lineage.
\item Source B, line 114: reviewed topic or figure-label match; not proof of edit lineage.
\end{itemize}

\begin{lstlisting}
The consequence is a shift in how drug hunters can interact with computational science. Scientists can increasingly shape software around the question, the stage of the project, the data available at that moment, and their preferred way of reasoning about the problem. A medicinal chemist optimizing a macrocycle may need a different interface and set of calculations than a structural biologist interrogating a cryptic pocket or a program team deciding which compound to advance. The interface can serve as a temporary, adaptive view over a more durable scientific substrate.
\end{lstlisting}

\paragraph{Version 2 --- earlier/alternative wording}

\begin{itemize}[leftmargin=*]
\item Source C, line 12: reviewed topic or figure-label match; not proof of edit lineage.
\end{itemize}

\begin{lstlisting}
The consequence is a shift in how drug hunters can interact with software. Scientists can increasingly customize software around project-specific questions, available data, and their preferred way of reasoning about a problem. A medicinal chemist optimizing a macrocycle may need a different interface and set of calculations than a structural biologist testing a mechanistic hypothesis or a program team deciding which compounds to advance from a virtual screening campaign. In this model, the interface becomes an adaptive and potentially transient view over a more durable scientific substrate. Changes to the interface do not imply changes to the underlying record of the work: molecular state, computational methods, parameters, intermediate results, user actions, and software versions are all persisted and linked through an auditable provenance history. This separation between a mutable presentation layer and a persistent scientific state is important for reproducibility, inspection, and reconstruction of prior decisions. Although Molarium presents as a relatively simple browser-based application, it is built on a substantially richer underlying architecture for managing versioned molecular state, tracking computational provenance, and executing adapted implementations of established molecular modeling algorithms efficiently within the constraints of a local web browser.
\end{lstlisting}

\paragraph{Version 3 --- CURRENT}

\begin{itemize}[leftmargin=*]
\item Source G1, line 114: exact text.
\item Source LIVE, line 114: exact text.
\end{itemize}

\begin{lstlisting}
The consequence is a shift in how drug hunters can interact with software. Scientists can increasingly customize software around project-specific questions, available data, and their preferred way of reasoning about a problem. A medicinal chemist optimizing a macrocycle may need a different interface and set of calculations than a structural biologist testing a mechanistic hypothesis or a program team deciding which compounds to advance from a virtual screening campaign. In this model, the interface becomes an adaptive and potentially transient view over a more durable scientific substrate. Changes to the interface do not imply changes to the underlying record of the work: molecular states, computational methods, parameters, intermediate results, user actions, and software versions should persist and link through an auditable provenance history. This separation between a mutable graphical layer and a persistent scientific state is important for reproducibility, inspection, and reconstruction of prior decisions. Although Molarium presents as a relatively simple browser-based application, it is built on a substantially richer underlying architecture for managing versioned molecular state, tracking computational provenance, and executing adapted implementations of established molecular modeling algorithms efficiently within the constraints of a local web browser.
\end{lstlisting}

\subsection{R009: From fixed software products to generated adaptive software}

Baseline main.tex line 116; paragraph.

\paragraph{Version 1 --- earlier/alternative wording}

\begin{itemize}[leftmargin=*]
\item Source A, line 116: inferred correspondence, similarity 0.71; verify before merging.
\item Source B, line 116: inferred correspondence, similarity 0.71; verify before merging.
\end{itemize}

\begin{lstlisting}
Molarium's development directly reflects this paradigm. It began as a need for a more effective environment for viewing and altering molecular structures and expanded through active use. Scientific challenges drove software changes: an unrealistic geometry led to refined editing semantics; questionable energies triggered explicit numerical reference checks; ambiguity about the execution backend led to clearer identification of the active engine; and unsupported methods were converted into explicit boundaries rather than unstated assumptions. The tool evolved alongside the scientific investigation. Appendix~\ref{app:architecture} provides the corresponding execution architecture, numerical validation, and executable design-history example. Appendix~\ref{app:manual} provides task-centered workflows for scientists and agents, including the relationship among skill contracts, public API actions, replay, and persistent Design History.
\end{lstlisting}

\paragraph{Version 2 --- earlier/alternative wording}

\begin{itemize}[leftmargin=*]
\item Source C, line 14: inferred correspondence, similarity 0.80; verify before merging.
\end{itemize}

\begin{lstlisting}
Molarium's development directly reflects this paradigm. It began as a need for a more effective environment for viewing and altering molecular structures and expanded through active use. Scientific challenges drove software changes: an unrealistic geometry led to refined editing semantics; questionable energies triggered explicit numerical reference checks; ambiguity about the execution backend led to clearer identification of the active engine; and unsupported methods were converted into explicit boundaries rather than unstated assumptions. The tool evolved alongside the scientific investigation. Appendix~\ref{app:architecture} provides the corresponding execution architecture, numerical validation, and executable design-history example. Appendix~\ref{app:manual} provides a module-by-module operating reference that distinguishes current, experimental, repository-only, proposed, and retired functionality.
\end{lstlisting}

\paragraph{Version 3 --- earlier/alternative wording}

\begin{itemize}[leftmargin=*]
\item Source D, line 3: inferred correspondence, similarity 0.93; verify before merging.
\end{itemize}

\begin{lstlisting}
Molarium’s development directly reflects this paradigm. It began as a need for a better environment to view and modify molecular structures, then evolved through active scientific use. Scientific problems drove software changes: unrealistic geometries led to refined editing semantics; questionable energies prompted explicit numerical reference checks; ambiguity about execution backends led to clearer identification of the active engine; and unsupported methods became explicit boundaries rather than implicit assumptions. The software evolved with the scientific questions it was built to address. Appendix~\ref{app} describes the execution architecture, numerical validation, and an executable example of this design history. Appendix~\ref{app} provides a module-by-module reference distinguishing current, experimental, repository-only, proposed, and retired functionality.
\end{lstlisting}

\paragraph{Version 4 --- CURRENT}

\begin{itemize}[leftmargin=*]
\item Source G1, line 116: exact text.
\item Source LIVE, line 116: exact text.
\end{itemize}

\begin{lstlisting}
Molarium’s development directly reflects this paradigm. It began as a need for a better environment to view and modify molecular structures, then evolved through active scientific use, all over the course of a weekend to create the first version. Scientific problems drove software changes: unrealistic geometries led to refined protocols; unconverged energies prompted explicit numerical reference checks; ambiguities about execution backends led to clearer identification of the active engine; and unvalidated methods became explicit boundaries rather than implicit assumptions. The software evolved with the scientific questions it was built to address. Appendix~\ref{app:architecture} describes the execution architecture, numerical validation, and an executable example of this design history. Appendix~\ref{app:manual} provides a module-by-module reference distinguishing current, experimental, repository-only, proposed, and retired functionality.
\end{lstlisting}

\subsection{R010: Molarium and the value of a persistent molecular state}

Baseline main.tex line 120; paragraph.

\paragraph{Version 1 --- earlier/alternative wording}

\begin{itemize}[leftmargin=*]
\item Source A, line 120: inferred correspondence, similarity 0.94; verify before merging.
\end{itemize}

\begin{lstlisting}
Molarium is an MIT-licensed, local-first molecular viewer, builder, and simulation workbench. Most supported calculations execute on the user's device using WebAssembly or WebGPU, allowing molecular graphics, editing, cheminformatics, conformer exploration, molecular mechanics, molecular dynamics, machine-learned potentials, and selected protein-structure prediction workflows to coexist in a standard browser environment. The browser reduces installation and orchestration overhead and, for supported calculations, can avoid remote job submission, cloud-compute charges, and application licensing while keeping computation inside the interactive design loop. WebAssembly provides a route for mature compiled scientific libraries to run locally, while WebGPU exposes modern accelerators through a portable browser API \citep{haas2017wasm,webgpu2025}.
\end{lstlisting}

\paragraph{Version 2 --- earlier/alternative wording}

\begin{itemize}[leftmargin=*]
\item Source B, line 120: inferred correspondence, similarity 0.97; verify before merging.
\end{itemize}

\begin{lstlisting}
Molarium is a local-first molecular viewer, builder, and simulation workbench released under the Apache License 2.0. Most supported calculations execute on the user's device using WebAssembly or WebGPU, allowing molecular graphics, editing, cheminformatics, conformer exploration, molecular mechanics, molecular dynamics, machine-learned potentials, and selected protein-structure prediction workflows to coexist in a standard browser environment. The browser reduces installation and orchestration overhead and, for supported calculations, can avoid remote job submission, cloud-compute charges, and application licensing while keeping computation inside the interactive design loop. WebAssembly provides a route for mature compiled scientific libraries to run locally, while WebGPU exposes modern accelerators through a portable browser API \citep{haas2017wasm,webgpu2025}.
\end{lstlisting}

\paragraph{Version 3 --- CURRENT}

\begin{itemize}[leftmargin=*]
\item Source G1, line 120: exact text.
\item Source LIVE, line 120: exact text.
\end{itemize}

\begin{lstlisting}
Molarium is a local-first molecular viewer, builder, and simulation workbench released under the Apache License 2.0. Calculations execute on the user's local device using WebAssembly or WebGPU, allowing molecular graphics, editing, cheminformatics, conformer exploration, molecular mechanics, molecular dynamics, machine-learned potentials, and selected protein-structure prediction workflows to coexist in a standard browser environment. The browser reduces installation and orchestration overhead and avoids remote job submission, cloud-compute charges, and application licensing while keeping computation inside the interactive design loop. WebAssembly provides a route for mature compiled scientific libraries to run locally, while WebGPU exposes modern accelerators through a portable browser API \citep{haas2017wasm,webgpu2025}.
\end{lstlisting}

\subsection{R011: Molarium and the value of a persistent molecular state}

Baseline main.tex line 122; structured block.

\paragraph{Version 1 --- earlier/alternative wording}

\begin{itemize}[leftmargin=*]
\item Source A, line 126: reviewed topic or figure-label match; not proof of edit lineage.
\end{itemize}

\begin{lstlisting}
\begin{figure}[htbp]
  \centering
  \safeincludegraphics[width=\linewidth]{figures/fig1_molarium_interface.png}
  \caption{The Molarium molecular studio during prospective BAY-293 pose generation from the frozen compound-21 state. The central view retains the local protein--ligand context and marks changed atoms; the left rail explains the active Designer Move; and the right rail reports 44 distinct poses, nine feasible candidates, and the eight-worker execution of the 64-chain search. View, Design, and Simulate reorganize controls around different scientific questions without replacing the authoritative molecular graph, coordinates, preparation state, provenance, or history beneath them.}
  \label{fig:views}
\end{figure}
\end{lstlisting}

\paragraph{Version 2 --- earlier/alternative wording}

\begin{itemize}[leftmargin=*]
\item Source B, line 126: reviewed topic or figure-label match; not proof of edit lineage.
\end{itemize}

\begin{lstlisting}
\begin{figure}[htbp]
  \centering
  \safeincludegraphics[width=\linewidth]{figures/fig1_molarium_interface.png}
  \caption{Molarium in its normal View workspace after browser-local Sage/CCD preparation of the experimental fragment F1-bound KRAS--SOS1 complex (PDB 6EPM). The official RCSB Chemical Component Dictionary topology supplies the BQ5/F1 bond orders and the chemically correct 2D depiction shown at lower left. The complete deposited protein assembly is shown as cartoons, while BQ5/F1 and protein residues within 5~\AA{} are rendered atomically to localize the allosteric design pocket. Preparation excluded crystallographic water; glycerol and display hydrogens are hidden for clarity. This scene provides structural context; the prospective hit-to-BAY-293 trajectory in Fig.~\ref{fig:sos1story} began independently from the 5OVE/AXE structure alone.}
  \label{fig:views}
\end{figure}
\end{lstlisting}

\paragraph{Version 3 --- earlier/alternative wording}

\begin{itemize}[leftmargin=*]
\item Source E, line 3: inferred correspondence, similarity 0.91; verify before merging.
\end{itemize}

\begin{lstlisting}
\caption{Molarium in its normal View workspace after browser-local Sage/CCD preparation of the experimental fragment F1-bound KRAS--SOS1 complex (PDB 6EPM). The official RCSB Chemical Component Dictionary topology supplies the BQ5/F1 bond orders and the chemically correct 2D depiction shown at lower left. The complete deposited protein assembly is shown as cartoons, while BQ5/F1 and protein residues within 5~\AA{} are rendered atomically to localize the allosteric design pocket. Sovlent, cofactors, and hydrogens are hidden for clarity. This scene provides structural context; the prospective hit-to-BAY-293 trajectory in Fig.~\ref{fig:sos1story} began independently from the 5OVE/AXE structure alone.}
\end{lstlisting}

\paragraph{Version 4 --- CURRENT}

\begin{itemize}[leftmargin=*]
\item Source G1, line 122: exact text.
\item Source LIVE, line 122: exact text.
\end{itemize}

\begin{lstlisting}
\begin{figure}[!htbp]
  \centering
  \safeincludegraphics[width=0.90\linewidth]{figures/fig1_molarium_interface.png}
  \caption{Molarium in its normal View workspace after browser-local Sage/CCD preparation of the experimental fragment F1-bound KRAS--SOS1 complex (PDB 6EPM). The official RCSB Chemical Component Dictionary topology supplies the BQ5/F1 bond orders and the chemically correct 2D depiction shown at lower left. The complete deposited protein assembly is shown as cartoons, while BQ5/F1 and protein residues within 5~\AA{} are rendered atomically to localize the allosteric design pocket. Solvent, cofactors, and hydrogens are hidden for clarity. This scene provides structural context; the prospective hit-to-BAY-293 trajectory in Fig.~\ref{fig:sos1story} began independently from the 5OVE/AXE structure alone.}
  \label{fig:views}
\end{figure}
\end{lstlisting}

\subsection{R012: Molarium and the value of a persistent molecular state}

Baseline main.tex line 131; paragraph.

\paragraph{Version 1 --- earlier/alternative wording}

\begin{itemize}[leftmargin=*]
\item Source A, line 122: inferred correspondence, similarity 0.95; verify before merging.
\item Source B, line 122: inferred correspondence, similarity 0.95; verify before merging.
\end{itemize}

\begin{lstlisting}
The architectural choice that matters most is the preservation of a common scientific state. Molarium maintains explicit topology, coordinates, protonation state, chemical edits, preparation history, trajectories, model identity, and provenance as properties of the molecular session rather than allowing each view or calculation to construct its own interpretation of the molecule. The interface can therefore change rapidly without redefining the scientific object beneath it.
\end{lstlisting}

\paragraph{Version 2 --- CURRENT}

\begin{itemize}[leftmargin=*]
\item Source G1, line 131: exact text.
\item Source LIVE, line 131: exact text.
\end{itemize}

\begin{lstlisting}
One of the advanced architectural capabilities is the preservation of a common scientific state. Molarium maintains explicit topology, coordinates, protonation state, chemical edits, preparation history, trajectories, model identity, and provenance as properties of the molecular session rather than allowing each view or calculation to construct its own interpretation of the molecule. The interface can therefore change rapidly without redefining the scientific object beneath it.
\end{lstlisting}

\subsection{R013: Molarium and the value of a persistent molecular state}

Baseline main.tex line 133; paragraph.

\paragraph{Version 1 --- earlier/alternative wording}

\begin{itemize}[leftmargin=*]
\item Source A, line 124: inferred correspondence, similarity 0.94; verify before merging.
\item Source B, line 124: inferred correspondence, similarity 0.94; verify before merging.
\end{itemize}

\begin{lstlisting}
Three task-specific views illustrate this separation. View emphasizes the molecular system, including protein, ligand, waters, ions, contacts, and selected residues. Design exposes chemical and geometrical interventions. Simulate emphasizes preparation, trajectories, and analysis. The controls differ because the scientific questions differ; topology, coordinates, method identity, and history remain attached to the same evolving molecular state. Detailed implementation and replay examples are provided in Appendix~\ref{app:architecture}, while the state mutations, exports, and failure behaviour of the corresponding user-facing modules are specified in Appendix~\ref{app:manual}.
\end{lstlisting}

\paragraph{Version 2 --- CURRENT}

\begin{itemize}[leftmargin=*]
\item Source G1, line 133: exact text.
\item Source LIVE, line 133: exact text.
\end{itemize}

\begin{lstlisting}
Three task-specific views illustrate this separation. ``View'' emphasizes the molecular system, including molecular representations of protein, ligand, waters, ions, contacts, measurements, and selected residues. ``Design'' exposes chemical and geometrical modifications. ``Simulate'' emphasizes preparation, simulation, and analysis. The controls differ because the scientific questions differ; topology, coordinates, method identity, and history remain attached to the same evolving molecular state. Detailed implementation and replay examples are provided in Appendix~\ref{app:architecture}, while the state mutations, exports, and failure behaviour of the corresponding user-facing modules are specified in Appendix~\ref{app:manual}.
\end{lstlisting}

\subsection{R014: Molarium and the value of a persistent molecular state}

Baseline main.tex line 135; structured block.

\paragraph{Version 1 --- earlier/alternative wording}

\begin{itemize}[leftmargin=*]
\item Source A, line 133: reviewed topic or figure-label match; not proof of edit lineage.
\end{itemize}

\begin{lstlisting}
\begin{figure}[htbp]
  \centering
  \safeincludegraphics[width=\linewidth]{figures/fig2_sos1_hit_to_bay293.png}
  \caption{Hit-only SOS1 design and post-freeze structural evaluation. (A) The prospective coordinate boundary was restricted to PDB 5OVE with its AXE ligand; structures 5OVF--5OVI remained unopened. (B) Sequential scaffold rewriting and growth to compound 21 reached the Phe890-in volume and introduced four steric clashes. (C) A three-basin side-chain search selected the Phe890-out branch, which introduced no additional clashes and moved $\chi_1$ from approximately $-73^{\circ}$ to $-172^{\circ}$. (D) The resulting open-pocket state was propagated to BAY-293 and frozen before 5OVI was opened. Post-freeze comparison with 5OVI gave a receptor-aligned, graph-symmetry-minimized ligand RMSD of 1.215~\AA{}; Phe890 $\chi_1$ was $-172.43^{\circ}$ in the prediction and $-166.55^{\circ}$ in the crystal. Later crystal structures evaluated the frozen trajectory but did not seed it.}
  \label{fig:sos1story}
\end{figure}
\end{lstlisting}

\paragraph{Version 2 --- CURRENT}

\begin{itemize}[leftmargin=*]
\item Source B, line 133: exact text.
\item Source G1, line 135: exact text.
\item Source LIVE, line 135: exact text.
\end{itemize}

\begin{lstlisting}
\begin{figure}[htbp]
  \centering
  \safeincludegraphics[width=\linewidth]{figures/fig2_sos1_hit_to_bay293.png}
  \caption{Hit-only SOS1 design replayed in the Molarium interface through the public Chemist Actions API. (A) The prepared 5OVE/AXE hit (teal) and its pale local pocket are centered before prospective design; 5OVE/AXE was the sole coordinate-bearing input. The camera established here remains fixed. Between panels A and B, the replay labels Phe890 and Lys898 while the first predicted pyrazole is inspected. (B) Sequential scaffold rewriting and growth produces compound 21 (purple) in the Phe890-in pocket. The cyan halo identifies the edited ligand region; gold identifies Phe890; magenta dotted connectors are severe ligand--protein short contacts recalculated from the live coordinates. (C) A discrete side-chain search rotates Phe890 out of the growth path. Only the six moved Phe890 heavy atoms are haloed, so the conformational change can be compared directly with B without a camera refit. (D) The final predicted BAY-293 complex is shown after coupled WebGPU relaxation with transient change marks cleared. The replay ends at this prediction and never opens 5OVI. Independent post-freeze comparison with 5OVI subsequently gave a receptor-aligned, graph-symmetry-minimized ligand RMSD of 1.215~\AA{} and Phe890 $\chi_1$ values of $-172.43^{\circ}$ in the prediction and $-166.55^{\circ}$ in the crystal; those validation coordinates did not seed the replay.}
  \label{fig:sos1story}
\end{figure}
\end{lstlisting}

\subsection{R015: Molarium and the value of a persistent molecular state}

Baseline main.tex line 142; paragraph.

\paragraph{Version 1 --- earlier/alternative wording}

\begin{itemize}[leftmargin=*]
\item Source A, line 140: reviewed topic or figure-label match; not proof of edit lineage.
\end{itemize}

\begin{lstlisting}
The executable replay is available at \url{https://molarium.org/sos1-hit-to-bay293}. A fixed 1600$\times$1000 interface rendering of the same 48-move public-API replay is available at \url{https://molarium.org/assets/media/sos1-designer-moves-molarium-interface.mp4}.
\end{lstlisting}

\paragraph{Version 2 --- CURRENT}

\begin{itemize}[leftmargin=*]
\item Source B, line 140: exact text.
\item Source G1, line 142: exact text.
\item Source LIVE, line 142: exact text.
\end{itemize}

\begin{lstlisting}
The executable replay is available at \url{https://molarium.org/sos1-hit-to-bay293}.
\end{lstlisting}

\subsection{R017: Molarium and the value of a persistent molecular state}

Baseline main.tex line 146; paragraph.

\paragraph{Version 1 --- earlier/alternative wording}

\begin{itemize}[leftmargin=*]
\item Source A, line 144: inferred correspondence, similarity 0.99; verify before merging.
\item Source B, line 144: inferred correspondence, similarity 0.99; verify before merging.
\end{itemize}

\begin{lstlisting}
Molarium extends this principle by preserving the sequence of molecular states and user actions that produced the current session. Edits, selections, manipulations, calculations, and other interactions form an event history rather than leaving only the final coordinates. The resulting architecture is closer to version control than to a traditional molecular project file. Each molecular state can be associated with the actions that produced it, the calculation applied to it, and the evidence available at that point in the design process.
\end{lstlisting}

\paragraph{Version 2 --- CURRENT}

\begin{itemize}[leftmargin=*]
\item Source G1, line 146: exact text.
\item Source LIVE, line 146: exact text.
\end{itemize}

\begin{lstlisting}
Molarium extends this principle by preserving the sequence of molecular states and user actions that produced the current state. Edits, selections, manipulations, calculations, and other interactions form an event history rather than leaving only the final coordinates. The resulting architecture is closer to version control than to a traditional molecular project file. Each molecular state can be associated with the actions that produced it, the calculation applied to it, and the evidence available at that point in the design process.
\end{lstlisting}

\subsection{R018: Molarium and the value of a persistent molecular state}

Baseline main.tex line 148; paragraph.

\paragraph{Version 1 --- earlier/alternative wording}

\begin{itemize}[leftmargin=*]
\item Source A, line 146: inferred correspondence, similarity 0.93; verify before merging.
\item Source B, line 146: inferred correspondence, similarity 0.93; verify before merging.
\end{itemize}

\begin{lstlisting}
The ``Git for molecules'' analogy is useful, although a drug-design history is richer than source code. A design trajectory includes structural changes together with evolving hypotheses, model outputs, uncertainty, and experimental evidence. In a discovery setting, such a record could capture which analogs were considered, which structures or properties influenced a design, which alternatives were rejected, and what was learned when a compound was synthesized and tested. This creates a foundation for recording medicinal-chemistry design trajectories rather than storing only the compounds that survived them. Much of the information that explains how a program learned lies in alternatives that were considered and abandoned, yet conventional compound databases usually preserve endpoints.
\end{lstlisting}

\paragraph{Version 2 --- CURRENT}

\begin{itemize}[leftmargin=*]
\item Source G1, line 148: exact text.
\item Source LIVE, line 148: exact text.
\end{itemize}

\begin{lstlisting}
Thia concept is analogous to ``Git for molecules'', although a drug design history is richer and more complex than tracking revisions of source code. A design history includes structural changes together with evolving hypotheses, model outputs, uncertainty, and experimental evidence. In a discovery setting, such a record could capture which analogs were considered, which structures or properties influenced a design, which alternatives were rejected, and what was learned when a compound was synthesized and tested. This creates a foundation for recording medicinal chemistry design trajectories rather than storing only the compounds that survived them. Much of the information that explains how a program learned lies in alternatives that were considered and abandoned, yet conventional compound databases usually only preserve experimental endpoints.
\end{lstlisting}

\subsection{R020: Molarium and the value of a persistent molecular state}

Baseline main.tex line 159; paragraph.

\paragraph{Version 1 --- earlier/alternative wording}

\begin{itemize}[leftmargin=*]
\item Source A, line 155: inferred correspondence, similarity 1.00; verify before merging.
\item Source B, line 155: inferred correspondence, similarity 1.00; verify before merging.
\end{itemize}

\begin{lstlisting}
The browser also provides a useful trust boundary for proprietary work. Molarium's Local Lab mode can restrict network access and verify reviewed application assets while keeping supported calculations on the local device. Because privacy controls depend on deployment configuration, this architecture makes data movement explicit and auditable. Organizations can evaluate and enforce how molecular data are handled rather than treating transfer to remote infrastructure as an invisible property of the software stack. Appendix~\ref{app:manual} specifies the connected and Local Lab network policies and makes clear that local calculation and fully offline operation are distinct claims.
\end{lstlisting}

\paragraph{Version 2 --- CURRENT}

\begin{itemize}[leftmargin=*]
\item Source G1, line 159: exact text.
\item Source LIVE, line 159: exact text.
\end{itemize}

\begin{lstlisting}
The browser also provides a useful trust boundary for proprietary work. Molarium's ``Local Lab'' mode can restrict network access and verify reviewed application assets while keeping supported calculations on the local device. Because privacy controls depend on deployment configuration, this architecture makes data movement explicit and auditable. Organizations can evaluate and enforce how molecular data are handled rather than treating transfer to remote infrastructure as an invisible property of the software stack. Appendix~\ref{app:manual} specifies the connected and Local Lab network policies and makes clear that local calculation and fully offline operation are distinct claims.
\end{lstlisting}

\subsection{R021: Scientific contracts: provenance, validation, and falsifiability}

Baseline main.tex line 163; paragraph.

\paragraph{Version 1 --- earlier/alternative wording}

\begin{itemize}[leftmargin=*]
\item Source A, line 159: inferred correspondence, similarity 0.91; verify before merging.
\item Source B, line 159: inferred correspondence, similarity 0.91; verify before merging.
\end{itemize}

\begin{lstlisting}
The emergence of cheap, adaptive software introduces a scientific vulnerability. Agentically generated workflows become especially risky when seamless execution obscures the boundaries of the underlying scientific contract. As the time required to build a custom analysis view drops to minutes, that view still needs to expose the molecule, model, parameters, assumptions, and evidence that produced the result. Reducing the friction of software development without those safeguards increases the rate at which scientific ambiguity can be created.
\end{lstlisting}

\paragraph{Version 2 --- CURRENT}

\begin{itemize}[leftmargin=*]
\item Source G1, line 163: exact text.
\item Source LIVE, line 163: exact text.
\end{itemize}

\begin{lstlisting}
The emergence of rapidly adaptive software introduces a scientific vulnerability. Agentically generated workflows become especially risky when ease-of-use obscures the underlying scientific methods. As the time required to build a custom analysis view drops, that view still needs to expose the molecule, model, parameters, assumptions, and evidence that produced the result. Reducing the friction of software development without those safeguards increases the rate at which scientific ambiguity and errors can be introduced.
\end{lstlisting}

\subsection{R022: Scientific contracts: provenance, validation, and falsifiability}

Baseline main.tex line 165; paragraph.

\paragraph{Version 1 --- earlier/alternative wording}

\begin{itemize}[leftmargin=*]
\item Source A, line 161: inferred correspondence, similarity 1.00; verify before merging.
\item Source B, line 161: inferred correspondence, similarity 1.00; verify before merging.
\end{itemize}

\begin{lstlisting}
We use the term scientific contract for the collection of invariants and evidence that gives a computational output its meaning. A molecular-mechanics result cannot be defined by the generic label ``molecular mechanics.'' Its identity depends on the molecular state, charge assignment, parameterization, energy function, treatment of nonbonded interactions, boundary conditions, constraints, integrator, numerical implementation, and precision. A learned potential depends on its architecture, weights, preprocessing, supported chemistry, and training domain. A protein-structure prediction depends on checkpoint selection, sequence and template information, inference protocol, and other model-specific choices. The surrounding interface may adapt dynamically, but these scientific identities must remain explicit and protected from silent drift.
\end{lstlisting}

\paragraph{Version 2 --- CURRENT}

\begin{itemize}[leftmargin=*]
\item Source G1, line 165: exact text.
\item Source LIVE, line 165: exact text.
\end{itemize}

\begin{lstlisting}
We use the term ``scientific contract'' for the collection of invariants and evidence that gives a computational output its meaning. A molecular-mechanics result cannot be defined by the generic label ``molecular mechanics.'' Its identity depends on the molecular state, charge assignment, parameterization, energy function, treatment of nonbonded interactions, boundary conditions, constraints, integrator, numerical implementation, and precision. A learned potential depends on its architecture, weights, preprocessing, supported chemistry, and training domain. A protein-structure prediction depends on checkpoint selection, sequence and template information, inference protocol, and other model-specific choices. The surrounding interface may adapt dynamically, but these scientific identities must remain explicit and protected from silent drift.
\end{lstlisting}

\subsection{R023: Scientific contracts: provenance, validation, and falsifiability}

Baseline main.tex line 167; paragraph.

\paragraph{Version 1 --- earlier/alternative wording}

\begin{itemize}[leftmargin=*]
\item Source A, line 163: inferred correspondence, similarity 1.00; verify before merging.
\item Source B, line 163: inferred correspondence, similarity 1.00; verify before merging.
\end{itemize}

\begin{lstlisting}
Molarium preserves distinctions among computational pathways even when they operate in the same workspace. RDKit supports chemical perception and conformer generation; OpenFF Sage defines a molecular-mechanics model; OpenMM Reference provides an independent numerical route for selected validation tests; STORMM-inspired WebGPU code advances homogeneous replicas; ANI-2x provides a machine-learned potential with a defined chemical domain; and OpenFold supports a separate protein-structure prediction pathway \citep{riniker2015etkdg,boothroyd2023sage,eastman2024openmm8,cerutti2024stormm,devereux2020ani2x,ahdritz2024openfold}. These methods can be compared or chained without implying that their energies, assumptions, or domains are interchangeable. Appendix~\ref{app:architecture} specifies the implemented numerical pathways and their validation boundaries; Appendix~\ref{app:manual} records how those pathways are exposed, invoked, and constrained in the current workbench.
\end{lstlisting}

\paragraph{Version 2 --- earlier/alternative wording}

\begin{itemize}[leftmargin=*]
\item Source E, line 6: inferred correspondence, similarity 1.00; verify before merging.
\end{itemize}

\begin{lstlisting}
Molarium preserves distinctions among computational pathways even when they operate in the same workspace. The RDKit supports chemical perception and conformer generation; OpenFF Sage defines a molecular-mechanics model; OpenMM Reference provides an independent numerical route for selected validation tests; STORMM-inspired WebGPU code advances homogeneous replicas; ANI-2x provides a machine-learned potential with a defined chemical domain; and OpenFold3 supports a separate protein-structure prediction pathway \citep{riniker2015etkdg,boothroyd2023sage,eastman2024openmm8,cerutti2024stormm,devereux2020ani2x,ahdritz2024openfold}. These methods can be compared or chained without implying that their energies, assumptions, or domains are interchangeable. Appendix~\ref{app:architecture} specifies the implemented numerical pathways and their validation boundaries; Appendix~\ref{app:manual} records how those pathways are exposed, invoked, and constrained in the current workbench.
\end{lstlisting}

\paragraph{Version 3 --- CURRENT}

\begin{itemize}[leftmargin=*]
\item Source G1, line 167: exact text.
\item Source LIVE, line 167: exact text.
\end{itemize}

\begin{lstlisting}
Molarium preserves distinctions among computational pathways even when they operate in the same workspace. RDKit supports chemical perception and conformer generation; OpenFF Sage defines a molecular-mechanics model; OpenMM Reference provides an independent numerical route for selected validation tests; STORMM-inspired WebGPU code advances homogeneous replicas; ANI-2x provides a machine-learned potential with a defined chemical domain; and OpenFold3 supports a separate protein-structure prediction pathway \citep{riniker2015etkdg,boothroyd2023sage,eastman2024openmm8,cerutti2024stormm,devereux2020ani2x,ahdritz2024openfold}. These methods can be compared or chained without implying that their energies, assumptions, or domains are interchangeable. Appendix~\ref{app:architecture} specifies the implemented numerical pathways and their validation boundaries; Appendix~\ref{app:manual} records how those pathways are exposed, invoked, and constrained in the current workbench.
\end{lstlisting}

\subsection{R024: Scientific contracts: provenance, validation, and falsifiability}

Baseline main.tex line 169; paragraph.

\paragraph{Version 1 --- earlier/alternative wording}

\begin{itemize}[leftmargin=*]
\item Source A, line 165: inferred correspondence, similarity 0.98; verify before merging.
\item Source B, line 165: inferred correspondence, similarity 0.98; verify before merging.
\end{itemize}

\begin{lstlisting}
The same principle applies to workflow provenance. Reproducible computational science requires retention of the data, computational process, execution environment, and provenance that produced a result \citep{cohenboulakia2017workflows,wilkinson2016fair}. More recent work formalizes workflow invariants and testable reproducibility criteria rather than treating provenance as passive metadata \citep{pritchard2025reproducibility}. Generated adaptive software increases the importance of this requirement because the workflow itself may be created or modified during the scientific interaction rather than existing as a stable artifact defined in advance.
\end{lstlisting}

\paragraph{Version 2 --- CURRENT}

\begin{itemize}[leftmargin=*]
\item Source G1, line 169: exact text.
\item Source LIVE, line 169: exact text.
\end{itemize}

\begin{lstlisting}
The same principle applies to workflow provenance. Reproducible computational science requires retention of the data, computational process, execution environment, and provenance that produced a result \citep{cohenboulakia2017workflows,wilkinson2016fair}. More recent work formalizes workflow invariants and testable reproducibility criteria rather than treating provenance as passive metadata \citep{pritchard2025reproducibility}. Generated adaptive software increases the importance of this requirement because the workflow itself may be created or modified during the scientific interaction rather than existing as a single release version defined in advance.
\end{lstlisting}

\subsection{R027: Scientific contracts: provenance, validation, and falsifiability}

Baseline main.tex line 180; paragraph.

\paragraph{Version 1 --- earlier/alternative wording}

\begin{itemize}[leftmargin=*]
\item Source A, line 176: inferred correspondence, similarity 0.98; verify before merging.
\item Source B, line 176: inferred correspondence, similarity 0.98; verify before merging.
\end{itemize}

\begin{lstlisting}
Molarium's development repeatedly followed this pattern. A chemically implausible geometry exposed a problem in representation and editing semantics. A questionable energy prompted comparison with an independent implementation. A discrepancy between CPU and GPU execution revealed both a unit error and an incorrect assumption about which backend was running. Requesting constrained dynamics led to implementation of SHAKE/RATTLE and comparison against a reference pathway \citep{ryckaert1977shake,andersen1983rattle}. In each case, the useful scientific result consisted of the feature together with a test that established the boundary of the claim.
\end{lstlisting}

\paragraph{Version 2 --- CURRENT}

\begin{itemize}[leftmargin=*]
\item Source G1, line 180: exact text.
\item Source LIVE, line 180: exact text.
\end{itemize}

\begin{lstlisting}
Molarium's development repeatedly followed this pattern. A chemically implausible geometry exposed a problem in representation and editing semantics and a questionable energy prompted comparison with an independent implementation. Similarly, a discrepancy between CPU and GPU execution revealed both a unit error and an incorrect assumption about which backend was running. Requesting constrained dynamics led to implementation of SHAKE/RATTLE and comparison against a reference pathway \citep{ryckaert1977shake,andersen1983rattle}. In each case, the useful scientific result consisted of the feature together with a test that established the boundary of the claim.
\end{lstlisting}

\subsection{R028: Scientific contracts: provenance, validation, and falsifiability}

Baseline main.tex line 182; paragraph.

\paragraph{Version 1 --- earlier/alternative wording}

\begin{itemize}[leftmargin=*]
\item Source A, line 178: inferred correspondence, similarity 0.78; verify before merging.
\item Source B, line 178: inferred correspondence, similarity 0.78; verify before merging.
\end{itemize}

\begin{lstlisting}
Fluent agentic software makes this distinction more consequential. A unit error can produce a smooth trajectory. Separate implementations can agree if they inherit the same flawed inputs. A numerically correct implementation can apply an inappropriate physical model. Successful execution and an attractive visualization therefore occupy the bottom of an evidence hierarchy rather than establishing that a molecular prediction is correct.
\end{lstlisting}

\paragraph{Version 2 --- CURRENT}

\begin{itemize}[leftmargin=*]
\item Source G1, line 182: exact text.
\item Source LIVE, line 182: exact text.
\end{itemize}

\begin{lstlisting}
Fluent agentic software makes this distinction more consequential, since problems are not always obvious and do not necessarily result in errors. For example, a unit error can still produce a smooth trajectory and independent implementations can agree if they inherit the same flawed inputs. Successful execution and an attractive visualization therefore occupy the bottom of an evidence hierarchy rather than establishing that a molecular prediction is correct.
\end{lstlisting}

\subsection{R031: Local computation changes the interaction model}

Baseline main.tex line 195; paragraph.

\paragraph{Version 1 --- earlier/alternative wording}

\begin{itemize}[leftmargin=*]
\item Source A, line 191: inferred correspondence, similarity 0.75; verify before merging.
\item Source B, line 191: inferred correspondence, similarity 0.75; verify before merging.
\end{itemize}

\begin{lstlisting}
A web browser does not replace production molecular simulation software or high-performance computing. Large periodic systems, long trajectories, advanced sampling, specialized force fields, and large-scale parallel studies can require substantially more compute than a personal machine can provide. Molarium addresses a different part of the workflow: calculations that are sufficiently bounded to remain inside the scientist's interactive design session.
\end{lstlisting}

\paragraph{Version 2 --- CURRENT}

\begin{itemize}[leftmargin=*]
\item Source G1, line 195: exact text.
\item Source LIVE, line 195: exact text.
\end{itemize}

\begin{lstlisting}
A web browser cannot always replace production molecular simulation software or high-performance computing. Large periodic systems, long trajectories, and large-scale screening studies can require substantially more compute resouirces than a personal machine can provide. Molarium currently addresses calculations that are sufficiently bounded to remain inside the scientist's interactive design session, although remote job submission to a HPC cluster or the cloud for specific calculations would be trival to add.
\end{lstlisting}

\subsection{R032: Local computation changes the interaction model}

Baseline main.tex line 197; paragraph.

\paragraph{Version 1 --- earlier/alternative wording}

\begin{itemize}[leftmargin=*]
\item Source A, line 193: inferred correspondence, similarity 0.88; verify before merging.
\item Source B, line 193: inferred correspondence, similarity 0.88; verify before merging.
\end{itemize}

\begin{lstlisting}
The practical benefit is reduced decision latency. Consider a medicinal chemist examining a bound ligand and asking whether a proposed three-dimensional edit creates an avoidable steric problem or opens a plausible alternative conformation. In a fragmented workflow, the scientist may export coordinates, request or run a separate calculation, move results back into a viewer, and reconstruct the context in which the question arose. In Molarium, supported editing, local relaxation, conformational exploration, and trajectory inspection can occur in the same session. The value comes from keeping the molecular state, the calculation, and the design question connected while the decision is still active.
\end{lstlisting}

\paragraph{Version 2 --- CURRENT}

\begin{itemize}[leftmargin=*]
\item Source G1, line 197: exact text.
\item Source LIVE, line 197: exact text.
\end{itemize}

\begin{lstlisting}
The practical benefit of local compute is the reduced decision latency. Consider a medicinal chemist examining a bound ligand and asking whether a proposed three-dimensional edit creates an avoidable steric problem or opens a plausible alternative conformation of the protien binding pocket. In a fragmented workflow, the scientist may export coordinates, request or run a separate calculation, move results back into a viewer, and reconstruct the context in which the question arose. In Molarium the editing, relaxation, conformational exploration, and trajectory inspection all run locally within the same session, tracking the molecular state, the calculation, associated parameters, the design question, and decisions.
\end{lstlisting}

\subsection{R034: Local computation changes the interaction model}

Baseline main.tex line 201; paragraph.

\paragraph{Version 1 --- earlier/alternative wording}

\begin{itemize}[leftmargin=*]
\item Source A, line 197: inferred correspondence, similarity 0.94; verify before merging.
\item Source B, line 197: inferred correspondence, similarity 0.94; verify before merging.
\end{itemize}

\begin{lstlisting}
The lower operational barrier increases the importance of choosing calculations deliberately. A medicinal chemist deciding between two substituent orientations may gain value from a constrained local refinement that changes which analog is synthesized next; running hundreds of additional simulations after the decision is already insensitive to the result adds computational activity without increasing scientific leverage. The operational gain matters when it helps the scientist focus on the right question, apply the appropriate method, and prioritize molecules that have a realistic chance to advance the program.
\end{lstlisting}

\paragraph{Version 2 --- CURRENT}

\begin{itemize}[leftmargin=*]
\item Source G1, line 201: exact text.
\item Source LIVE, line 201: exact text.
\end{itemize}

\begin{lstlisting}
The lower operational barrier increases the importance of defining the project questions appropriately and choosing calculations deliberately. A medicinal chemist deciding between two substituent orientations may gain value from a constrained local refinement that changes which analog is synthesized next while running hundreds of additional simulations after the decision is already insensitive to the result adds computational activity without increasing scientific leverage. The operational efficiency helps the scientist focus on the right question, apply the appropriate method, and prioritize molecules that have a realistic chance to advance the program.
\end{lstlisting}

\subsection{R035: Local computation changes the interaction model}

Baseline main.tex line 203; paragraph.

\paragraph{Version 1 --- earlier/alternative wording}

\begin{itemize}[leftmargin=*]
\item Source A, line 199: inferred correspondence, similarity 0.95; verify before merging.
\item Source B, line 199: inferred correspondence, similarity 0.95; verify before merging.
\end{itemize}

\begin{lstlisting}
Local execution also redistributes control. The scientist can determine which analysis should exist, how data should be juxtaposed, and which calculation should be available at the point of decision. Molarium illustrates this shift: a scientist did not need to wait for every useful workflow to appear on a conventional product roadmap. Missing interfaces and bounded implementations could be added while the underlying question was still active.
\end{lstlisting}

\paragraph{Version 2 --- CURRENT}

\begin{itemize}[leftmargin=*]
\item Source G1, line 203: exact text.
\item Source LIVE, line 203: exact text.
\end{itemize}

\begin{lstlisting}
Local execution also redistributes control. The scientist can determine which analysis should exist, how data should be juxtaposed, and which calculation should be available at the point of decision. Molarium illustrates this shift: a scientist did not need to wait for every useful workflow to appear on a conventional product roadmap. Missing interfaces, algorithms, andmethods can be added while the underlying question was still active.
\end{lstlisting}

\subsection{R060: Conclusion}

Baseline main.tex line 266; paragraph.

\paragraph{Version 1 --- earlier/alternative wording}

\begin{itemize}[leftmargin=*]
\item Source A, line 262: inferred correspondence, similarity 0.97; verify before merging.
\end{itemize}

\begin{lstlisting}
\callout{The tool layer is becoming cheap and adaptive. Scientific advantage moves toward better models, better discovery intelligence, and better decisions.}
\end{lstlisting}

\paragraph{Version 2 --- CURRENT}

\begin{itemize}[leftmargin=*]
\item Source B, line 262: exact text.
\item Source G1, line 266: exact text.
\item Source LIVE, line 266: exact text.
\end{itemize}

\begin{lstlisting}
The tool layer is becoming cheap and adaptive. Scientific advantage moves toward better models, better discovery intelligence, and better decisions.
\end{lstlisting}

\subsection{R062: AI assistance and availability}

Baseline main.tex line 272; paragraph.

\paragraph{Version 1 --- earlier/alternative wording}

\begin{itemize}[leftmargin=*]
\item Source A, line 268: inferred correspondence, similarity 0.89; verify before merging.
\end{itemize}

\begin{lstlisting}
Molarium is available at \url{https://molarium.org}. Its source code and technical documentation are publicly available at \url{https://github.com/rafwiewiora/molarium}. Quantitative benchmark results are scoped by the relevant implementation state, fixtures, model and asset hashes, browser and hardware environment, numerical precision, and timing protocol. Appendix~\ref{app:architecture} provides the architecture, numerical validation, and executable design-history evidence underlying these results; Appendix~\ref{app:manual} records the inspected operating surface, module boundaries, failure behaviour, and export/persistence contracts.
\end{lstlisting}

\paragraph{Version 2 --- CURRENT}

\begin{itemize}[leftmargin=*]
\item Source B, line 268: exact text.
\item Source G1, line 272: exact text.
\item Source LIVE, line 272: exact text.
\end{itemize}

\begin{lstlisting}
Molarium is available at \url{https://molarium.org}. Its source code and technical documentation are publicly available under the Apache License 2.0 at \url{https://github.com/rafwiewiora/molarium}. Bundled third-party software, model parameters, force fields, and scientific data retain their separately identified terms. Quantitative benchmark results are scoped by the relevant implementation state, fixtures, model and asset hashes, browser and hardware environment, numerical precision, and timing protocol. Appendix~\ref{app:architecture} provides the architecture, numerical validation, and executable design-history evidence underlying these results; Appendix~\ref{app:manual} records the inspected operating surface, module boundaries, failure behaviour, and export/persistence contracts.
\end{lstlisting}

\subsection{R096: Molarium architecture, numerical validation, and executable design history}

Baseline main.tex line 421; paragraph.

Subsection: Workflow orchestration and local execution boundary.

\paragraph{Version 1 --- earlier/alternative wording}

\begin{itemize}[leftmargin=*]
\item Source A, line 417: inferred correspondence, similarity 1.00; verify before merging.
\item Source B, line 417: inferred correspondence, similarity 1.00; verify before merging.
\end{itemize}

\begin{lstlisting}
\item \textbf{Local small-molecule design.} Synchronized 2D/3D construction, graph validation, RDKit or direct-WebGPU calculations, aligned trajectory/energy inspection, and export. OpenMM/WebAssembly is used for internal validation and fixed-scaffold operations rather than as a selectable Simulate-panel method.
\item \textbf{Protein--ligand pose propagation.} A trusted pose is captured with stable atom lineage; changed ligand branches are sampled while the receptor remains fixed; optional contact hypotheses constrain search; and a selected pose can be applied and relaxed. An experimental induced-fit action releases the ligand and all atoms of protein residues within a 6 \AA{} pocket shell while outer atoms remain fixed.
\item \textbf{Agent-operated design.} Agents invoke the versioned Chemist Actions API rather than arbitrary page functions. The API exposes the same bounded state-changing verbs available to a chemist---including select, edit, search, apply, relax, and undo---plus read-only inspection. Requests execute serially and record sequence, timestamps, arguments, status, and duration.
\end{lstlisting}

\paragraph{Version 2 --- CURRENT}

\begin{itemize}[leftmargin=*]
\item Source G1, line 421: exact text.
\item Source LIVE, line 421: exact text.
\end{itemize}

\begin{lstlisting}
\item \textbf{Local small-molecule design.} Synchronized 2D/3D construction, graph validation, RDKit or direct-WebGPU calculations, aligned trajectory/energy inspection, and export. OpenMM/WebAssembly is used for internal validation and fixed-scaffold operations rather than as a selectable Run-panel method.
\item \textbf{Protein--ligand pose propagation.} A trusted pose is captured with stable atom lineage; changed ligand branches are sampled while the receptor remains fixed; optional contact hypotheses constrain search; and a selected pose can be applied and relaxed. An experimental induced-fit action releases the ligand and all atoms of protein residues within a 6 \AA{} pocket shell while outer atoms remain fixed.
\item \textbf{Agent-operated design.} Agents invoke the versioned Chemist Actions API rather than arbitrary page functions. The API exposes the same bounded state-changing verbs available to a chemist---including select, edit, search, apply, relax, and undo---plus read-only inspection. Requests execute serially and record sequence, timestamps, arguments, status, and duration.
\end{lstlisting}

\subsection{R099: Molarium architecture, numerical validation, and executable design history}

Baseline main.tex line 432; paragraph.

Subsection: Content-addressed molecular and decision provenance.

\paragraph{Version 1 --- earlier/alternative wording}

\begin{itemize}[leftmargin=*]
\item Source A, line 428: reviewed topic or figure-label match; not proof of edit lineage.
\item Source B, line 428: reviewed topic or figure-label match; not proof of edit lineage.
\end{itemize}

\begin{lstlisting}
The current live application retains at most 500 Chemist Actions audit records in memory; clearing the scene clears that audit unless the relevant actions have been exported or attached to a campaign commit. The Design History interface and public \texttt{campaign.*} actions can create a local campaign, atomically commit the visible molecule, create and restore branches, record decisions, merge branch histories, and verify the append-only record. Campaign JSON and the active branch are persisted in browser IndexedDB. Each commit contains a complete molecular snapshot; an attached action script records the completed public operations observed since the previous commit but is explicitly marked as incomplete coverage.
\end{lstlisting}

\paragraph{Version 2 --- CURRENT}

\begin{itemize}[leftmargin=*]
\item Source G1, line 432: exact text.
\item Source LIVE, line 432: exact text.
\end{itemize}

\begin{lstlisting}
The current live application retains at most 500 action records in memory; clearing the scene clears that audit. Durable provenance therefore requires explicit export of API moves or a replay log. The append-only design-campaign ledger, commit, branch, decision, and verification functions exist in the repository and support standalone builders and reviews, while one-click promotion of a live session into ledger commits is not currently exposed in the main interface.
\end{lstlisting}

\subsection{R101: Molarium architecture, numerical validation, and executable design history}

Baseline main.tex line 449; paragraph.

Subsection: Content-addressed molecular and decision provenance.

\paragraph{Version 1 --- earlier/alternative wording}

\begin{itemize}[leftmargin=*]
\item Source A, line 445: inferred correspondence, similarity 1.00; verify before merging.
\item Source B, line 445: inferred correspondence, similarity 1.00; verify before merging.
\end{itemize}

\begin{lstlisting}
Canonical JSON and SHA-256 fingerprints bind records to exact serialized representations. The \path{molarium.design-campaign/v1} schema stores snapshots, commits, branches, events, and decisions. Each commit identifies its selected snapshot, parent commit(s), optional action script, hypotheses, evidence, tags, and branch. Events are append-only and contain the previous event fingerprint, actor class (human, agent, system, or import), and source. Supported decision states are \texttt{progressed}, \texttt{not-progressed}, \texttt{deferred}, \texttt{failed}, \texttt{duplicate}, \texttt{superseded}, and \texttt{archived}. Finalization appends a completion event and hashes the campaign; subsequent record changes fail verification. Hashes provide alteration detection for the recorded representation, not authorship, chemical equivalence, or scientific validity.
\end{lstlisting}

\paragraph{Version 2 --- CURRENT}

\begin{itemize}[leftmargin=*]
\item Source G1, line 449: exact text.
\item Source LIVE, line 449: exact text.
\end{itemize}

\begin{lstlisting}
Canonical JSON and SHA-256 fingerprints bind records to exact serialized representations. The \texttt{molarium.design-campaign/v1} schema stores snapshots, commits, branches, events, and decisions. Each commit identifies its selected snapshot, parent commit(s), optional action script, hypotheses, evidence, tags, and branch. Events are append-only and contain the previous event fingerprint, actor class (human, agent, system, or import), and source. Supported decision states are \texttt{progressed}, \texttt{not-progressed}, \texttt{deferred}, \texttt{failed}, \texttt{duplicate}, \texttt{superseded}, and \texttt{archived}. Finalization appends a completion event and hashes the campaign; subsequent record changes fail verification. Hashes provide alteration detection for the recorded representation, not authorship, chemical equivalence, or scientific validity.
\end{lstlisting}

\subsection{R102: Molarium architecture, numerical validation, and executable design history}

Baseline main.tex line 451; paragraph.

Subsection: Content-addressed molecular and decision provenance.

\paragraph{Version 1 --- earlier/alternative wording}

\begin{itemize}[leftmargin=*]
\item Source A, line 447: inferred correspondence, similarity 0.99; verify before merging.
\item Source B, line 447: inferred correspondence, similarity 0.99; verify before merging.
\end{itemize}

\begin{lstlisting}
Registered design routes use a separate \path{molarium.registered-design-route/v1} schema that specifies the coordinate boundary, hash-pinned starting state, and ordered graph edits. The SOS1, BCL-xL, and CDK2 routes use this schema. Route actions are \texttt{designRoute.load}, \texttt{designRoute.applyStep}, and \texttt{designRoute.inspect}; campaign-ledger operations remain separate. Cross-schema validation rejects a route presented as a campaign ledger or a campaign ledger presented as a route. Frozen predictions made before the schema split retain their original input hash, with migration records storing both hashes and accepting only the schema-identifier substitution that reconstructs the frozen bytes.
\end{lstlisting}

\paragraph{Version 2 --- CURRENT}

\begin{itemize}[leftmargin=*]
\item Source G1, line 451: exact text.
\item Source LIVE, line 451: exact text.
\end{itemize}

\begin{lstlisting}
Registered design routes use a separate \texttt{molarium.registered-design-route/v1} schema that specifies the coordinate boundary, hash-pinned starting state, and ordered graph edits. The SOS1, BCL-xL, and CDK2 routes use this schema. Route actions are \texttt{designRoute.load}, \texttt{designRoute.applyStep}, and \texttt{designRoute.inspect}; campaign-ledger operations remain separate. Cross-schema validation rejects a route presented as a campaign ledger or a campaign ledger presented as a route. Frozen predictions made before the schema split retain their original input hash, with migration records storing both hashes and accepting only the schema-identifier substitution that reconstructs the frozen bytes.
\end{lstlisting}

\subsection{R105: Molarium architecture, numerical validation, and executable design history}

Baseline main.tex line 477; paragraph.

Subsection: Content-addressed molecular and decision provenance.

\paragraph{Version 1 --- earlier/alternative wording}

\begin{itemize}[leftmargin=*]
\item Source A, line 473: inferred correspondence, similarity 0.90; verify before merging.
\end{itemize}

\begin{lstlisting}
Executable replays use \texttt{molarium.chemist-action-script/v1}. Scripts contain ordered calls to the allow-listed Chemist Actions API rather than arbitrary code or direct coordinate replacement. Imports are capped at 2 MiB; one action's arguments are capped at 32 KiB, eight nested levels, and 2,048 values. Actions can name returned stable identifiers for reuse by later steps. The runner executes one action at a time through \texttt{api.execute(\{action,args\})}, derives a repeatable request ID from the script fingerprint, and stops on the first failure. Replay results use \texttt{molarium.chemist-action-replay/v1} and record the source fingerprint and each step outcome. Presentation actions can alter focus or styling but cannot change molecular graph or coordinates except through the same allow-listed state-changing actions.
\end{lstlisting}

\paragraph{Version 2 --- CURRENT}

\begin{itemize}[leftmargin=*]
\item Source B, line 473: exact text.
\item Source G1, line 477: exact text.
\item Source LIVE, line 477: exact text.
\end{itemize}

\begin{lstlisting}
Executable replays use \texttt{molarium.chemist-action-script/v1}. Scripts contain ordered calls to the allow-listed Chemist Actions API rather than arbitrary code or direct coordinate replacement. Designer Moves JSON files must be smaller than 2 MB; each serialized API request and its argument object are capped at 8 MiB for ASCII JSON, eight nested levels, and 2,048 JSON nodes. Actions can name returned stable identifiers for reuse by later steps. The runner executes one action at a time through \texttt{api.execute(\{action,args\})}, derives a repeatable request ID from the script fingerprint, and stops on the first failure. Replay results use \texttt{molarium.chemist-action-replay/v1} and record the source fingerprint and each step outcome. Presentation actions can alter focus or styling but cannot change molecular graph or coordinates except through the same allow-listed state-changing actions.
\end{lstlisting}

\subsection{R108: Molarium architecture, numerical validation, and executable design history}

Baseline main.tex line 485; paragraph.

Subsection: SOS1 hit-to-BAY-293 executable replay.

\paragraph{Version 1 --- earlier/alternative wording}

\begin{itemize}[leftmargin=*]
\item Source A, line 481: inferred correspondence, similarity 1.00; verify before merging.
\end{itemize}

\begin{lstlisting}
The selected route contains 33 recorded API actions: 29 molecular or protocol operations and four workspace switches. The full exploration audit contains 89 completed API calls, including three Phe890/pose branches, nine undo calls, and deterministic reselection. The rendered interface sequence adds 15 display-only actions for 48 replay steps. Later crystal structures were opened only after four prediction checkpoints and the exploration audit had been hashed. Display and focus actions are non-mutating.
\end{lstlisting}

\paragraph{Version 2 --- CURRENT}

\begin{itemize}[leftmargin=*]
\item Source B, line 481: exact text.
\item Source G1, line 485: exact text.
\item Source LIVE, line 485: exact text.
\end{itemize}

\begin{lstlisting}
The selected route contains 33 recorded API actions: 29 molecular or protocol operations and four workspace switches. The full exploration audit contains 89 completed API calls, including three Phe890/pose branches, nine undo calls, and deterministic reselection. The rendered interface sequence adds 16 display-only actions for 49 replay steps. Later crystal structures were opened only after four prediction checkpoints and the exploration audit had been hashed. Display and focus actions are non-mutating.
\end{lstlisting}

\subsection{R111: Molarium architecture, numerical validation, and executable design history}

Baseline main.tex line 505; paragraph.

Subsection: SOS1 hit-to-BAY-293 executable replay.

\paragraph{Version 1 --- earlier/alternative wording}

\begin{itemize}[leftmargin=*]
\item Source A, line 501: inferred correspondence, similarity 0.88; verify before merging.
\item Source B, line 501: inferred correspondence, similarity 0.88; verify before merging.
\end{itemize}

\begin{lstlisting}
Each ``64-chain'' search comprises 64 independently seeded local pose searches, distributed over eight browser workers in the recorded replay. Parameterization without motion assigns the numerical force-field representation while preserving coordinates; only the subsequent relaxation changes coordinates. The SOS1 route uses whole-system Sage 2.1.0 with deterministic Gasteiger charges. AWT, AWZ, AWW, and AXH contain 7,929, 7,933, 7,935, and 7,937 particles, respectively. This charge assignment is part of the recorded protocol and is distinct from standard Sage workflows using AM1-BCC or NAGL charges.
\end{lstlisting}

\paragraph{Version 2 --- CURRENT}

\begin{itemize}[leftmargin=*]
\item Source G1, line 505: exact text.
\item Source LIVE, line 505: exact text.
\end{itemize}

\begin{lstlisting}
Each ``64-chain'' search comprises 64 independently seeded local pose searches. Parameterization without motion assigns the numerical force-field representation while preserving coordinates; only the subsequent relaxation changes coordinates. The SOS1 route uses whole-system Sage 2.1.0/Gasteiger preparation rather than the prepared Rosemary validation fixtures. AWT, AWZ, AWW, and AXH contain 7,929, 7,933, 7,935, and 7,937 particles, respectively. This charge assignment is part of the recorded protocol and is distinct from standard Sage workflows using AM1-BCC or NAGL charges.
\end{lstlisting}

\subsection{R113: Molarium architecture, numerical validation, and executable design history}

Baseline main.tex line 509; structured block.

Subsection: SOS1 hit-to-BAY-293 executable replay.

\paragraph{Version 1 --- earlier/alternative wording}

\begin{itemize}[leftmargin=*]
\item Source A, line 505: inferred correspondence, similarity 0.99; verify before merging.
\item Source B, line 505: inferred correspondence, similarity 0.99; verify before merging.
\end{itemize}

\begin{lstlisting}
\begin{table}[htbp]
\centering\small
\caption{Sequence of molecular states in the SOS1 replay. Predictions, rather than future crystal structures, seed the next step.}
\label{tab:appA_sos1_beats}
\begin{tabularx}{\linewidth}{@{}p{0.07\linewidth}p{0.25\linewidth}X@{}}
\toprule
\textbf{Step} & \textbf{Operation} & \textbf{State/evidence}\\
\midrule
0 & Load and prepare & 5OVE/AXE only; capture the compound 1 reference and coordinate boundary\\
1 & Rewrite to compound 17 (AWT) & Search independent reference-guided poses; apply the selected pose; parameterize; relax\\
2 & Rewrite to compound 18 (AWZ) & Propagate from predicted AWT, not 5OVF; apply, parameterize, relax\\
3 & Grow compound 21 (AWW) & New ring enters the Phe890-in volume; fixed-core representative gains four clashes\\
4 & Resolve pocket state & Enumerate discrete Phe890 rotamers; select outward branch; retain the moved receptor as the next reference\\
5 & Predict BAY-293 (AXH) & Grow from frozen AWW; 64-chain pose search; apply; parameterize without motion; relax; freeze\\
6 & Evaluate holdouts & Compare AWT/AWZ/AWW/BAY-293 with 5OVF/5OVG/5OVH/5OVI\\
\bottomrule
\end{tabularx}
\end{table}
\end{lstlisting}

\paragraph{Version 2 --- CURRENT}

\begin{itemize}[leftmargin=*]
\item Source G1, line 509: exact text.
\item Source LIVE, line 509: exact text.
\end{itemize}

\begin{lstlisting}
\begin{table}[htbp]
\centering\small
\caption{Sequence of molecular states in the SOS1 replay. Predictions, rather than future crystal structures, seed the next step.}
\label{tab:appA_sos1_beats}
\begin{tabularx}{\linewidth}{@{}p{0.07\linewidth}p{0.25\linewidth}X@{}}
\toprule
\textbf{Step} & \textbf{Operation} & \textbf{State/evidence}\\
\midrule
0 & Load and prepare & 5OVE/AXE only; capture the compound 1 reference and coordinate boundary\\
1 & Rewrite to compound 17 (AWT) & Run 64 independent reference-guided pose searches; apply selected pose; parameterize; relax\\
2 & Rewrite to compound 18 (AWZ) & Propagate from predicted AWT, not 5OVF; apply, parameterize, relax\\
3 & Grow compound 21 (AWW) & New ring enters the Phe890-in volume; fixed-core representative gains four clashes\\
4 & Resolve pocket state & Enumerate discrete Phe890 rotamers; select outward branch; retain the moved receptor as the next reference\\
5 & Predict BAY-293 (AXH) & Grow from frozen AWW; 64-chain pose search; apply; parameterize without motion; relax; freeze\\
6 & Evaluate holdouts & Compare AWT/AWZ/AWW/BAY-293 with 5OVF/5OVG/5OVH/5OVI\\
\bottomrule
\end{tabularx}
\end{table}
\end{lstlisting}

\subsection{R118: Molarium architecture, numerical validation, and executable design history}

Baseline main.tex line 550; paragraph.

Subsection: SOS1 hit-to-BAY-293 executable replay.

\paragraph{Version 1 --- earlier/alternative wording}

\begin{itemize}[leftmargin=*]
\item Source A, line 546: inferred correspondence, similarity 0.96; verify before merging.
\item Source B, line 546: inferred correspondence, similarity 0.96; verify before merging.
\end{itemize}

\begin{lstlisting}
The interface replay executes from a blank canvas through the public Chemist Actions API. Pose and rotamer generation leave coordinates unchanged until an explicit Apply action. The replay therefore represents discrete state transitions rather than interpolated trajectories or molecular dynamics. Its 48-action rendered sequence lasts 69.25~s while the action audit records approximately 420.7~s of operation durations; timing remains attached to the audit and render manifest, and the rendered sequence is not used as a performance benchmark. Comparison of fragment-bound 6EPM with the 5OVE starting structure shows that the inhibitory Tyr884 arrangement is already present at the coordinate boundary; the hit-only route does not predict a Tyr884 transition.
\end{lstlisting}

\paragraph{Version 2 --- CURRENT}

\begin{itemize}[leftmargin=*]
\item Source G1, line 550: exact text.
\item Source LIVE, line 550: exact text.
\end{itemize}

\begin{lstlisting}
The interface replay executes from a blank canvas through the public Chemist Actions API. Candidate-generation operations leave coordinates unchanged until an explicit Apply action. The replay therefore represents discrete state transitions rather than interpolated trajectories or molecular dynamics. Its 48-action rendered sequence lasts 62.58 s while the underlying audit records approximately 908.5 s of operation durations; timing remains attached to the audit and render manifest, and the rendered sequence is not used as a performance benchmark. Comparison of fragment-bound 6EPM with the 5OVE starting structure shows that the inhibitory Tyr884 arrangement is already present at the coordinate boundary; the hit-only route does not predict a Tyr884 transition.
\end{lstlisting}

\subsection{Earlier material without a current-text match}

Absence of a match does not establish intentional rejection. The following include directive is preserved, but its target file is unavailable in the ZIP.

Source A, line 552; Working with Molarium: scientific workflows for humans and agents.

\begin{lstlisting}
\begin{small}
\setlength{\emergencystretch}{3em}

Molarium is designed so that a scientist using the graphical interface and an agent using the Chemist Actions API act on the same authoritative molecular state and encounter the same scientific boundaries. This appendix is organized by scientific task. Each workflow states the question, the human procedure, the corresponding agent skill contract, the evidence that must be retained, and the point at which the workflow must stop.

\textbf{Implementation basis.} The descriptions below refer to the inspected browser implementation on 2 September 2026, feature branch \texttt{feature/sos1-interface-story} through commit \texttt{f5a8900}. ``Current'' means reachable in the main browser application or installed public API. ``Experimental'' means executable but explicitly bounded by the implementation or its scientific validation. Repository-only and proposed facilities are identified rather than presented as operating instructions. User-facing workspace names are View, Design, and Simulate; serialized mode values remain \texttt{view}, \texttt{build}, and \texttt{run} so saved action scripts remain replayable.

\subsection{One molecular state, two interfaces}\label{appB:state}

\textbf{Status: current browser surface.}

\textbf{Question.} What changes the molecule, and what merely changes how it is displayed?

Molarium separates the authoritative molecular state from its presentation. Camera motion, representation, focus, labels, highlighting, and component visibility are view-only. Loading a structure, committing chemistry, applying a protonation state, applying a refined pose, or running a coordinate-producing Simulate job changes the state. Current Simulate calculations are state-mutating: non-energy jobs apply returned coordinates on completion, and selecting or playing a saved frame writes that frame's coordinates into the current molecule without a separate Apply control or undo snapshot. Refined-pose and side-chain candidates remain non-mutating until their explicit Apply actions.

For a reproducible workflow, begin with \texttt{?blank}. A scientist can then paste XYZ or PDB text, upload supported XYZ or PDB data, enter a SMILES string or PDB identifier in connected mode, or choose a prepared example. After loading, inspect atom count, formal charge, connected components, bonding, and coordinates. PDB input loads the first model and resolves alternate locations using the implemented occupancy policy; XYZ bonding is inferred from proximity and must therefore be checked.

The same operating rule applies to a human and an agent: \emph{inspect, act, verify}. A successful operation is not automatically a scientifically appropriate one.

\paragraph{Agent skill contract: \texttt{operate\_molecular\_state}.}

\begin{description}[leftmargin=2.2cm,style=nextline]
\item[Goal] Establish the intended starting object without confusing a display change, candidate, or imported approximation with accepted molecular state.
\item[Preconditions] A supported input or registered story/route is available; the network boundary is known.
\item[Procedure] Start blank when reproducibility requires it; load by a supported route; inspect graph, components, charge, and coordinates; report ambiguities before mutation.
\item[Evidence] Source identity or hash when available, atom/component summary, and warnings about inferred bonds, alternate locations, or retrieval.
\item[Stop] Do not reinterpret unsupported input or silently substitute another structure.
\end{description}

Generic file loading is presently a graphical-interface operation rather than a Chemist Actions route. Once a molecule is loaded, the installed API can inspect it directly:

\begin{lstlisting}
await api.execute({action:"session.inspect",
  args:{scope:"all", includeCoordinates:false, maximumAtoms:500}})
await api.execute({action:"view.setDisplay",
  args:{representation:"both", showHydrogens:false,
        showPocketAtoms:true, showInteractions:true}})
\end{lstlisting}

Registered prospective routes and public structure stories have guarded loaders (\texttt{designRoute.load} and \texttt{structureStory.load}); neither is a general-purpose arbitrary-file loader.

\subsection{Prepare a protein for local molecular design}\label{appB:prepare}

\textbf{Status: experimental current surface.}

\textbf{Question.} What must happen before an experimental protein structure can support a bounded local calculation?

A deposited PDB entry may contain missing atoms, alternate locations, ambiguous protonation, crystallographic waters, unusual residues, or an incomplete ligand. Molarium's preparation procedure is conservative: it previews a bounded set of supported repairs and stops when the structure exceeds them.

\begin{enumerate}[leftmargin=*]
\item Load and inspect the protein, ligand, waters, other components, and intended binding site.
\item In Protein Preparation choose pH, histidine policy, supported missing-heavy-atom repair, ligand treatment, water treatment, and gap policy, then run Prepare structure.
\item Molarium constructs an internal preview. Unresolved blockers stop and expose it for review; a blocker-free preview is immediately parameterized and applied.
\item Inspect the resulting graph, coordinates, warnings, and attached preparation audit. Export that audit when its decisions must persist beyond the live session.
\end{enumerate}

Inspect ligand identity and completeness, protonation around the active site, bridging waters, reconstructed geometry, gaps, clashes, and the chosen parameter model. A successful preview does not certify production simulation readiness; the implementation records \texttt{readyForProductionSimulation:false}. Missing-loop construction, membranes, general metal coordination, covalent ligands, and periodic explicit solvent remain outside this workflow.

\paragraph{Agent skill contract: \texttt{prepare\_protein}.}

\begin{description}[leftmargin=2.2cm,style=nextline]
\item[Goal] Create an explicitly reviewed molecular state suitable for a supported local Molarium calculation.
\item[Preconditions] A PDB-derived structure is loaded and the downstream use is known.
\item[Inputs] pH, histidine, repair, ligand, water, and gap policies.
\item[Evidence] Options, proposed repairs, warnings, blockers, report, and parameterization identity.
\item[Stop] Do not manufacture missing loops, parameters, coordination, or unsupported chemistry.
\end{description}

Representative public actions are:

\begin{lstlisting}
await api.execute({action:"protein.prepare", args:{
  pH:7.4, histidine:"auto", repairMissingHeavy:true,
  ligandPolicy:"ccd", waterPolicy:"crucial", gapPolicy:"block"}})
await api.execute({action:"session.inspect",
  args:{scope:"pocket", includeCoordinates:false, maximumAtoms:500}})
\end{lstlisting}

The preparation action already parameterizes and applies a blocker-free result. \texttt{protein.parameterize} is a separate alternative for an already prepared or subsequently edited complex. The exact accepted argument vocabulary is returned by \texttt{api.describe()} and is the runtime authority; the sequence above illustrates the contract rather than replacing that manifest.

\subsection{Choose the protonation state being designed}\label{appB:protonation}

\textbf{Status: current ligand workflow; protein choices are part of experimental preparation.}

\textbf{Question.} Which charged and protonated species does the graph represent?

Protonation changes chemical identity, formal charge, hydrogen count, hydrogen bonding, and parameterization. For a ligand, choose pH, enumerate RDKit candidate states, inspect charge and changed sites, and apply one state explicitly. The reported weights are independent Henderson--Hasselbalch estimates from the implemented site table; they are not microscopic p$K_a$ predictions or protein-conditioned populations. For a protein, choose pH and histidine policy during preparation and inspect residues near ligands, metals, catalytic centers, and salt bridges.

Applying a ligand state replaces the graph and coordinates and invalidates topology-dependent results. A failed sanitization or embedding must leave the previous accepted state intact.

\paragraph{Agent skill contract: \texttt{select\_protonation\_state}.}

\begin{description}[leftmargin=2.2cm,style=nextline]
\item[Goal] Make protonation and formal charge explicit before editing or calculation.
\item[Preconditions] A valid structure and relevant pH or stated assumption are available.
\item[Evidence] pH, candidates, chosen total charge, changed sites, and rationale where several states remain plausible.
\item[Stop] Do not choose a different state merely because it sanitizes or runs.
\end{description}

Ligand protonation enumeration is not yet a dedicated Chemist Actions route. An agent must therefore treat it as a human-reviewed interface step, while protein protonation is carried within \texttt{protein.prepare}. The absence of an API verb is an execution boundary, not permission to synthesize one.

\subsection{Edit atoms, bonds, charges, and local geometry}\label{appB:edit}

\textbf{Status: current browser surface.}

\textbf{Question.} How can chemistry be changed without corrupting the graph or losing undoability?

The synchronized 2D and 3D views select the same persistent atom identifiers. Chemistry-changing operations are grouped into a transaction. In Design, select the intended editable atom or bond, apply related element, bond-order, charge, hydrogen, deletion, or addition operations, then choose Finish chemistry once. RDKit sanitizes the graph, reconciles hydrogens, polishes eligible local coordinates, remaps reference contacts where possible, invalidates obsolete parameterization, and creates one undo snapshot. Discard restores the pre-transaction state.

Coordinate-only distance, angle, and dihedral controls operate on selected connected paths. The moving side is determined from graph connectivity; ambiguous ring cuts are blocked. After any committed edit, inspect valence, charge, aromaticity, hydrogen count, stereochemistry, local geometry, close contacts, and agreement between 2D and 3D representations. Distance, angle, dihedral, and general stereochemistry controls are currently graphical-interface operations; the public API exposes the listed atom/bond chemistry actions but no general stereochemistry-edit action.

\paragraph{Agent skill contract: \texttt{edit\_molecular\_graph}.}

\begin{description}[leftmargin=2.2cm,style=nextline]
\item[Goal] Commit an explicit graph or local-geometry change while preserving atomic lineage and undoability.
\item[Preconditions] The editable component and persistent atom identifiers are unambiguous; no unrelated transaction is pending.
\item[Evidence] Pre-edit identity, ordered actions, changed atoms/bonds, sanitization, cleanup, and final inspection.
\item[Stop] Never bypass sanitization or inject coordinates as a substitute for a graph edit.
\end{description}

\begin{lstlisting}
await api.execute({action:"view.setMode", args:{mode:"build"}})
await api.execute({action:"selection.replace", args:{atomIds:[anchorId]}})
await api.execute({action:"chemistry.addAtom",
  args:{attachedToAtomId:anchorId, element:"Cl"}})
await api.execute({action:"chemistry.finish", args:{}})
await api.execute({action:"session.inspect",
  args:{scope:"ligand", includeCoordinates:true, maximumAtoms:200}})
\end{lstlisting}

The serialized value \texttt{build} selects the user-facing Design workspace; its persistence is intentional action-script compatibility.

\subsection{Grow an analogue and preserve its structural context}\label{appB:grow}

\textbf{Status: current builder; pose-aware extension is experimental.}

\textbf{Question.} How can a substituent be added while retaining what was learned from the parent structure?

Fragment growth combines two separate tasks: constructing the correct graph and placing new atoms in a plausible initial geometry. If the parent is bound and its pose matters, first capture the reference pose and any explicit core or interaction hypotheses. Select the attachment atom, stage the fragment, and make associated bond, leaving-group, charge, hydrogen, and stereochemical changes in one transaction. Use geometry controls only to create a starting placement; they do not score the pocket. Finish chemistry, inspect graph and stereochemistry, then relax or search poses as a separate step. Generic fragment-template staging and the geometry controls are currently interface-only. The public API supports individual atom/bond edits and registered \texttt{designRoute.applyStep} operations.

A successful build establishes graph validity, not the preferred rotamer, ring conformation, pocket pose, synthetic feasibility, or affinity. General combinatorial enumeration, automated campaign selection, and automatic MCS transfer of an independently imported analogue are not current workbench workflows.

\paragraph{Agent skill contract: \texttt{grow\_fragment}.}

\begin{description}[leftmargin=2.2cm,style=nextline]
\item[Goal] Build one intended analogue with a reviewable starting geometry and preserved parent atom lineage.
\item[Preconditions] Ligand, attachment atom, fragment chemistry, and need for reference continuity are explicit.
\item[Evidence] Parent/product inspection, atom mapping, transaction actions, initial placement rationale, and captured contacts or reference.
\item[Stop] Treat relaxation and pose selection as new scientific decisions; do not infer them from successful graph construction.
\end{description}

For a registered, hash-pinned prospective route, an agent can stage a predefined product graph. \texttt{attachmentAtomId} is required only when that route explicitly declares designer-directed \texttt{spatialIntent}; the SOS1 route does not:

\begin{lstlisting}
await api.execute({action:"designRoute.load",
  args:{routeId:"sos1-hit-only"}})
await api.execute({action:"view.setMode", args:{mode:"build"}})
await api.execute({action:"protein.prepare", args:{
  pH:7.4, histidine:"auto", repairMissingHeavy:true,
  ligandPolicy:"ccd", waterPolicy:"retain", gapPolicy:"cap"}})
await api.execute({action:"pose.captureReference",
  args:{mode:"propagate"}})
await api.execute({action:"designRoute.applyStep",
  args:{stepId:"scaffold-rewrite"}})
await api.execute({action:"designRoute.inspect", args:{}})
\end{lstlisting}

Registered routes constrain prospective inputs and permitted edits; they are not retrospective campaign ledgers.

\subsection{Relax a structure, run bounded dynamics, and inspect motion}\label{appB:simulate}

\textbf{Status: current and experimental methods with method-specific limits.}

\textbf{Question.} Does a designed geometry survive a declared local model, and what moves during a bounded calculation?

Choose the method from the question rather than treating all engines as interchangeable:

\begin{center}
\begin{tabularx}{\linewidth}{@{}p{0.23\linewidth}p{0.19\linewidth}X@{}}
\toprule
\textbf{Question} & \textbf{Method} & \textbf{Current boundary}\\
\midrule
Small-molecule cleanup & RDKit force field & Workbench optimization and short motion; not a production trajectory.\\
Sage energy or local relaxation & Direct WebGPU & Experimental, 32-bit kernels, nonperiodic; no PME.\\
Many replicas of one topology & WebGPU ensemble & Homogeneous replicas; limits on atoms, replicas, steps, and pair work.\\
Eligible neutral-molecule electronic energy & ANI-2x & H/C/N/O/S/F/Cl, explicit H, neutral singlet, at most 96 atoms.\\
Independent internal reference & OpenMM 8.2 Reference & Internal parameterization, checking, and bounded relaxation; not a general selectable method.\\
\bottomrule
\end{tabularx}
\end{center}

Finish pending chemistry and establish complete parameterization. In Simulate, record job type, model, charge model, solvent, constraints, movable atoms, step count, saved frames, and ensemble settings. Inspect engine identity, finite energies, convergence or diagnostics, elapsed time, warnings, and motion in the region relevant to the question. Because coordinate-producing jobs and subsequent frame selection immediately update the current molecule, export the starting coordinates first when reversibility matters. Display alignment does not overwrite the raw coordinates retained in the calculation record.

Automatic Sage parameterization uses deterministic Gasteiger charges. Unsupported chemistry or missing terms fail explicitly. The current local workflows do not provide periodic solvent, PME, a barostat, production protein dynamics, or binding free energies.

\paragraph{Agent skill contract: \texttt{relax\_or\_simulate}.}

\begin{description}[leftmargin=2.2cm,style=nextline]
\item[Goal] Use a declared model to test local geometry or bounded motion while retaining interpretable metadata.
\item[Preconditions] Chemistry is finished, method eligibility is established, and an acceptance observable is defined.
\item[Inputs] Method, steps, solvent, constraints, fixed/movable atoms, saved frames, and ensemble settings.
\item[Evidence] Model/version, force field, charge model, solvent, constraints, timestep, energies, diagnostics, timing, frames, and warnings.
\item[Stop] A minimized pose or stable short trace is not an affinity estimate, equilibrium ensemble, free energy, or production trajectory.
\end{description}

\begin{lstlisting}
await api.execute({action:"view.setMode", args:{mode:"build"}})
await api.execute({action:"optimization.run",
  args:{method:"induced-fit-webgpu"}})
await api.execute({action:"session.inspect",
  args:{scope:"pocket", includeCoordinates:true, maximumAtoms:500}})
\end{lstlisting}

This public action is a bounded Design optimizer, not the agent equivalent of a general Simulate calculation. The current API has no action for general Simulate jobs, dynamics, or trajectories; those remain human-operated until a dedicated action is added.

\subsection{Explore conformations with Conformer Arena}\label{appB:conformers}

\textbf{Status: experimental current surface.}

\textbf{Question.} Which conformational basins can a finite search find, and how sensitive are they to the generation and refinement method?

Begin with finished chemistry, explicit hydrogens, and a chosen protonation state. Generate deterministic RDKit seeds, refine them with the declared finite protocol, cluster by RMSD, and compare energies, distinct clusters, recurring torsions, ring states, diagnostics, and failures. The STORMM-inspired path uses homogeneous replicas through staged relaxation, heating, cooling, and final relaxation. Conformer Arena compares supported pathways on the same molecular graph and settings.

The generated ensemble and raw energy records remain calculation evidence, but the current Conformer Arena workflow changes displayed/current coordinates as frames or conformers are selected. Cluster populations from this anneal/minimize search are not thermodynamic populations; failure to find a basin does not establish inaccessibility; energies from different force fields are not interchangeable.

\paragraph{Agent skill contract: \texttt{explore\_conformers}.}

\begin{description}[leftmargin=2.2cm,style=nextline]
\item[Goal] Generate, refine, cluster, and compare a finite set of conformational alternatives for one topology.
\item[Preconditions] Chemistry and protonation are final; method eligibility and intended diversity metric are stated.
\item[Evidence] Seed, methods, requested/returned counts, energies, clusters, constraint diagnostics, timings, and exported ensemble where applicable.
\item[Stop] Do not describe a finite heuristic search as exhaustive or Boltzmann weighted.
\end{description}

Conformer Arena does not yet have a dedicated Chemist Actions route. An agent may operate it only through a separately reviewed interface automation or future public action; it must not imitate the result through direct state injection.

\subsection{Predict a short single-chain protein structure}\label{appB:fold}

\textbf{Status: experimental current surface; connected mode only.}

\textbf{Question.} Can Molarium predict one protein chain from sequence, and is this co-folding?

The implemented path predicts one protein entity from one sequence; it is not ligand co-folding or multimer prediction. A configured external service produces the MSA, after which OpenFold ONNX inference runs locally in the browser, preferring WebGPU and falling back to WebAssembly. Local Lab disables the external MSA route.

Enter one sequence of at most 128 residues, disclose and accept transmission to the selected MSA endpoint, choose recycle settings, run inference, and inspect residue mapping, chain completeness, geometry, and downstream suitability. First use can download approximately 540~MB of model assets. A successful prediction replaces the current molecule; export PDB before loading other work. The path does not use templates or perform Amber relaxation. It creates no server-side job ledger.

\paragraph{Agent skill contract: \texttt{predict\_protein\_structure}.}

\begin{description}[leftmargin=2.2cm,style=nextline]
\item[Goal] Predict coordinates for one short protein sequence using the declared MSA endpoint and local model.
\item[Preconditions] Connected mode, accepted endpoint, sequence length at most 128, and accepted model-asset download.
\item[Evidence] Input sequence, endpoint, bucket/model identity, recycle settings, progress or failure record, and exported PDB when persistence matters.
\item[Stop] Do not claim ligand co-folding, multimer prediction, experimental validation, or simulation readiness.
\end{description}

This workflow has no dedicated Chemist Actions verb in the current manifest. The contract therefore documents how an agent must reason about a human-visible workflow; it does not imply an installed autonomous route.

\subsection{Maintain an experimental pose while changing ligand chemistry}\label{appB:pose}

\textbf{Status: experimental current surface.}

\textbf{Question.} How can an edited analogue inherit trusted information from an experimental ligand pose without pretending that the answer is known?

Reference-guided pose maintenance is narrower than global docking. Capture the prepared reference complex before editing. The record includes ligand lineage, receptor atoms within 8~\AA{}, reference contacts, selected interaction hypotheses, settings, and provenance. Edit and finish the analogue, inspect remapped or unavailable contacts, then search only the degrees of freedom released by the edit. Molarium generates deterministic seeds, parameterizes the edited ligand, rejects infeasible candidates, relaxes feasible candidates in a rigid local receptor site, and ranks distinct results.

Inspect more than rank 1: mapped-core displacement, preserved and replaced contacts, ligand strain, clashes, ring/side-chain choices, and rejection reasons. Only \texttt{pose.apply} changes authoritative coordinates. The score contains local Lennard--Jones and Coulomb cross terms and relative vacuum ligand strain; it omits general receptor relaxation, water displacement, entropy, binding free energy, and broad concerted macrocycle/fused-ring search.

\paragraph{Agent skill contract: \texttt{maintain\_reference\_pose}.}

\begin{description}[leftmargin=2.2cm,style=nextline]
\item[Goal] Generate reviewable local poses for an edited analogue while retaining explicit relationships to a trusted reference.
\item[Preconditions] Prepared complex, captured reference, finished chemistry, unchanged receptor, and available required contacts.
\item[Evidence] Settings, hashes, atom mapping, contacts, constraints, rejected candidates, scores, selection, and chained labbook.
\item[Stop] If no feasible candidate survives, apply nothing and retain the negative record.
\end{description}

\begin{lstlisting}
await api.execute({action:"pose.captureReference",
  args:{mode:"propagate"}})
// perform and finish the intended molecular-graph edit
await api.execute({action:"pose.refine",
  args:{searchChains:16}})
await api.execute({action:"pose.apply", args:{index:0}})
await api.execute({action:"session.inspect",
  args:{scope:"pocket", includeCoordinates:true, maximumAtoms:500}})
\end{lstlisting}

For an explicit receptor-side-chain alternative, call \path{pose.enumerateSidechainRotamers}, apply one with \path{pose.applySidechainRotamer}, and then call \path{pose.updateReceptorReference} before another \texttt{pose.refine}. Candidate enumeration itself remains non-mutating.

\subsection{Record, replay, and hand off a design process}\label{appB:history}

\textbf{Status: current action/replay and Design History surfaces.}

\textbf{Question.} Which records allow another scientist or agent to reproduce the procedure, inspect the result, and understand the decision?

Molarium retains distinct evidence objects because they answer different questions:

\begin{center}
\begin{tabularx}{\linewidth}{@{}p{0.23\linewidth}X@{}}
\toprule
\textbf{Record} & \textbf{Meaning}\\
\midrule
Chemist Actions audit & Serialized requests, outcomes, failures, and timings for recognized public operations.\\
Calculation/preparation record & Model, protocol, result, warnings, and method-specific audit.\\
Docking labbook & Mapping, constraints, feasibility, ranking, rejection, selection, and chained integrity records.\\
Designer action script & JSON sequence of completed public actions replayed through the same guarded API. It is not necessarily a complete or standalone reconstruction.\\
Registered design route & Prospective, hash-pinned coordinate boundary and permitted ordered graph edits. It is not retrospective history.\\
Design History campaign & Content-addressed molecular snapshots, commits, branches, merge parents, decisions, events, and verification state.\\
\bottomrule
\end{tabularx}
\end{center}

Designer Moves import begins from a blank canvas and performs script-level validation of schema, recognized names, binding order, and forbidden shortcut keys. Individual action arguments and molecular preconditions are checked when each step executes. Audit conversion includes only completed public actions and excludes campaign administration, but a resulting script does not contain an arbitrary loaded starting molecule. Standalone replay therefore requires a reproducible loader or a separately supplied starting state. Failed audit entries remain evidence but are not converted silently into successful replay steps. Replay has play/pause, previous/next-step controls, stable central captions, and a separate replay audit.

Design History is a live browser feature rather than repository-only scaffolding. \texttt{campaign.create} creates a locally persisted campaign and atomically commits the visible molecule. \texttt{campaign.commitCurrent} stores a complete graph-and-coordinate snapshot together with completed public actions since the prior commit. Branch switching restores the exact head molecule and refuses to proceed when the visible molecular state is dirty. Merge records the currently visible target-branch molecule as an explicit result with both parents; it is not an automatic structural merge. Decisions attach a disposition, rationale, reason codes, and optional evidence to a commit. Verification recomputes content hashes, the append-only event chain, commit relationships, and derived branch heads.

Campaign JSON and the active branch are stored in browser IndexedDB. Reload restores the committed graph and coordinates at the active branch head; it does not restore parameter tables, docking candidates, calculation frames, view state, or the complete workbench session. Import verifies the campaign and requires a restorable graph/coordinate snapshot before mutating the live state. Closing a campaign clears the active pointer without deleting its commits. A live commit may attach an action script marked \texttt{coverage.kind:"public-actions-only"} and \texttt{coverage.complete:false}; the complete molecular snapshot is authoritative. Campaign-management actions are intentionally excluded from molecule replay scripts to avoid recursively replaying history administration.

\paragraph{Agent skill contract: \texttt{export\_and\_handoff}.}

\begin{description}[leftmargin=2.2cm,style=nextline]
\item[Goal] Preserve the smallest complete molecular, procedural, numerical, and decision record needed by the recipient.
\item[Preconditions] The authoritative result and downstream purpose are explicit; unsaved work has not been cleared.
\item[Procedure] Inventory records; export reusable molecular state and method evidence; create an action script when replay is required; commit decision history when project continuity is required; verify hashes and replay.
\item[Evidence] Files, schemas, versions, hashes, replay audit, decision rationale, and unresolved failures.
\item[Stop] Do not describe an image, share URL, partial action script, or failed replay as a complete scientific handoff.
\end{description}

\begin{lstlisting}
await api.execute({action:"campaign.create", args:{
  campaignId:"sos1-local-design", title:"SOS1 local design",
  initialCommitMessage:"Register 5OVE/AXE starting state"}})
// chemist or agent performs guarded molecular actions
await api.execute({action:"campaign.commitCurrent", args:{
  message:"Accept Phe890-out design state", tags:["prospective"]}})
await api.execute({action:"campaign.recordDecision", args:{
  disposition:"progressed", rationale:"No new clashes after pocket-state search"}})
await api.execute({action:"campaign.verify", args:{}})
\end{lstlisting}

\subsection{How skill contracts, API actions, and Design History fit together}\label{appB:contractapi}

The three layers serve different purposes and should not be collapsed:

\begin{enumerate}[leftmargin=*]
\item \textbf{The skill contract defines scientific intent.} It states preconditions, permissible operations, required evidence, success, failure, and escalation. In this appendix these names denote proposed narrative operating contracts stored with the manuscript. They are not installed runtime skills, Chemist Actions, schemas, or executable enforcement.
\item \textbf{The Chemist Actions API defines executable authority.} \texttt{api.describe()} lists the installed action names and bounded arguments. Calls are validated, serialized through one queue, and audited with sequence, status, result or error, and duration. An agent cannot gain extra molecular authority by phrasing a request differently.
\item \textbf{Design History defines durable meaning over time.} It commits complete molecular states and may link them to an explicitly partial public-action script, branches, decisions, evidence, actors, and append-only integrity records. The snapshot answers ``what state was accepted?''; the action script answers ``which guarded operations were captured and can be replayed?''; the decision answers ``why was it accepted?''
\end{enumerate}

The practical agent pattern is therefore:

\begin{lstlisting}
const manifest = api.describe();             // discover installed authority
await api.execute({action:"session.inspect", args:{scope:"all"}});
// choose only actions permitted by the relevant scientific contract
await api.execute({action:"...", args:{...}});
await api.execute({action:"session.inspect", args:{scope:"all"}});
await api.execute({action:"campaign.commitCurrent",
  args:{message:"Accepted state after declared check"}});
await api.execute({action:"campaign.recordDecision",
  args:{disposition:"progressed", rationale:"..."}});
await api.execute({action:"campaign.verify", args:{}});
\end{lstlisting}

The ellipsis is not a wildcard action: it represents one explicit name returned by the manifest. Workflows without a dedicated public action remain human-reviewed or unavailable to autonomous execution. This distinction prevents an agent skill from overstating what the installed software can actually do.

\textbf{Handoff rule.} Export the structure when the recipient needs coordinates; export the calculation record when the recipient must judge the method; export the Designer Moves script together with its required starting state when the recipient must replay operations; and commit to Design History when the recipient must recover branches, decisions, and accepted molecular states. Hashes detect alteration of the serialized record. They do not establish authorship, chemical equivalence, or scientific validity.

\textbf{Deliberate boundaries.} The current workbench does not provide missing-loop construction, membrane or general-metal preparation, covalent-ligand parameterization, periodic explicit solvent, PME, production-scale trajectories, binding free energy, global docking, general remote CUDA execution, GFN/xTB, general combinatorial enumeration, or automatic campaign selection. Their absence is part of the operating contract rather than an invitation to approximate them silently.

\end{small}

\end{lstlisting}

\subsection{Index of records with one recovered version}

These passages are already present in the manuscript at the baseline line indicated. Only one distinct wording was recovered from the sources listed; this does not prove a passage was never edited. The index includes paragraphs, list items, tables, figures, examples, and formatting definitions, as classified by source-block extraction.

\begin{longtable}{@{}p{0.10\linewidth}p{0.10\linewidth}p{0.18\linewidth}p{0.54\linewidth}@{}}
\toprule Record & Line & Sources & Location\\\midrule\endhead
R004 & 104 & A, B, G1, LIVE & Abstract \\
R016 & 144 & A, B, G1, LIVE & Molarium and the value of a persistent molecular state \\
R019 & 152 & A, B, G1, LIVE & Molarium and the value of a persistent molecular state \\
R025 & 171 & A, B, G1, LIVE & Scientific contracts: provenance, validation, and falsifiability \\
R026 & 173 & A, B, G1, LIVE & Scientific contracts: provenance, validation, and falsifiability \\
R029 & 184 & A, B, G1, LIVE & Scientific contracts: provenance, validation, and falsifiability \\
R030 & 191 & A, B, G1, LIVE & Scientific contracts: provenance, validation, and falsifiability \\
R033 & 199 & A, B, G1, LIVE & Local computation changes the interaction model \\
R036 & 207 & A, B, G1, LIVE & What becomes scarce when the tool layer becomes cheap \\
R037 & 209 & A, B, G1, LIVE & What becomes scarce when the tool layer becomes cheap \\
R038 & 211 & A, B, G1, LIVE & What becomes scarce when the tool layer becomes cheap \\
R039 & 213 & A, B, G1, LIVE & What becomes scarce when the tool layer becomes cheap \\
R040 & 215 & A, B, G1, LIVE & What becomes scarce when the tool layer becomes cheap \\
R041 & 217 & A, B, G1, LIVE & What becomes scarce when the tool layer becomes cheap \\
R042 & 224 & A, B, G1, LIVE & What becomes scarce when the tool layer becomes cheap \\
R043 & 226 & A, B, G1, LIVE & What becomes scarce when the tool layer becomes cheap \\
R044 & 228 & A, B, G1, LIVE & What becomes scarce when the tool layer becomes cheap \\
R045 & 232 & A, B, G1, LIVE & The changing role of scientific expertise \\
R046 & 234 & A, B, G1, LIVE & The changing role of scientific expertise \\
R047 & 236 & A, B, G1, LIVE & The changing role of scientific expertise \\
R048 & 238 & A, B, G1, LIVE & The changing role of scientific expertise \\
R049 & 240 & A, B, G1, LIVE & The changing role of scientific expertise \\
R050 & 244 & A, B, G1, LIVE & Limitations and outlook \\
R051 & 246 & A, B, G1, LIVE & Limitations and outlook \\
R052 & 248 & A, B, G1, LIVE & Limitations and outlook \\
R053 & 250 & A, B, G1, LIVE & Limitations and outlook \\
R054 & 252 & A, B, G1, LIVE & Limitations and outlook \\
R055 & 254 & A, B, G1, LIVE & Limitations and outlook \\
R056 & 258 & A, B, G1, LIVE & Conclusion \\
R057 & 260 & A, B, G1, LIVE & Conclusion \\
R058 & 262 & A, B, G1, LIVE & Conclusion \\
R059 & 264 & A, B, G1, LIVE & Conclusion \\
R061 & 270 & A, B, G1, LIVE & AI assistance and availability \\
R063 & 278 & A, B, G1, LIVE & Molarium architecture, numerical validation, and executable design history \\
R064 & 280 & A, B, G1, LIVE & Molarium architecture, numerical validation, and executable design history \\
R065 & 284 & A, B, G1, LIVE & Execution architecture \\
R066 & 286 & A, B, G1, LIVE & Execution architecture \\
R067 & 288 & A, B, G1, LIVE & Execution architecture \\
R068 & 300 & A, B, G1, LIVE & Execution architecture \\
R069 & 302 & A, B, G1, LIVE & Execution architecture \\
R070 & 320 & A, B, G1, LIVE & Execution architecture \\
R071 & 322 & A, B, G1, LIVE & Execution architecture \\
R072 & 326 & A, B, G1, LIVE & OpenMM WebAssembly reference \\
R073 & 328 & A, B, G1, LIVE & OpenMM WebAssembly reference \\
R074 & 330 & A, B, G1, LIVE & OpenMM WebAssembly reference \\
R075 & 332 & A, B, G1, LIVE & OpenMM WebAssembly reference \\
R076 & 336 & A, B, G1, LIVE & Direct WebGPU mechanics \\
R077 & 338 & A, B, G1, LIVE & Direct WebGPU mechanics \\
R078 & 340 & A, B, G1, LIVE & Direct WebGPU mechanics \\
R079 & 342 & A, B, G1, LIVE & Direct WebGPU mechanics \\
R080 & 360 & A, B, G1, LIVE & Direct WebGPU mechanics \\
R081 & 362 & A, B, G1, LIVE & Direct WebGPU mechanics \\
R082 & 366 & A, B, G1, LIVE & Replica-parallel WebGPU \\
R083 & 368 & A, B, G1, LIVE & Replica-parallel WebGPU \\
R084 & 370 & A, B, G1, LIVE & Replica-parallel WebGPU \\
R085 & 372 & A, B, G1, LIVE & Replica-parallel WebGPU \\
R086 & 374 & A, B, G1, LIVE & Replica-parallel WebGPU \\
R087 & 378 & A, B, G1, LIVE & Reference-driven implementation and validation \\
R088 & 381 & A, B, G1, LIVE & Reference-driven implementation and validation \\
R089 & 389 & A, B, G1, LIVE & Reference-driven implementation and validation \\
R090 & 391 & A, B, G1, LIVE & Reference-driven implementation and validation \\
R091 & 393 & A, B, G1, LIVE & Reference-driven implementation and validation \\
R092 & 410 & A, B, G1, LIVE & Reference-driven implementation and validation \\
R093 & 412 & A, B, G1, LIVE & Reference-driven implementation and validation \\
R094 & 416 & A, B, G1, LIVE & Workflow orchestration and local execution boundary \\
R095 & 418 & A, B, G1, LIVE & Workflow orchestration and local execution boundary \\
R097 & 426 & A, B, G1, LIVE & Workflow orchestration and local execution boundary \\
R098 & 430 & A, B, G1, LIVE & Content-addressed molecular and decision provenance \\
R100 & 434 & A, B, G1, LIVE & Content-addressed molecular and decision provenance \\
R103 & 453 & A, B, G1, LIVE & Content-addressed molecular and decision provenance \\
R104 & 455 & A, B, G1, LIVE & Content-addressed molecular and decision provenance \\
R106 & 479 & A, B, G1, LIVE & Content-addressed molecular and decision provenance \\
R107 & 483 & A, B, G1, LIVE & SOS1 hit-to-BAY-293 executable replay \\
R109 & 487 & A, B, G1, LIVE & SOS1 hit-to-BAY-293 executable replay \\
R110 & 489 & A, B, G1, LIVE & SOS1 hit-to-BAY-293 executable replay \\
R112 & 507 & A, B, G1, LIVE & SOS1 hit-to-BAY-293 executable replay \\
R114 & 528 & A, B, G1, LIVE & SOS1 hit-to-BAY-293 executable replay \\
R115 & 530 & A, B, G1, LIVE & SOS1 hit-to-BAY-293 executable replay \\
R116 & 532 & A, B, G1, LIVE & SOS1 hit-to-BAY-293 executable replay \\
R117 & 548 & A, B, G1, LIVE & SOS1 hit-to-BAY-293 executable replay \\
R119 & 558 & B, G1, LIVE & Working with Molarium: scientific workflows for humans and agents \\
R120 & 563 & B, G1, LIVE & Working with Molarium: scientific workflows for humans and agents \\
R121 & 566 & B, G1, LIVE & Working with Molarium: scientific workflows for humans and agents \\
R122 & 572 & B, G1, LIVE & Working with Molarium: scientific workflows for humans and agents \\
R123 & 578 & B, G1, LIVE & Working with Molarium: scientific workflows for humans and agents \\
R124 & 584 & B, G1, LIVE & Working with Molarium: scientific workflows for humans and agents \\
R125 & 600 & B, G1, LIVE & Working with Molarium: scientific workflows for humans and agents \\
R126 & 603 & B, G1, LIVE & Working with Molarium: scientific workflows for humans and agents \\
R127 & 605 & B, G1, LIVE & Working with Molarium: scientific workflows for humans and agents \\
R128 & 608 & B, G1, LIVE & Working with Molarium: scientific workflows for humans and agents \\
R129 & 610 & B, G1, LIVE & Working with Molarium: scientific workflows for humans and agents \\
R130 & 612 & B, G1, LIVE & Working with Molarium: scientific workflows for humans and agents \\
R131 & 614 & B, G1, LIVE & Working with Molarium: scientific workflows for humans and agents \\
R132 & 618 & B, G1, LIVE & B.1: One molecular state, two interfaces \\
R133 & 620 & B, G1, LIVE & B.1: One molecular state, two interfaces \\
R134 & 622 & B, G1, LIVE & B.1: One molecular state, two interfaces \\
R135 & 626 & B, G1, LIVE & Starting from a structure \\
R136 & 628 & B, G1, LIVE & Starting from a structure \\
R137 & 632 & B, G1, LIVE & State and persistence \\
R138 & 634 & B, G1, LIVE & State and persistence \\
R139 & 646 & B, G1, LIVE & State and persistence \\
R140 & 648 & B, G1, LIVE & State and persistence \\
R141 & 663 & B, G1, LIVE & State and persistence \\
R142 & 667 & B, G1, LIVE & B.2: Prepare a protein for local molecular design \\
R143 & 672 & B, G1, LIVE & Human workflow \\
R144 & 682 & B, G1, LIVE & What should I inspect? \\
R145 & 689 & B, G1, LIVE & What should I inspect? \\
R146 & 691 & B, G1, LIVE & What should I inspect? \\
R147 & 693 & B, G1, LIVE & What should I inspect? \\
R148 & 706 & B, G1, LIVE & What should I inspect? \\
R149 & 708 & B, G1, LIVE & What should I inspect? \\
R150 & 716 & B, G1, LIVE & What should I inspect? \\
R151 & 720 & B, G1, LIVE & B.3: Choose the protonation state you are designing \\
R152 & 725 & B, G1, LIVE & Ligand workflow \\
R153 & 733 & B, G1, LIVE & Protein workflow \\
R154 & 735 & B, G1, LIVE & Protein workflow \\
R155 & 737 & B, G1, LIVE & Protein workflow \\
R156 & 749 & B, G1, LIVE & Protein workflow \\
R157 & 751 & B, G1, LIVE & Protein workflow \\
R158 & 759 & B, G1, LIVE & Protein workflow \\
R159 & 763 & B, G1, LIVE & B.4: Edit atoms, bonds, charges, and stereochemistry in 3D \\
R160 & 768 & B, G1, LIVE & Human workflow \\
R161 & 778 & B, G1, LIVE & What should I inspect? \\
R162 & 784 & B, G1, LIVE & What should I inspect? \\
R163 & 786 & B, G1, LIVE & What should I inspect? \\
R164 & 788 & B, G1, LIVE & What should I inspect? \\
R165 & 800 & B, G1, LIVE & What should I inspect? \\
R166 & 813 & B, G1, LIVE & What should I inspect? \\
R167 & 817 & B, G1, LIVE & B.5: Grow an analogue by adding and positioning fragments \\
R168 & 822 & B, G1, LIVE & Human workflow \\
R169 & 831 & B, G1, LIVE & How should I judge the result? \\
R170 & 833 & B, G1, LIVE & How should I judge the result? \\
R171 & 835 & B, G1, LIVE & How should I judge the result? \\
R172 & 847 & B, G1, LIVE & How should I judge the result? \\
R173 & 849 & B, G1, LIVE & How should I judge the result? \\
R174 & 856 & B, G1, LIVE & How should I judge the result? \\
R175 & 870 & B, G1, LIVE & How should I judge the result? \\
R176 & 874 & B, G1, LIVE & B.6: Relax a structure, run bounded dynamics, and inspect motion \\
R177 & 877 & B, G1, LIVE & B.6: Relax a structure, run bounded dynamics, and inspect motion \\
R178 & 878 & B, G1, LIVE & B.6: Relax a structure, run bounded dynamics, and inspect motion \\
R179 & 899 & B, G1, LIVE & Human workflow \\
R180 & 908 & B, G1, LIVE & Interpreting the model \\
R181 & 910 & B, G1, LIVE & Interpreting the model \\
R182 & 922 & B, G1, LIVE & Interpreting the model \\
R183 & 924 & B, G1, LIVE & Interpreting the model \\
R184 & 932 & B, G1, LIVE & Interpreting the model \\
R185 & 934 & B, G1, LIVE & Interpreting the model \\
R186 & 944 & B, G1, LIVE & B.7: Explore conformations with Conformer Arena \\
R187 & 949 & B, G1, LIVE & Human workflow \\
R188 & 959 & B, G1, LIVE & What does the result establish? \\
R189 & 961 & B, G1, LIVE & What does the result establish? \\
R190 & 973 & B, G1, LIVE & What does the result establish? \\
R191 & 975 & B, G1, LIVE & What does the result establish? \\
R192 & 985 & B, G1, LIVE & B.8: Predict a short single-chain protein structure \\
R193 & 987 & B, G1, LIVE & B.8: Predict a short single-chain protein structure \\
R194 & 992 & B, G1, LIVE & Human workflow \\
R195 & 999 & B, G1, LIVE & Human workflow \\
R196 & 1001 & B, G1, LIVE & Human workflow \\
R197 & 1013 & B, G1, LIVE & Human workflow \\
R198 & 1015 & B, G1, LIVE & Human workflow \\
R199 & 1024 & B, G1, LIVE & B.9: Maintain an experimental pose while changing ligand chemistry \\
R200 & 1029 & B, G1, LIVE & Human workflow \\
R201 & 1039 & B, G1, LIVE & What is being scored? \\
R202 & 1041 & B, G1, LIVE & What is being scored? \\
R203 & 1043 & B, G1, LIVE & What is being scored? \\
R204 & 1055 & B, G1, LIVE & What is being scored? \\
R205 & 1066 & B, G1, LIVE & What is being scored? \\
R206 & 1068 & B, G1, LIVE & What is being scored? \\
R207 & 1076 & B, G1, LIVE & What is being scored? \\
R208 & 1080 & B, G1, LIVE & B.10: Record, replay, and hand off a design process \\
R209 & 1083 & B, G1, LIVE & B.10: Record, replay, and hand off a design process \\
R210 & 1084 & B, G1, LIVE & B.10: Record, replay, and hand off a design process \\
R211 & 1105 & B, G1, LIVE & Replay workflow \\
R212 & 1108 & B, G1, LIVE & Replay workflow \\
R213 & 1115 & B, G1, LIVE & Replay workflow \\
R214 & 1119 & B, G1, LIVE & Live Design History \\
R215 & 1121 & B, G1, LIVE & Live Design History \\
R216 & 1126 & B, G1, LIVE & Choose the export that matches the handoff \\
R217 & 1127 & B, G1, LIVE & Choose the export that matches the handoff \\
R218 & 1147 & B, G1, LIVE & Choose the export that matches the handoff \\
R219 & 1159 & B, G1, LIVE & Choose the export that matches the handoff \\
R220 & 1174 & B, G1, LIVE & B.11: How skill contracts, API actions, and Design History fit together \\
R221 & 1177 & B, G1, LIVE & B.11: How skill contracts, API actions, and Design History fit together \\
R222 & 1182 & B, G1, LIVE & B.11: How skill contracts, API actions, and Design History fit together \\
R223 & 1184 & B, G1, LIVE & B.11: How skill contracts, API actions, and Design History fit together \\
R224 & 1197 & B, G1, LIVE & B.11: How skill contracts, API actions, and Design History fit together \\
R225 & 1199 & B, G1, LIVE & B.11: How skill contracts, API actions, and Design History fit together \\
R226 & 1201 & B, G1, LIVE & B.11: How skill contracts, API actions, and Design History fit together \\
R227 & 1205 & B, G1, LIVE & Capabilities deliberately outside these workflows \\
R228 & 1207 & B, G1, LIVE & Capabilities deliberately outside these workflows \\
\bottomrule\end{longtable}

\subsection{Coverage and limitations}

The reconstruction examines 7 named sources and 228 current paragraph or structured records. It preserves 38 groups with multiple versions and indexes 190 groups with one version. Every extracted source-record occurrence is assigned either to a baseline group or to the unmatched inventory. Each baseline group contains exactly one LIVE record. The preamble, author metadata, section-heading history, bibliography entries, and image pixels are outside this paragraph inventory.

The first two appendices and the main text remain the scientific manuscript. This third appendix is an editorial aid containing superseded and potentially contradictory claims; it should not be treated as an extension of the software's operating contract.

\endgroup
\clearpage

\begin{thebibliography}{99}

\bibitem{bran2024chemcrow}
Bran, A. M.; Cox, S.; Schilter, O.; Baldassari, C.; White, A. D.; Schwaller, P.
Augmenting Large Language Models with Chemistry Tools.
\emph{Nature Machine Intelligence} \textbf{2024}, \emph{6}, 525--535.

\bibitem{boiko2023coscientist}
Boiko, D. A.; MacKnight, R.; Kline, B.; Gomes, G.
Autonomous Chemical Research with Large Language Models.
\emph{Nature} \textbf{2023}, \emph{624}, 570--578.

\bibitem{pham2026chemgraph}
Pham, T. D.; Tanikanti, A.; Ke{\c{c}}eli, M.
ChemGraph as an Agentic Framework for Computational Chemistry Workflows.
\emph{Communications Chemistry} \textbf{2026}, \emph{9}, 33.

\bibitem{macknight2025rethinking}
MacKnight, R.; Boiko, D. A.; Regio, J. E.; Gallegos, L. C.; Neukomm, T. A.; et al.
Rethinking Chemical Research in the Age of Large Language Models.
\emph{Nature Computational Science} \textbf{2025}, \emph{5}, 715--726.

\bibitem{haas2017wasm}
Haas, A.; Rossberg, A.; Schuff, D. L.; et al.
Bringing the Web up to Speed with WebAssembly.
\emph{Proc. PLDI} \textbf{2017}, 185--200.

\bibitem{webgpu2025}
W3C GPU for the Web Working Group.
\emph{WebGPU Specification}, 2025.

\bibitem{riniker2015etkdg}
Riniker, S.; Landrum, G. A.
Better Informed Distance Geometry: Using What We Know to Improve Conformation Generation.
\emph{J. Chem. Inf. Model.} \textbf{2015}, \emph{55}, 2562--2574.

\bibitem{boothroyd2023sage}
Boothroyd, S.; Madin, O. C.; Mobley, D. L.; et al.
Development and Benchmarking of Open Force Field 2.0.0: The Sage Small-Molecule Force Field.
\emph{J. Chem. Theory Comput.} \textbf{2023}, \emph{19}, 3251--3275.

\bibitem{eastman2024openmm8}
Eastman, P.; Galvelis, R.; Pel{\'a}ez, R. P.; et al.
OpenMM 8: Molecular Dynamics Simulation with Machine Learning Potentials.
\emph{J. Phys. Chem. B} \textbf{2024}, \emph{128}, 109--116.

\bibitem{cerutti2024stormm}
Cerutti, D. S.; Wiewiora, R. P.; Boothroyd, S.; Sherman, W.
STORMM: Structure and Topology Replica Molecular Mechanics for Chemical Simulations.
\emph{J. Chem. Phys.} \textbf{2024}, \emph{161}, 032501.

\bibitem{devereux2020ani2x}
Devereux, C.; Smith, J. S.; Huddleston, K. K.; et al.
Extending the Applicability of the ANI Deep Learning Molecular Potential to Sulfur and Halogens.
\emph{J. Chem. Theory Comput.} \textbf{2020}, \emph{16}, 4192--4202.

\bibitem{ahdritz2024openfold}
Ahdritz, G.; Bouatta, N.; Floristean, C.; et al.
OpenFold: Retraining AlphaFold2 Yields New Insights into Its Learning Mechanisms and Capacity for Generalization.
\emph{Nature Methods} \textbf{2024}, \emph{21}, 1514--1524.

\bibitem{ryckaert1977shake}
Ryckaert, J.-P.; Ciccotti, G.; Berendsen, H. J. C.
Numerical Integration of the Cartesian Equations of Motion of a System with Constraints: Molecular Dynamics of n-Alkanes.
\emph{J. Comput. Phys.} \textbf{1977}, \emph{23}, 327--341.

\bibitem{andersen1983rattle}
Andersen, H. C.
RATTLE: A ``Velocity'' Version of the SHAKE Algorithm for Molecular Dynamics Calculations.
\emph{J. Comput. Phys.} \textbf{1983}, \emph{52}, 24--34.

\bibitem{cohenboulakia2017workflows}
Cohen-Boulakia, S.; Belhajjame, K.; Collin, O.; et al.
Scientific Workflows for Computational Reproducibility in the Life Sciences: Status, Challenges and Opportunities.
\emph{Future Gener. Comput. Syst.} \textbf{2017}, \emph{75}, 284--298.

\bibitem{pritchard2025reproducibility}
Pritchard, N. J.; Wicenec, A.
Formal Definition and Implementation of Reproducibility Tenets for Computational Workflows.
\emph{Future Gener. Comput. Syst.} \textbf{2025}, \emph{166}, 107684.

\bibitem{wilkinson2016fair}
Wilkinson, M. D.; Dumontier, M.; Aalbersberg, I. J.; et al.
The FAIR Guiding Principles for Scientific Data Management and Stewardship.
\emph{Sci. Data} \textbf{2016}, \emph{3}, 160018.

\bibitem{bender2021realistic1}
Bender, A.; Cort{\'e}s-Ciriano, I.
Artificial Intelligence in Drug Discovery: What Is Realistic, What Are Illusions? Part 1.
\emph{Drug Discov. Today} \textbf{2021}, \emph{26}, 511--524.

\bibitem{bender2021realistic2}
Bender, A.; Cort{\'e}s-Ciriano, I.
Artificial Intelligence in Drug Discovery: What Is Realistic, What Are Illusions? Part 2.
\emph{Drug Discov. Today} \textbf{2021}, \emph{26}, 1040--1052.

\bibitem{wigh2022representations}
Wigh, D. S.; Goodman, J. M.; Lapkin, A. A.
A Review of Molecular Representation in the Age of Machine Learning.
\emph{WIREs Comput. Mol. Sci.} \textbf{2022}, \emph{12}, e1603.

\bibitem{bainbridge1983ironies}
Bainbridge, L.
Ironies of Automation.
\emph{Automatica} \textbf{1983}, \emph{19}, 775--779.


\bibitem{webgpuWGSL2026}
W3C GPU for the Web Working Group.
\emph{WebGPU Shading Language}, 2026. \url{https://www.w3.org/TR/WGSL/}.

\bibitem{wang2020macrocycle}
Wang, S.; Witek, J.; Landrum, G. A.; Riniker, S.
Improving Conformer Generation for Small Rings and Macrocycles Based on Distance Geometry and Experimental Torsional-Angle Preferences.
\emph{J. Chem. Inf. Model.} \textbf{2020}, \emph{60}, 2044--2058.

\bibitem{halgren1996mmff94}
Halgren, T. A.
Merck Molecular Force Field. I. Basis, Form, Scope, Parameterization, and Performance of MMFF94.
\emph{J. Comput. Chem.} \textbf{1996}, \emph{17}, 490--519.

\bibitem{rappe1992uff}
Rapp\'{e}, A. K.; Casewit, C. J.; Colwell, K. S.; Goddard, W. A.; Skiff, W. M.
UFF, a Full Periodic Table Force Field for Molecular Mechanics and Molecular Dynamics Simulations.
\emph{J. Am. Chem. Soc.} \textbf{1992}, \emph{114}, 10024--10035.

\bibitem{gasteiger1980charges}
Gasteiger, J.; Marsili, M.
Iterative Partial Equalization of Orbital Electronegativity---A Rapid Access to Atomic Charges.
\emph{Tetrahedron} \textbf{1980}, \emph{36}, 3219--3228.

\bibitem{eastman2017openmm7}
Eastman, P.; Swails, J.; Chodera, J. D.; et al.
OpenMM 7: Rapid Development of High Performance Algorithms for Molecular Dynamics.
\emph{PLoS Comput. Biol.} \textbf{2017}, \emph{13}, e1005659.

\bibitem{onufriev2004gb}
Onufriev, A.; Bashford, D.; Case, D. A.
Exploring Protein Native States and Large-Scale Conformational Changes with a Modified Generalized Born Model.
\emph{Proteins} \textbf{2004}, \emph{55}, 383--394.

\bibitem{glaser2015gpu}
Glaser, J.; Nguyen, T. D.; Anderson, J. A.; et al.
Strong Scaling of General-Purpose Molecular Dynamics Simulations on GPUs.
\emph{Comput. Phys. Commun.} \textbf{2015}, \emph{192}, 97--107.

\bibitem{hillig2019sos1}
Hillig, R. C.; Sautier, B.; Schroeder, J.; et al.
Discovery of Potent SOS1 Inhibitors That Block RAS Activation via Disruption of the RAS--SOS1 Interaction.
\emph{Proc. Natl. Acad. Sci. U.S.A.} \textbf{2019}, \emph{116}, 2551--2560.

\end{thebibliography}

\end{document}
